\documentclass[11pt]{article} \usepackage[utf8]{inputenc}
\usepackage[margin=1in]{geometry} \usepackage{amsmath} \usepackage{amssymb}
\usepackage{amsthm} \usepackage{mathtools} \usepackage{microtype}
\usepackage{graphicx} \usepackage{tikz} \usepackage{booktabs} \usepackage{array}
\usepackage{float} \usepackage{caption} \usepackage{subcaption}
\usepackage{xcolor} \usepackage{algorithm} \usepackage{algpseudocode}
\usepackage{aliascnt} \usepackage[numbers,sort&compress]{natbib}
\usepackage{hyperref} \usepackage[nameinlink,capitalise,noabbrev]{cleveref}
\makeatletter
\def\theHALG@line{\thealgorithm.\arabic{ALG@line}}
\makeatother
\hypersetup{hidelinks,pdftitle={Interlacing and Additive Inapproximability of Deterministic Communication Complexity}}
\usetikzlibrary{arrows.meta,positioning,calc,matrix,fit,backgrounds}
\theoremstyle{plain} \newtheorem{theorem}{Theorem}[section]
\newaliascnt{lemma}{theorem} \newtheorem{lemma}[lemma]{Lemma}
\aliascntresetthe{lemma} \newaliascnt{proposition}{theorem}
\newtheorem{proposition}[proposition]{Proposition}
\aliascntresetthe{proposition} \newaliascnt{corollary}{theorem}
\newtheorem{corollary}[corollary]{Corollary} \aliascntresetthe{corollary}
\newaliascnt{claim}{theorem} \newtheorem{claim}[claim]{Claim}
\aliascntresetthe{claim} \newaliascnt{observation}{theorem}
\newtheorem{observation}[observation]{Observation}
\aliascntresetthe{observation} \newaliascnt{fact}{theorem}
\newtheorem{fact}[fact]{Fact} \aliascntresetthe{fact}
\newaliascnt{conjecture}{theorem}
\newtheorem{conjecture}[conjecture]{Conjecture} \aliascntresetthe{conjecture}
\theoremstyle{definition} \newaliascnt{definition}{theorem}
\newtheorem{definition}[definition]{Definition} \aliascntresetthe{definition}
\newaliascnt{assumption}{theorem}
\newtheorem{assumption}[assumption]{Assumption} \aliascntresetthe{assumption}
\newaliascnt{example}{theorem} \newtheorem{example}[example]{Example}
\aliascntresetthe{example} \theoremstyle{remark} \newaliascnt{remark}{theorem}
\newtheorem{remark}[remark]{Remark} \aliascntresetthe{remark}

%% Setup carried over from notes.tex
\usetikzlibrary{arrows.meta,positioning,calc,matrix,fit,decorations.pathreplacing}
\definecolor{aliceblue}{RGB}{32,96,151} \definecolor{boborange}{RGB}{183,96,27}
\definecolor{answergreen}{RGB}{34,111,81} \definecolor{linkblue}{RGB}{32,71,116}
\crefname{lemma}{Lemma}{Lemmas}
\crefname{proposition}{Proposition}{Propositions}
\crefname{corollary}{Corollary}{Corollaries}
\crefname{definition}{Definition}{Definitions}
\crefname{example}{Example}{Examples} \crefname{remark}{Remark}{Remarks}

%% Macros carried over from notes.tex
\newcommand{\D}{\operatorname{D}}
\newcommand{\Int}[2]{\mathsf{I}_{#1}(#2)}
\newcommand{\RP}[3]{\mathcal{P}(#1;#2,#3)}
\newcommand{\card}[1]{\lvert#1\rvert}
\newcommand{\ceil}[1]{\left\lceil#1\right\rceil}
\newcommand{\floor}[1]{\left\lfloor#1\right\rfloor}
\newcommand{\ind}{\mathbf{1}}
\newcommand{\NP}{\mathsf{NP}}
\newcommand{\BPP}{\mathsf{BPP}}
\newcommand{\SAT}{\mathsf{SAT}}
\DeclareMathOperator{\rank}{rank} \crefname{claim}{Claim}{Claims}
\crefname{observation}{Observation}{Observations} \crefname{fact}{Fact}{Facts}
\crefname{conjecture}{Conjecture}{Conjectures}
\crefname{assumption}{Assumption}{Assumptions}

\begin{document}
\title{Interlacing and Additive Inapproximability of Deterministic Communication Complexity}
\author{}
\date{}
\maketitle

\begin{abstract}
Given the full truth table of a Boolean function, how accurately can one
estimate its deterministic communication complexity? Earlier versions of
this paper established exact NP-hardness and ETH-based additive
inapproximability. We prove that approximation within every fixed additive
error is NP-hard under deterministic polynomial-time reductions, both for
unrestricted protocols and for Yao's strictly alternating model, including
receiver output. For unrestricted protocols, randomised sparse ETH rules
out bounded-error randomised polynomial-time estimation within
$c\log T/\log\log T$ for every fixed $0\le c<1/8$, where $T$ is the number
of table entries. Deterministic ETH rules out deterministic polynomial-time
estimation within $c_*\log\log T$ for an effective absolute constant $c_*>0$.

These results bring together the interlacing program, from its origins in
a direct-sum counterexample through relaxed interlacing and exact hardness
to repeatable gap amplification and sparsification. The central obstacle is
to enlarge a communication gap while keeping the table small enough to
print. We develop an interlacing step whose extra bit survives input
deletions and whose lower-bound guarantees are restored for the next step.
Short explicit input lists make every fixed gap polynomial in size;
keeping the full products implicit and sampling once gives the stronger
randomised bound. The same size-control principle extends the program to
strict alternation: an explicit array list preserves a different lookup
lower bound that exploits the fixed speaking schedule.
\end{abstract}

\section{Introduction}
\label{sec:intro}

Alice knows a row of a Boolean matrix and Bob knows a column. They must
determine the entry where the two meet, communicating as few bits as possible.
Their local computation is free, and their protocol must answer every input
correctly. The deterministic communication complexity $\D(M)$ is the smallest
worst-case cost of such a protocol. We study the algorithmic problem in which
the \emph{whole table} is given: how accurately can this optimum be estimated?
The answers are already in the input; the difficulty is finding how little
communication suffices to produce them.

Yao asked this question in 1979, when he introduced communication
complexity~\cite{yao1979distributed}. Exact NP-hardness alone says little
about approximation. Costs $B$ and $B+1$ can receive
the same estimate if an error of one bit is allowed. Moreover, Yao required
the players to alternate after every bit, whereas unrestricted protocols
choose the speaker at each node. These are different optimisation problems.
We prove that approximation within every fixed additive error is NP-hard
in both models, and obtain growing-error hardness in the unrestricted model.

\paragraph{The interlacing program.}
This paper is the culmination of a line of work that connects the structure
of optimal protocols to the difficulty of finding their cost. Its organising
operation is \emph{interlacing}: Alice receives a base row and a selector,
Bob receives a tuple of base columns, and they must compute the entry at
the selected column. Mackenzie and Saffidine introduced this framework in
their direct-sum counterexample, first circulated in
2024~\cite{mackenzie2025refuting}. Interlacing made it possible to prove
that solving several independent instances together can cost less than
solving them separately.

The next obstacle was the size of the construction. Interlacing with all
column tuples creates an exponential input family. \emph{Relaxed
interlacing} replaces that family by a small pseudorandom one while
preserving the lower-bound argument. Earlier versions of this paper,
first circulated in 2025~\cite{gaspers-he-mackenzie-npcc}, used this method
to establish exact NP-hardness, concurrently and independently with
Hirahara, Ilango and Loff~\cite{hirahara2025communication}. The reduction
turns a violated combinatorial constraint into one unavoidable extra bit of
communication.

Approximation required that extra bit to accumulate. This led to
\emph{repeatable gap amplification}: make the lower bound survive the
input restrictions created by a protocol, and restore the guarantees
needed for another round. The earlier ETH stage of this
paper~\cite{gaspers2026interlacingeth}, prepared as version 5 of
arXiv:2508.05597, rules out every fixed additive error under ETH and gives
a computable unbounded error allowance. Its fixed-gap constructions take
subexponential time in the source length; they do not give the
polynomial-time reductions needed for fixed-error NP-hardness.

\paragraph{What this paper adds.}
The present work makes the amplified gap survive in a small enough table
for every fixed gap to be produced in deterministic polynomial time.
Thus fixed-error hardness now follows from NP-hardness itself. We also
obtain explicit growing-error rates under deterministic and randomised
ETH, replacing the earlier unspecified unbounded allowance by quantitative
bounds. The key advance is preserving the properties needed by the
lower-bound proof while reducing its input families. One construction
keeps short explicit lists at every round; the other keeps full products
implicit and samples only after amplification. Interlacing supplies the
unrestricted lower bounds in both constructions, and restoration makes
the gap repeatable. Sparsification controls the size of the table that
must ultimately be printed.

This manuscript continues the same paper under arXiv:2508.05597. Our
version-specific references identify earlier stages of that work; the
approximation and strict-alternation results are later contributions.
The direct-sum paper is a separate precursor to this continuing program.

\paragraph{From sparsification to strict alternation.}
Separating the lower-bound structure from the size of its input family
also opens a route to Yao's model. The short protocols in the unrestricted
construction use blocks of consecutive bits from one player; inserting
the missing turns can consume their additive advantage. For strict
alternation, we build a lookup table whose lower bound uses the fixed
speaking schedule. Allowing Bob every Boolean array initially gives an
exponential family of columns. The proof only needs a bound on how often
an array matches a prescribed pattern of bits. A short explicit array
list preserves that bound and makes the table polynomial in size for
every fixed gap. The shared principle is sparsification; the alternating
proof counts response strategies and does not invoke the unrestricted
amplifier. This construction gives fixed-error hardness, with no
corresponding growing-error bound established here.

Write $T$ for the number of table entries; polynomial time means polynomial
in $T$. Additive error measures accuracy in bits, the natural scale here:
on an $n\times n$ table, Alice can send her row index and Bob the answer, so
$\D(M)\le\ceil{\log n}+1$. An estimator returns a number, not necessarily a
protocol. A bounded-error randomised estimator must meet its error guarantee
with probability at least $2/3$ on every input. The protocols whose costs
we estimate are always deterministic and exact.

We write $D_A^{\mathrm{alt}}(M)$ for the strictly alternating public-output
cost with Alice starting. Yao's receiver-output cost $Y_{\mathrm{free}}(M)$
allows an input-independent choice of starter and lets the receiver of the
last bit determine the answer. Stopping remains public.
\Cref{sec:models} fixes these conventions, including depth zero.
The growing-error results use sparse ETH: $3$-SAT on $N$ variables and a
fixed linear number of clauses has no $2^{o(N)}$-time algorithm. Its
deterministic and bounded-error randomised forms are distinguished in
\cref{def:reth}.

\subsection{Results}
\label{sec:results}

The central conclusion is NP-hardness for every fixed additive error in
both speaking models. The two quantitative growing-error bounds concern
unrestricted protocols and use the stated forms of ETH. The table separates
these conclusions and their assumptions.

\begin{center}
\small
\begin{tabular}{@{}p{.18\linewidth}p{.27\linewidth}p{.25\linewidth}p{.19\linewidth}@{}}
\toprule
Protocol model & Additive error & Basis of hardness & Estimator excluded\\
\midrule
Unrestricted & $c\log T/\log\log T$, $c<1/8$ & Randomised sparse ETH & Bounded-error randomised\\
\addlinespace
Unrestricted & $c_*\log\log T$ & Deterministic ETH & Deterministic\\
\addlinespace
Unrestricted & Every fixed error & NP-hard by deterministic reduction & See the consequences below\\
\addlinespace
Strict alternation & Every fixed error & NP-hard by deterministic reduction & See the consequences below\\
\bottomrule
\end{tabular}
\end{center}

\begin{theorem}[Randomised growing-error hardness]
\label{thm:quantitative}
Assume randomised sparse ETH and fix $0\le c<1/8$. No bounded-error randomised
estimator running in time polynomial in $T$ estimates $\D(M)$ within
additive error $c\log T/\log\log T$ on every sufficiently large explicit
total Boolean table. This remains true for square tables whose side length
is a power of two.
\end{theorem}

\begin{theorem}[Deterministic growing-error hardness]
\label{det:eth}
There is an effective absolute constant $c_*>0$ such that deterministic ETH
rules out every deterministic polynomial-time estimator that estimates
$\D(M)$ within additive error $c_*\log\log T$ on every sufficiently large
explicit total Boolean table.
\end{theorem}

For fixed gaps, no ETH assumption is needed. In each reduction below the
threshold $B$ is computed from the formula and requested gap, without
deciding satisfiability.

\begin{theorem}[Fixed gaps for unrestricted protocols]
\label{intro:deterministic}
For every fixed integer $k\ge1$, a deterministic polynomial-time algorithm
maps a $3$-CNF formula $\Phi$ to an explicit total Boolean table $M$ and an
integer $B$ such that
\[
\Phi\in\SAT\ \Longrightarrow\ \D(M)\le B,\qquad
\Phi\notin\SAT\ \Longrightarrow\ \D(M)\ge B+k.
\]
The polynomial may depend on $k$. The output may be square with
power-of-two side length.
\end{theorem}

\begin{theorem}[Fixed gaps under strict alternation]
\label{yao:main}
For every fixed integer $k\ge1$, a deterministic polynomial-time algorithm
maps a $3$-CNF formula $\Phi$ to an explicit total Boolean table $M$ and an
integer $B$ such that
\[
\Phi\in\SAT\ \Longrightarrow\ D_A^{\mathrm{alt}}(M)\le B,\qquad
\Phi\notin\SAT\ \Longrightarrow\ D_A^{\mathrm{alt}}(M)\ge B+k.
\]
The polynomial bounds on output size and running time may depend on $k$.
\end{theorem}

\begin{observation}\label{obs:gap-estimation}
A reduction with YES cost at most $B$ and NO cost at least $B+k$ defeats
additive error $a$ whenever $k>2a$. Thus, for every fixed error, estimating
either cost in the preceding two theorems in deterministic polynomial time
implies $\mathsf P=\NP$; bounded-error randomised estimation implies
$\NP\subseteq\BPP$.
\end{observation}
\begin{proof}
An accurate estimate is at most $B+a$ in the YES case and at least $B+k-a$
in the NO case. Since $B+a<B+k/2<B+k-a$, comparison with the public
midpoint decides satisfiability whenever the estimate is accurate. For
fixed $a$, choose a fixed integer $k>2a$ and compose this comparison with
the corresponding deterministic reduction.
\end{proof}

For example, a gap of three bits separates estimates with error one.
Changing the starter or announcing a receiver's answer costs at most one
bit each (\cref{yao:comparison}), so increasing the fixed gap also gives
Yao's output convention.

\begin{corollary}\label{yao:bpp}
For every fixed nonnegative additive error, a deterministic polynomial-time
estimator of $Y_{\mathrm{free}}(M)$ on all explicit total Boolean tables
implies $\mathsf P=\NP$. A bounded-error randomised estimator implies
$\NP\subseteq\BPP$. The same conclusions hold for public-output strict
alternation with Alice first, Bob first, or an input-independent choice
of starter.
\end{corollary}

The strict-alternation construction has polynomial size for each fixed
gap, with a degree that may depend on that gap. It gives no growing-error
bound. The stronger unrestricted rate uses both a randomised ETH assumption
and randomised estimators; under deterministic ETH our deterministic
obstruction is $c_*\log\log T$.

\subsection{Proof overview}
\label{sec:overview}

The two unrestricted reductions combine repeatable gap amplification with
sparsification. Recall the interlacing operation: Alice receives a base row
and a selector $i$, Bob receives a tuple of base columns, and they must
answer the $i$-th question.
Repetition enlarges the gap but also creates large products of input sets.
Deleting inputs shrinks the table, yet can erase the selector's effect:
if Bob's questions are all equal, the selector no longer matters.
\Cref{fig:lookup} gives the smallest example. A protocol itself creates
deletions as messages remove incompatible rows or columns.

\begin{figure}[H]
\centering
\begin{tikzpicture}[
  >=Stealth,
  every node/.style={font=\footnotesize,align=center},
  box/.style={draw,rounded corners=2pt,inner sep=5pt,minimum height=15mm},
  source/.style={box,text width=25mm},
  route/.style={box,text width=43mm},
  result/.style={box,text width=49mm}]
\node[font=\footnotesize\bfseries] at (0,2.25) {Starting problem};
\node[font=\footnotesize\bfseries] at (4.5,2.25) {Build the gap};
\node[font=\footnotesize\bfseries] at (10,2.25) {Obtain an explicit table};
\node[source] (graph) at (0,0) {Graph colouring};
\node[route] (full) at (4.5,1) {Repeat full products\\keep the table implicit};
\node[route,fill=black!4] (lists) at (4.5,-1)
  {Repeat with short lists\\control size at every round};
\node[result,fill=black!4] (random) at (10,1)
  {Sample once, after all rounds\\randomised unrestricted reduction};
\node[result] (det) at (10,-1)
  {Print the final table\\deterministic\\unrestricted reduction};
\node[source] (hyper) at (0,-3.35) {Hypergraph\\colouring};
\node[route] (lookup) at (4.5,-3.35) {Lookup on a short\\explicit array list};
\node[result] (alt) at (10,-3.35) {Print the full table\\strict-alternation reduction};
\coordinate (branch) at ($(graph.east)!0.5!(graph.east -| full.west)$);
\draw (graph.east) -- (branch);
\draw[->] (branch) |- (full.west);
\draw[->] (branch) |- (lists.west);
\draw[->] (full.east) -- (random.west);
\draw[->] (lists.east) -- (det.west);
\draw[->] (hyper.east) -- (lookup.west);
\draw[->] (lookup.east) -- (alt.west);
\draw[black!35,dashed] (-1.4,-2.2) -- (12.6,-2.2);
\end{tikzpicture}
\caption{Where size is controlled (shaded boxes). The two unrestricted routes
share the internal-interlace argument: one samples only after growing the gap,
while the other keeps the input lists short throughout. Strict alternation
uses a separate lookup construction. All sources come from $3$-SAT.}
\label{fig:reduction-chain}
\end{figure}

\paragraph{A repeatable extra bit.}
The graph supplies a simple obstruction to a short protocol. At each
edge coordinate, its two endpoints and a few auxiliary items are active.
A failed colouring forces two active items into a bin that also contains
an inactive item. These items request three distinct lookup positions.
First forcing the protocol to fix an assignment to bins makes the same
obstruction available after inputs are deleted.

To repeat the construction, we track two deletion thresholds. At a fixed
minimum row mass, retaining $3/5$ of the columns preserves the current
lower bound; retaining $3/10$ may lose one bit.
The lower bound then follows by tracking where the protocol sends its
inputs. Most copies retain a large unsplit share at a distinct
continuation; we call that continuation the copy's \emph{owner}.
Special diagonal inputs expose the collision at an owner, forcing some
of these protected inputs to go elsewhere.
Neither other owners nor the few exceptional continuations can receive
all this displaced mass. This contradiction handles arbitrary speaking
order. One final interlace restores the bounds after deletion that the
next round needs.

Both unrestricted constructions use the same argument for this internal
interlace. It bounds how much restricting the tuple inputs can change the
law of a few coordinates. The inner source estimates and final restoration
remain separate: their smaller retained events require stronger estimates
for full products.

\paragraph{One final sample.}
For the randomised reduction, we keep every product implicit through all
$k$ rounds. Rank bounds the number of column sets on which a short protocol
can succeed, uniformly over retained row sets. This lets a union bound
protect a column sample, followed by an independent row sample.
The samples have the same power-of-two size, so the output is already
square. Recursive entry evaluation prints it without expanding the products.
With our choice of $k$, the size is $\log T=(4+o(1))k\log N$ for a
sparse formula on $N$ variables,
which yields the coefficient $1/8$ after allowing error on both sides of
the gap. The sample need not support any further amplification.

\paragraph{Explicit lists at every round.}
For the deterministic reduction, short weighted lists replace the products.
They agree with the product law on every small set of coordinates.
We build a list one coordinate at a time and remove redundant support
points while preserving those marginals.
Conditioning on a transcript can destroy this agreement: Bob can announce
whether two independent questions agree, making them identical on that
branch. A moment bound shows that only a bounded set of coordinates has
badly changed small marginals. This supplies the common internal argument
and, at larger widths, the list-specific source and restoration bounds.
The list sizes compound over the rounds, giving at most $N^{2^{O(k)}}$
entries and the $\log\log T$ error scale.

\paragraph{Strict alternation.}
Here we sparsify a lookup table that encodes a hypergraph colouring
instance. If Bob could hold every Boolean array, the
column family would be exponential, but a uniform array would match a
prescribed pattern on $r$ positions with probability exactly $2^{-r}$.
A short explicit list retains this bound up to a small additive error.
This is the property needed by the following lower-bound argument.

A hypergraph-colouring argument forces many distinct lookup requests to
share Alice's early response strategy: she responds identically for
\emph{every} possible early history of Bob's bits. The remaining turns
allow too few output patterns to cover the surviving arrays, even with
the small error allowed by the list. The fixed speaking schedule makes
this counting argument possible. We print the whole list, obtaining a
deterministic reduction for each fixed gap.

\subsection{Related work}

Yao introduced the strictly alternating model and asked how hard its optimum
is to compute from a truth table~\cite{yao1979distributed}.
Kushilevitz and Weinreb proved approximation hardness under cryptographic
assumptions~\cite{kushilevitz2009complexity}. Ilango, Loff and Oliveira
proved NP-hardness and strong inapproximability for partial
matrices~\cite{ilango2020multioutput}. For total matrices, Hirahara,
Ilango and Loff proved additive bit-cost hardness for three-message protocols
and leaf-count hardness for every fixed number of messages at least
three~\cite{hirahara2021constant}. Messages may contain many bits, so
these results cover neither bit-by-bit alternation nor the unrestricted
optimum. Earlier versions of this paper~\cite{gaspers-he-mackenzie-npcc}
established exact unrestricted NP-hardness, concurrently and independently
with Hirahara, Ilango and Loff~\cite{hirahara2025communication}. The present
fixed-gap results address the additive-one and alternating questions left
open by those exact-hardness results.

The amplifier uses the interlacing operation of Mackenzie and
Saffidine~\cite{mackenzie2025refuting}. Its graph source combines Dinur's
gap amplification with the bounded-degree colouring reduction of
Bogdanov, Obata and Trevisan~\cite{dinur2007pcp,bogdanov2002lower}.
The alternating source is the balanced rainbow-colouring hardness of
Guruswami and Lee~\cite{guruswami2018balanced}, combined with an explicit
finite-field array construction based on the small-bias construction of
Alon et al.~\cite{alon1992smallbias}. The deterministic lists use
sequential extension and support compression in the tradition of Koller
and Megiddo~\cite{koller1994sample}. Rank-based counts have antecedents
in rankwidth and clique-width algorithms~\cite{bui2010rankwidth,oum2014faster}
and hard-submatrix arguments~\cite{hrubes2025hard}. In particular, the
row-sum method of Oum, S{\ae}ther and Vatshelle~\cite{oum2014faster}
motivates our short leaf-condition count. Mackenzie's lossless hardness
condensation~\cite{mackenzie2026condensation} also counts inputs served
by short protocols; it preserves hardness relative to the new input
length, whereas our final sample preserves an absolute threshold.
We prove the sampling theorem needed here directly.

\subsection{Formal verification}
\label{sec:formal-verification}

Formal verification accompanies this program. The independent Beyond the
Library team~\cite{moakhar2026beyond} has formalised the corrected
Mackenzie--Saffidine direct-sum result~\cite{mackenzie2025refuting} and its
supporting arguments in Lean. Their
\href{https://beyondthelibrary.github.io/formal_archive/p/Refuting/}{public
proof inspector} gives the statements and formal proofs.

A separate companion formalisation of the present paper is being developed
for its sparsified constructions and the strictly alternating model.
The final artifact will specify the verified theorem coverage, source
version, assumptions and correspondence with the statements here.
% FORMALISATION_ARTIFACT_LINK: before submission replace the visible placeholder
% with the permanent artifact link and confirm the final coverage and build.
\textbf{[Companion formalisation link to be inserted before submission.]}
The direct-sum formalisation and this companion development have distinct
theorem scopes; the latter remains in progress.

\paragraph{Reading guide.}
The proof order follows the three constructions rather than their chronology.
After the basic definitions, we state the repeatable amplifier and derive
the randomised result by one final sample. We then prove the amplifier's
lower bounds, replace its products by explicit lists for the deterministic
results, and give the alternating lookup construction. The appendices
supply the interlacing estimates, source conversions, implementation details
and effective parameter choices.

\section{Preliminaries}
\label{sec:models}

\subsection{Communication and estimation}

An estimator sees the whole table; a communication protocol sees one row
and one column, held by different players.

\begin{definition}\label{def:table-estimation}
A table $M:X\times Y\to\{0,1\}$ is total and its indexed input sets are
finite and nonempty. Equal rows or columns may have different indices.
Its size is $T=|X||Y|$.
A deterministic public-output protocol is a binary tree whose internal
nodes specify a speaker and a bit-valued function of that player's input
and the preceding transcript. Each leaf has a fixed output, and the
protocol must compute every entry correctly. Its cost is its maximum
depth, charging every bit, including fixed bits. Local computation is free.
Write $\D(M)$ for the least cost with unrestricted speaking order.

An additive-$a(T)$ estimator returns $\widehat D(M)$ satisfying
$|\widehat D(M)-\D(M)|\le a(T)$; its output must admit polynomial-time
comparison with polynomial-length binary rationals. A bounded-error
randomised estimator satisfies this guarantee with probability at least
$2/3$ on every input. Polynomial time means worst-case time polynomial
in $T$. The same convention applies to the alternating costs below.
\end{definition}

\begin{definition}\label{yao:model}
In strict alternation the players speak one bit at a time, in turn.
Let $D_A^{\mathrm{alt}}$ and $D_B^{\mathrm{alt}}$ be the minimum
public-output costs when Alice or Bob starts, respectively, and let
$D_{\mathrm{free}}^{\mathrm{alt}}=\min\{D_A^{\mathrm{alt}},D_B^{\mathrm{alt}}\}$.
The starter is chosen independently of the inputs. Every bit is charged,
and stopping is public: the transcript determines whether communication
has ended.

Yao's receiver-output cost $Y_{\mathrm{free}}$ permits the same choice of
starter, but the receiver of the last bit determines the answer from
that player's input and the transcript. At depth zero, the protocol
designates one player whose input determines the answer. Stopping is
still public and correctness is required on every input.
\end{definition}

\begin{definition}\label{def:reth}
Randomised sparse ETH asserts that, for some fixed constant $c_s$, no
bounded-error randomised algorithm decides $3$-SAT on $N$ variables and
at most $c_sN$ clauses in worst-case time $2^{o(N)}$. Deterministic sparse
ETH makes the corresponding assertion for deterministic algorithms.
We use this deterministic form, related to the usual formulation by
sparsification~\cite{impagliazzo2001ksat,impagliazzo2001strongly}.
\end{definition}

All communication protocols in this paper are deterministic and exact.
Randomness belongs only to an estimator or a reduction. All logarithms
are binary unless written $\ln$. Write $[q]=\{1,\ldots,q\}$ and
$\operatorname{pow}_2(a)=2^{\ceil{\log a}}$ for the least power of two
at least $a\ge1$.

\begin{observation}\label{obs:row-index-upper}
For every $n\times n$ table, $\D(M)\le\ceil{\log n}+1$: Alice sends
her row index and Bob sends the entry.
\end{observation}

\subsection{Relabelling and probability laws}

We attach positive rational probability laws to finite input sets to
measure how much of each set survives a restriction. These weights do
not change protocol cost: correctness is still required at every entry.
Repeated labels have their own weights even when their rows or columns
coincide.

\begin{observation}[Local decoding]\label{obs:local-decoding}
Suppose $F'(x',y')=F(f(x'),g(y'))$ for maps $f:X'\to X$ and $g:Y'\to Y$.
In every communication model above, the cost of $F'$ is at most that of
$F$; it is equal when both maps are onto. Also
$\rank(F')\le\rank(F)$, with equality when both maps are onto.
Transposition preserves rank and unrestricted cost, and swaps the
starters in the alternating model.
\end{observation}
\begin{proof}
Each player decodes their own input and follows the original protocol,
with the same bits, speakers, stopping rule and output convention.
For onto maps, choosing one preimage of each original label gives the
reverse simulation. Selecting or repeating rows and columns cannot
increase rank, and the same sections give equality in the onto case.
Exchanging the players proves the transposition assertions.
\end{proof}

\begin{corollary}[Square padding]\label{cor:dyadic-padding}
Let $M:X\times Y\to\{0,1\}$ have finite nonempty indexed input sets, and
let $P\ge1$ be an integer. Put
\[
S=\operatorname{pow}_2(\max\{|X|,|Y|,P\}).
\]
Repeating input labels gives an $S\times S$ table with the same cost in
every model above. If $\ell$ is the complete explicit encoding length of
$M$, including any input laws, the padded table can be printed in time
polynomial in $\ell+P$.
\end{corollary}
\begin{proof}
Choose surjections $f:[S]\to X$ and $g:[S]\to Y$ and use entries
$M(f(i),g(j))$. Local decoding preserves every cost.
Since $|X|,|Y|\le\ell$ and $S<2(\ell+P)$, printing the table takes a
fixed polynomial in $\ell+P$. The internal weights do not affect its entries.
\end{proof}

Under a decoder $f:X'\to X$, the \emph{pushforward law} assigns a set
$A\subseteq X$ the probability $\mu'(f^{-1}(A))$. Thus every retained
set $A'\subseteq X'$ has image of pushforward mass at least $\mu'(A')$.
Choosing one surviving label per image preserves the corresponding matrix
entries; it need not preserve those representatives' original weights.
When a projection is assigned its original product law below, this is an
explicit choice of law on the projected inputs.

\paragraph{Counting surviving blocks.}
We will repeatedly partition an input set into $q$ blocks of equal weight.
If a restriction retains total mass $\beta$, the sum of its within-block
masses is $q\beta$. Suppose at most $e$ blocks are exceptional and at most
$r$ other blocks retain mass at least $t$. Bounding these $r+e$ masses by
one and every remaining mass by $t$ gives
\[
q\beta\le r+e+qt.
\]
Thus a larger total mass forces more than $r$ nonexceptional blocks above
the threshold. If the blocks instead form one component of a mixture,
of total weight $\lambda$, that component retains fraction at least
$(\beta-(1-\lambda))/\lambda$: the other component can contribute at most
$1-\lambda$.

\subsection{Lookup and interlacing}
\label{sec:interlacing}

Alice holds two bits $x=x_1x_2$ and Bob holds a question $y\in\{1,2\}$;
the lookup table has entry $F(x,y)=x_y$. Give Alice a selector $i\in\{1,2\}$
and Bob two questions $(y_1,y_2)$, and ask for $F(x,y_i)$.
The selector can matter, but deleting questions can make it irrelevant.
The following restriction to $x\in\{01,10\}$ shows both effects.

\begin{figure}[H]
\centering
\begin{tikzpicture}[>=Stealth,font=\small,
 every node/.style={align=center}]
\matrix (before) at (0,0) [matrix of math nodes,
 nodes={minimum width=9mm,minimum height=6mm,inner sep=1pt},
 column 1/.style={nodes={minimum width=7mm}},
 column 2/.style={nodes={minimum width=6mm}},
 column 4/.style={nodes={fill=black!10}},
 column 5/.style={nodes={fill=black!10}},
 row 1/.style={nodes={font=\footnotesize}}] {
 x&i&(1,1)&(1,2)&(2,1)&(2,2)\\
 01&1&0&0&1&1\\
 01&2&0&1&0&1\\
 10&1&1&1&0&0\\
 10&2&1&0&1&0\\
};
\matrix (after) at (8,0) [matrix of math nodes,
 nodes={minimum width=10mm,minimum height=6mm,inner sep=1pt},
 column 1/.style={nodes={minimum width=7mm}},
 column 2/.style={nodes={minimum width=6mm}},
 row 1/.style={nodes={font=\footnotesize}}] {
 x&i&(1,1)&(2,2)\\
 01&1&0&1\\
 01&2&0&1\\
 10&1&1&0\\
 10&2&1&0\\
};
\node[above=3mm of before,font=\small\bfseries] {Entry $x_{y_i}$: selector can matter};
\node[above=3mm of after,font=\small\bfseries] {Selector no longer matters};
\draw[->,semithick] ([xshift=4mm]before.east)
 -- node[above=3mm,font=\footnotesize] {keep $y_1=y_2$}
    ([xshift=-4mm]after.west);
\draw (before-1-1.south west) -- (before-1-6.south east);
\draw (after-1-1.south west) -- (after-1-4.south east);
\draw[black!35] (before-3-1.south west) -- (before-3-6.south east);
\draw[black!35] (after-3-1.south west) -- (after-3-4.south east);
\draw[draw=answergreen,rounded corners=2pt]
 ([xshift=-1mm,yshift=-.6mm]after-2-3.north west) rectangle
 ([xshift=1mm,yshift=.6mm]after-3-4.south east);
\draw[draw=answergreen,rounded corners=2pt]
 ([xshift=-1mm,yshift=-.6mm]after-4-3.north west) rectangle
 ([xshift=1mm,yshift=.6mm]after-5-4.south east);
\node[right=3mm of after,align=left,font=\footnotesize] {equal rows\\within each\\fixed $x$};
\node[below=3mm of before,font=\footnotesize] {Shaded columns are deleted.};
\end{tikzpicture}
\caption{Deleting the unequal-question columns erases the effect of Alice's
selector $i$. In each outlined pair, changing $i$ leaves every retained entry
unchanged. Column labels are Bob's question pairs $(y_1,y_2)$.}
\label{fig:lookup}
\end{figure}

\begin{observation}\label{ex:lookup-deletion}
In the equal-question columns of \cref{fig:lookup}, the answer is independent
of Alice's selector, since $F(x,y_1)=F(x,y_2)$ whenever $y_1=y_2$.
\end{observation}

The general selector construction is the interlacing operation of
Mackenzie and Saffidine~\cite{mackenzie2025refuting}.

\begin{definition}[Interlacing]\label{def:interlace}
For a Boolean matrix $R:X\times Y\to\{0,1\}$ with specified input laws,
the full $q$-fold interlace $\Int{q}{R}$ has rows $(i,u)\in[q]\times X$,
columns $(y_1,\ldots,y_q)\in Y^q$, and entry
\[
\Int{q}{R}((i,u),(y_1,\ldots,y_q))=R(u,y_i).
\]
The row law is uniform on $[q]$ times the base row law; the column law is
the product of the base column law.
\end{definition}

\begin{definition}[Exact tuple lists]\label{def:exact-tuples}
Fix integers $q,D\ge1$. For a finite alphabet $Y$ with a prescribed
rational law $\mu$,
an indexed tuple list consists of a finite nonempty label set $\mathcal L$,
a positive rational law $\nu$ on its labels, and a decoder
$z:\mathcal L\to Y^q$. It is \emph{exactly $D$-wise independent} if every
set $J\subseteq[q]$ of at most $D$ decoded coordinates has law $\mu^J$.
Exactness refers to equality of these small marginals; higher-order
independence is not required. Events are sets of labels and may distinguish
labels with the same decoded tuple.

For a base matrix $R:X\times Y\to\{0,1\}$ with column law $\mu$,
interlacing on this list gives column $\ell$ the entry $R(u,z_i(\ell))$
at row $(i,u)$, with the row law of \cref{def:interlace} and column law $\nu$.
Within-block row mass is measured in the base row law, without division by
$q$. Distinct labels remain distinct inputs even if their decoded tuples agree.
\end{definition}

\section{Amplification and the randomized result}
\label{sec:unrestricted}

The main construction turns a one-bit gap into an arbitrary gap while keeping
it stable under deletion of inputs. We state its interface first and derive
the randomized hardness result; the construction and its proof follow in
\cref{sec:amplification-proof}.

\subsection{Lower bounds that survive restrictions}

The restriction profile records how much of a lower bound remains after
inputs are removed. We need only two column thresholds: one for the
current restriction and one after a single bit halves its column mass.

\begin{definition}[Restriction profile]
\label{def:profile}
Let \(F:X\times Y\to\{0,1\}\) have positive rational input laws. For
\(0<\rho\le1\) and an integer \(d\), the restriction profile \(\RP{F}{\rho}{d}\)
holds if, for every \(j\in\{0,1\}\) and all \(U\subseteq X\), \(V\subseteq
Y\),
\[
\mu_X(U)\ge\rho\ \text{and}\ \mu_Y(V)\ge(3/5)2^{-j} \quad\Longrightarrow\quad
\D(F[U,V])\ge d-j.
\]
The two column thresholds are $3/5$ and $3/10$: halving the retained column
mass may lose one bit of the lower bound. No matching upper bound is required.
\end{definition}

We next state the bounded-degree colouring source used by the construction. It
supplies the activity pattern, while the restriction profile carries the
communication lower bound to the next round. We use the bounded-degree colouring
promise below; its conversion from sparse \(3\)-SAT, including the size
accounting, is proved in \cref{sec:sources} and uses the standard sparse gap and
bounded-degree reductions~\cite{dinur2007pcp,bogdanov2002lower}.

\begin{theorem}[Bounded-degree colouring source]
\label{src:gap-colouring}
There are constants \(\Delta\ge1\) and rational \(\gamma>0\), and a
deterministic polynomial-time reduction from \(3\)-SAT to nonempty ordered
loopless multigraphs \(G=(V,E)\) with no isolated vertices and maximum degree at
most \(\Delta\), such that:
\begin{enumerate}
\item a satisfiable formula produces a three-colourable graph;
\item an unsatisfiable formula produces a graph for which every three-colouring has at least \(\gamma|E|\) monochromatic edge occurrences.
\end{enumerate}
For sparse formulas on \(N\ge1\) variables, the reduction can use a fixed
positive integer \(p\) such that, for \(n=|V|\),
\[
N\le |E|,\qquad n,|E|\le N(\log(N+2))^p,\qquad n/2\le |E|\le\Delta n/2.
\]
\end{theorem}

The amplifier combines this graph promise with a restriction profile; when the
base cost is at least three, it restores that profile while charging one
additional bit on a NO instance.

\begin{theorem}[Deletion-robust amplification]
\label{thm:amplifier}
Fix the graph promise of \cref{src:gap-colouring} and write \(L=\ceil{\log n}\).
For every source larger than an effectively computable constant, there is a
public integer \(h=L+O_{\Delta,\gamma}(\log L)\) and a deterministic, implicitly
represented construction \(R\mapsto R^+\) with the following properties:
\begin{enumerate}
\item[(i)] For every source,
\[
\D(R^+)\le\D(R)+h+1,
\]
and, if \(G\) is YES, the stronger bound \(\D(R^+)\le\D(R)+h\) holds.
\item[(ii)] If \(G\) is NO, \(\D(R)=d\ge3\), and \(\RP{R}{2^{-L}}{d}\) holds, then
\[
\D(R^+)=d+h+1,\qquad \RP{R^+}{2^{-L}}{d+h+1}.
\]
\item[(iii)] On either source promise, \(\rank(R^+)\le2^h\rank(R)\).
\end{enumerate}
All intermediate products are complete. When iterated from the lookup seed
below, the construction has the recursive input generators and entry evaluator,
with their uniform resource bounds, proved in \cref{lem:sample-access}. Its parameter
choices and admissibility test are those of \cref{lem:full-schedule}; they are
independent of the base matrix and of the number of rounds.
\end{theorem}

The guarantee is conditional on the graph promise and restriction profile, so it
does not apply to an arbitrary base matrix. The deletion argument, including the
part that handles whichever player speaks first, is given in
\cref{sec:amplification-proof}; the interlacing estimates used there are proved
in \cref{sec:interlacing-details}.


\subsection{Iteration}

One round adds $h$ bits to the YES upper bound and $h+1$ bits to the NO
lower bound. Starting from a lookup table makes the accumulated gap explicit.

\begin{lemma}[Lookup seed]\label{cl:seed}
For $L\ge1$, let $p_0=\operatorname{pow}_2(64(L+1))$ and
$R_0(z,j)=z_j$ on $\{0,1\}^{p_0}\times[p_0]$, with uniform laws. Then
\[
d_0:=\log p_0+1,\qquad \D(R_0)=d_0,\qquad
\rank(R_0)=p_0=2^{d_0-1},\qquad \RP{R_0}{2^{-L}}{d_0}.
\]
For $L=\ceil{\log n}$ and $n\ge4$, the seed has polynomial size in $n$
and can be printed in polynomial time.
\end{lemma}
The proof is in \cref{sec:interlacing_details}.

\begin{proposition}\label{prop:iteration}
Iterate \cref{thm:amplifier} from $R_0$ using the same admissible graph and
schedule. Put $B_i=d_0+ih$, $H_i=B_i+i$, and $r_i=2^{B_i-1}$. A YES graph
gives $\D(R_i)\le B_i$. A NO graph gives
$\D(R_i)=H_i$ and $\RP{R_i}{2^{-L}}{H_i}$. On either promise,
$\rank(R_i)\le r_i$ and $\D(R_i)\le H_i$.
\end{proposition}
\begin{proof}
The seed supplies the case $i=0$. At each round the two cost bounds increase
by $h$ and $h+1$, respectively, while rank grows by at most $2^h$.
The NO profile is restored at every round, so the induction continues.
\end{proof}

\section{Final sampling and growing error}
\label{sec:quantitative}

The final sampled table keeps the absolute communication threshold from the
full-product construction while reducing the printed size to a rank-sensitive
sample.

\subsection{Counting columns served by short protocols}

We first identify the objects that the sampling union bound must count. For a
target threshold $H$, making the final table hard requires every depth-$(H-1)$
protocol to miss a sampled column, so the count must be uniform over retained
row sets. The right objects to count are therefore the column sets that one
protocol with Bob's maps unspecified can serve on one retained row set.

\begin{definition}[Compatible-column set]\label{def:compatibility}
A \emph{partial protocol tree} fixes Alice's behaviour and the leaf outputs
while leaving Bob's maps to be chosen separately for each column. More
precisely, a depth-$t$ protocol tree $P$ specifies a complete binary tree of
depth $t$, the speaker at each internal node, all Alice maps, and all leaf
outputs; Bob's maps are unspecified. For a Boolean matrix
$F:X\times Y\to\{0,1\}$ and $U\subseteq X$, let $G_P(U)$ be the set of columns
for which some assignment of bits to Bob's nodes makes the tree correct on every
row in $U$. The assignment may depend on the column, but not on Alice's row.
\end{definition}

This definition turns the restriction question into a containment question: a
short protocol on a restriction is exactly one such tree whose compatible set
contains all retained columns.

\begin{observation}\label{obs:compatible-protocol}
For retained sets $U,V$, with correctness on an empty restriction understood
vacuously,
\begin{equation}\label{eq:compatibility}
 \D(F[U,V])\le t
 \quad\Longleftrightarrow\quad
 V\subseteq G_P(U)\text{ for some depth-}t\text{ protocol tree }P\text{ of this kind}.
\end{equation}
\end{observation}

\begin{proof}
Pad a protocol of depth at most $t$ with fixed dummy messages and extend Alice's
maps to all rows. Its Bob maps witness the required containment. Conversely,
for each retained column choose a witnessing assignment at Bob's nodes; these
assignments define Bob's maps for a correct protocol on $U\times V$. Empty
restrictions satisfy both assertions vacuously.
\end{proof}

The two one-bit trees in \cref{fig:compatible-trees} use the same small lookup
table. They show why leaf requirements combine by intersection at an Alice
node and by union at a Bob node: Bob's choice can depend on the column, but
must work for both retained rows.

\begin{figure}[H]
\centering
{\small
One lookup table for both trees: $F(x,y)=x_y$.

\smallskip
\begin{tabular}{c@{\qquad}*{3}{>{\centering\arraybackslash}p{8mm}}}
\toprule
&\multicolumn{3}{c}{Bob's question $y$}\\
\cmidrule(l){2-4}
Alice's row $x$&$1$&$2$&$3$\\
\midrule
$001$&$0$&$0$&$1$\\
$011$&$0$&$1$&$1$\\
\bottomrule
\end{tabular}
\par}
\medskip
\captionsetup[subfigure]{justification=centering}
\begin{subfigure}[t]{.48\linewidth}
\centering
\begin{tikzpicture}[font=\small,>=Stealth,
 box/.style={draw,rounded corners=2pt,align=center,inner sep=4pt},
 leaf/.style={box,text width=2.55cm}]
\node[box,draw=aliceblue] (a) at (0,0) {Alice sends $x_2$};
\node[leaf,anchor=north] (a0) at (-1.65,-1.55) {row $001$\\output $0$\\columns $\{1,2\}$};
\node[leaf,anchor=north] (a1) at (1.65,-1.55) {row $011$\\output $1$\\columns $\{2,3\}$};
\coordinate (am) at (0,-.95);
\coordinate (al) at (-1.65,-.95);
\coordinate (ar) at (1.65,-.95);
\draw[aliceblue] (a.south) -- (am);
\draw[aliceblue] (am) -- node[midway,above=3pt,inner sep=1pt] {$0$} (al);
\draw[aliceblue] (am) -- node[midway,above=3pt,inner sep=1pt] {$1$} (ar);
\draw[->,aliceblue] (al) -- (a0.north);
\draw[->,aliceblue] (ar) -- (a1.north);
\end{tikzpicture}
\caption{\textbf{Both required: intersection.}\\
$\{1,2\}\cap\{2,3\}=\{2\}$.}
\end{subfigure}\hfill
\begin{subfigure}[t]{.48\linewidth}
\centering
\begin{tikzpicture}[font=\small,>=Stealth,
 box/.style={draw,rounded corners=2pt,align=center,inner sep=4pt},
 leaf/.style={box,text width=2.55cm}]
\node[box,draw=boborange] (b) at (0,0) {Bob chooses using $y$};
\node[leaf,anchor=north] (b0) at (-1.65,-1.55) {rows $001,011$\\output $0$\\columns $\{1\}$};
\node[leaf,anchor=north] (b1) at (1.65,-1.55) {rows $001,011$\\output $1$\\columns $\{3\}$};
\coordinate (bm) at (0,-.95);
\coordinate (bl) at (-1.65,-.95);
\coordinate (br) at (1.65,-.95);
\draw[boborange] (b.south) -- (bm);
\draw[boborange] (bm) -- node[midway,above=3pt,inner sep=1pt] {$0$} (bl);
\draw[boborange] (bm) -- node[midway,above=3pt,inner sep=1pt] {$1$} (br);
\draw[->,boborange] (bl) -- (b0.north);
\draw[->,boborange] (br) -- (b1.north);
\end{tikzpicture}
\caption{\textbf{Either works: union.}\\
$\{1\}\cup\{3\}=\{1,3\}$.}
\end{subfigure}
\caption{Alice's branch depends on the row, so a compatible column must
satisfy both leaf conditions. Bob's branch depends on the column, so either
leaf may work, but his chosen leaf must answer both rows correctly.}
\label{fig:compatible-trees}
\end{figure}

The count below uses the row-sum argument of Oum, S{\ae}ther and
Vatshelle~\cite[Theorem~4.2]{oum2014faster}. Counting sums, rather than the rows
that produce them, is what keeps the final sample small.

\begin{theorem}[Compatible-column count]\label{thm:rank-count}
Let $r,t$ be nonnegative integers. If a Boolean matrix $F$ has real rank at most
$r$, then, over all depth-at-most-$t$ partial protocol trees and all retained
row subsets, there are at most
\begin{equation}\label{eq:rank-count}
 2^{(r+1)\ceil{\log(r+2)}\,2^{t+1}}
\end{equation}
distinct compatible-column sets. Repeated row and column labels are allowed.
\end{theorem}
\begin{proof}
Stack the rows of $F$ and $\mathbf1-F$ into one Boolean matrix $A$. Its row
space lies in the span of the rows of $F$ and the all-one vector, so its rank
is at most $r+1$. Fix a set of basis columns of $A$. For any subset of its
rows, choose a basis of their span and sum those basis rows. Nonnegativity
makes the zero coordinates of this sum exactly the columns vanishing on the
chosen row subset. The sum is determined by its basis-column coordinates,
each an integer between $0$ and $r+1$. Hence there are at most
$(r+2)^{r+1}$ common-zero sets. The zero sum includes empty row subsets.

Both leaf outputs are covered: output $0$ asks for a common-zero set of rows
of $F$, and output $1$ asks for one of their complements. Write
$a=(r+1)\ceil{\log(r+2)}\ge1$, so there are at most $2^a$ leaf conditions.
At an Alice node the child conditions intersect; at a Bob node they unite.
These operations use the rows consistent with Alice's messages on each path,
and therefore include every retained row set and every choice of Alice's maps.
After padding to depth $t$, the count is at most
\[
2^{2^t-1}(2^a)^{2^t}\le2^{a2^{t+1}}.
\]
The argument also covers rank zero and repeated row or column labels.
\end{proof}

\subsection{Sampling columns and then rows}
\label{sec:sampling}

We sample the same number of rows and columns and round that number up to a
power of two. This gives the required square table directly. The proof exposes
the columns first: their guarantee must hold for every sufficiently heavy row
set before the independent row sample is drawn.

\begin{theorem}[Two-sided compression]\label{thm:compression}
Let $F:X\times Y\to\{0,1\}$ have real rank at most $r$, where $r\ge0$ is an
integer, and let $\mu_X,\mu_Y$ be input laws on nonempty finite indexed sets.
Suppose an integer $H\ge1$ satisfies
\[
\mu_X(U)\ge3/5,\quad\mu_Y(V)\ge3/5
\quad\Longrightarrow\quad \D(F[U,V])\ge H.
\]
Put
\[
m=\operatorname{pow}_2\bigl(16(r+1)\ceil{\log(r+2)}\,2^H\bigr),
\]
where $\operatorname{pow}_2(u)$ is the least power of two at least $u$.
Draw $m$ rows and $m$ columns independently, retaining separate indices for
repeated draws. The resulting $m\times m$ matrix has cost at least $H$ except
with probability less than $1/100$. Every protocol for $F$ works on every
sample outcome.
\end{theorem}
\begin{proof}
Repeated labels can always be decoded locally; conversely, a protocol on
repeated labels can be restricted to one occurrence of each label. We use this
observation throughout.

For every $U$ with $\mu_X(U)\ge3/5$ and every depth-$(H-1)$ partial protocol
$P$, the compatible set $G_P(U)$ has mass less than $3/5$: otherwise
\cref{obs:compatible-protocol} contradicts the restriction hypothesis. The
\emph{whole family} of these sets, as both $U$ and $P$ vary, has size at most
$2^{m/16}$ by \cref{thm:rank-count}. Thus the probability that all sampled
columns lie in any one of them is at most $2^{m/16}(3/5)^m$.
Outside this event, the sampled-column matrix
$F_Y(x,j)=F(x,y_j)$ has cost at least $H$ on every original row set of mass at
least $3/5$.

Condition on such a column outcome. The transpose $F_Y^{\mathsf T}$ still has
rank at most $r$. Every compatible set of its columns, using all its rows,
has $\mu_X$-mass less than $3/5$, or transposing the resulting protocol would
contradict the preceding guarantee. There are again at most $2^{m/16}$ such sets.
The row draws remain independent with law $\mu_X$, so the probability that
they all lie in one of these sets is at most $2^{m/16}(3/5)^m$. Any short protocol
on the final matrix would give precisely such a containment. Since $m\ge32$,
the total failure probability is at most
\[
2\cdot2^{m/16}(3/5)^m\le2^{1-7m/16}\le2^{-13}<1/100.
\]
The universal upper bound follows by decoding sampled indices and running a
protocol for $F$.
\end{proof}

Compression is used only after the full-product amplification is complete.
It preserves the final threshold; no further amplification of the sample is
needed. The probability bound does not give a test for whether a particular
sample succeeds.

\subsection{Printing the sample and comparing errors}
\label{sec:reduction}
\label{sec:randomised_reduction}

The full-product table is too large to print. We instead generate complete
labels for the sampled inputs, then evaluate only their intersections.

Take a source graph on $n$ vertices that passes the computable admissibility
test in \cref{lem:full-schedule}. Every sufficiently large source does so.
Use its schedule at every round and write
\[
L=\ceil{\log n},\qquad B_k=d_0+kh,\qquad H_k=B_k+k,\qquad r_k=2^{B_k-1}
\]
as in \cref{prop:iteration}. The known bound $r_k$ suffices; the reduction
never computes the rank of the implicit table.

\begin{lemma}\label{lem:sample-access}
For the level-$k$ iteration from the lookup seed in \cref{prop:iteration}, use
$r_k,H_k$ in \cref{thm:compression}. A joint sampler and entry evaluator use $n^{O(k+1)}$
time and random bits on every tape, with an absolute implied constant
independent of $k$. Every output label is an actual input of $R_k$, and the
sample counts and tape length are fixed. The joint law is within total
variation less than $1/100$ of the ideal independent sample.
\end{lemma}

Here total variation is the maximum difference in probabilities assigned to an
event. The recursive construction and its fixed-bit implementation are proved
in \cref{app:implementation}. Sampling is approximate; every printed entry is
exact.

\begin{theorem}[Uniform randomised reduction]\label{thm:uniform-reduction}
One randomised algorithm, given an admissible source graph on $n$ vertices and
an integer $k\ge1$, outputs a total Boolean matrix $M$ and an integer $B$.
The table is square with power-of-two side length. On every tape the algorithm
uses $n^{O(k+1)}$ time and random bits, with an absolute implied
constant independent of $k$. For a YES source, $\D(M)\le B$ on every tape.
For a NO source,
\[
\Pr[\D(M)=B+k]\ge49/50.
\]
Here $B=B_k$ and the tape-independent entry count is
\[
m=\operatorname{pow}_2\bigl(16(r_k+1)\ceil{\log(r_k+2)}\,2^{H_k}\bigr),
\qquad T=m^2.
\]
\end{theorem}

\begin{algorithm}[H]
\caption{Printing the sampled table and its threshold}
\begin{algorithmic}[1]
\Require An admissible source graph $G$ and an integer $k\ge1$.
\State Compute its seed and schedule using \cref{lem:full-schedule}.
\State Compute $B_k,H_k,r_k,m$ and reserve the fixed random tape.
\State Generate $m$ columns and $m$ rows using \cref{lem:sample-access}.
\State Evaluate all stored row--column pairs, keeping repeated labels as distinct indices.
\State \Return $(M,B_k)$.
\end{algorithmic}
\end{algorithm}

\begin{proof}
The resource and dimension guarantees follow from \cref{lem:sample-access}.
On a YES source, local decoding transfers the bound $\D(R_k)\le B_k$ to
every sample, since all generated labels are legal.

On a NO source, \cref{prop:iteration} gives
\[
\D(R_k)=H_k,\qquad \RP{R_k}{2^{-L}}{H_k},\qquad \rank(R_k)\le r_k.
\]
Because $2^{-L}\le1/2$, the $j=0$ profile bound applies whenever both retained
masses are at least $3/5$. Thus \cref{thm:compression} gives ideal sampling
failure less than $1/100$. Replacing its law by the implemented sampler adds
less than $1/100$, by \cref{lem:sample-access}. Except with probability less
than $1/50$, the output therefore has cost at least $H_k$; local decoding gives
cost at most $H_k$ on every tape. Hence its cost is exactly $B_k+k$.
\end{proof}

\begin{proof}[Proof of \cref{thm:quantitative}]
Suppose a bounded-error randomised polynomial-time estimator achieves additive
error $a(T)=c\log T/\log\log T$ for some fixed $0\le c<1/8$ on every
sufficiently large square table of power-of-two side length. We obtain a
bounded-error $2^{o(N)}$ algorithm for sparse $3$-SAT.

For a formula with $N$ variables, \cref{src:gap-colouring} produces a source
graph with $n$ vertices, where
\[
2N/\Delta\le n\le N(\log(N+2))^p
\]
for a fixed $p$. The lower bound follows from $N\le|E|\le\Delta n/2$.
Thus the graph is admissible for every sufficiently large $N$, and
\[
L=\ceil{\log n}=\log N+O(\log\log N).
\]
Choose $k=\floor{N/\ceil{\log N}^{2}}$. By
\cref{lem:full-schedule,cl:seed}, $h=L+O(\log L)$ and $d_0=O(\log L)$, so
\[
B_k=d_0+kh=(1+o(1))k\log N.
\]
Substituting $r_k=2^{B_k-1}$ and $H_k=B_k+k$ into
\cref{thm:uniform-reduction} gives
$\log m=2B_k+k+O(\log(B_k+1))$. Hence the printed size satisfies
\begin{align}
\log T&=2\log m=4B_k+2k+O(\log(B_k+1))
             =(4+o(1))k\log N,\label{eq:quantitative-size}\\
\log\log T&=(1+o(1))\log N,
\qquad k=\left(\frac14+o(1)\right)\frac{\log T}{\log\log T}.
\end{align}
Consequently $k>2a(T)$ for all sufficiently large $N$. The table already has
the square power-of-two form required by the assumed estimator.

For each estimator run, test whether its output is below the rational midpoint
$B_k+k/2$. An accurate estimate answers this test correctly on every YES output
and every successful NO output. Repeat independently a constant number of times
and take the majority answer, reducing the decision error on each such table
below $1/100$. This uses only the rational-threshold comparison required by
the estimator model. The reduction
fails with probability at most $1/50$, so the decision error is less than
$3/100<1/3$.

Both construction time and random-bit usage are
$n^{O(k+1)}=2^{o(N)}$ on every tape. The estimator also takes
$2^{o(N)}$ time because
$\log T=O(N/\log N)$. For all sufficiently large $N$, the chosen $k$ is
positive, the graph is admissible, and $T$ meets the estimator's size threshold.
The remaining finitely many input sizes can be decided exhaustively within
a fixed bound. This contradicts randomised sparse
ETH. In particular, an estimator working on all explicit total tables is also
impossible.
\end{proof}

\section{Proof of robust amplification}
\label{sec:amplification-proof}

\subsection{Encoding the graph}

We combine the graph's activity pattern with lookup in a base matrix.
The widths \(s,q_2,w,Q\) remain symbolic here; \cref{lem:full-schedule}
fixes their public dyadic values and admissibility tests.

In a YES instance, the players will reduce the new lookup to one lookup in
$R$ using three messages: Alice names a bin, Bob reveals its activity
coordinate, and Alice selects the relevant base question. A proper colouring
ensures that an active item is unique within its bin. In a NO instance, some
bin must contain a collision, and the rest of the proof shows that the collision
cannot be avoided cheaply. The template rows below force the first messages;
the source rows carry the graph constraint.

\begin{definition}[Source construction]
\label{def:source-matrix}
Let \(R:X\times Y\to\{0,1\}\) have positive rational input laws, and let
\(s\ge4\) be dyadic. Integer coordinates below are uniform; products use
the specified factor laws.

\smallskip\noindent\textit{Activity.}
An item is a vertex or auxiliary label in one of \(4s\) groups:
\(\mathcal I=[4s]\times(V\sqcup[s-3])\).
In each group, each edge occurrence has a coordinate activating its two
endpoints and all \(s-3\) auxiliary items. All other items are inactive.
Pad these \(4s|E|\) coordinates with inactive coordinates to length
\(q_0=\operatorname{pow}_2(4s|E|)\), then repeat each coordinate
\(q_2/q_0\) times, where \(q_2\) is a dyadic multiple of \(q_0\).
Write \(v_i(\alpha)\) for item \(i\)'s activity at coordinate \(\alpha\),
and \(\pi_\alpha\) for the order-preserving injection of active items into
\([s]\). Activity does not depend on any assignment of items to bins.

\smallskip\noindent\textit{Template lookup.}
Put \(q_1=s-2\), leaving two codes for active and inactive source items.
Define
\[
\Gamma=Y^{2s-1},\qquad J=[q_1]\times X,\qquad S=J^{q_2},\qquad
J_2=[q_2]\times\Gamma.
\]
The first \(q_1\) positions of a word \(\gamma\in\Gamma\) answer
template questions; its last \(s+1\) positions will answer source questions.
The two lookup matrices are
\[
A(\gamma,(r,u))=R(u,\gamma_r),\qquad
R_{\rm temp}(z,(\alpha,\gamma))=A(\gamma,z_\alpha).
\]

\smallskip\noindent\textit{Raw entries.}
The raw matrix has template rows \((p,z)\) and source rows \((i,u)\):
\[
\mathcal K=([s]\times S)\sqcup(\mathcal I\times X).
\]
Set \(\lambda=255/256\) and \(\mu=1/256\).
Template rows have total mass \(\lambda\); vertex rows and auxiliary rows
each have mass \(\mu/2=1/512\). Within a type, bin or item labels are uniform
and the other labels follow their specified laws.
For a column \(Z=(\xi,j_1,\ldots,j_s)\), with \(\xi\in Y^{s+1}\) and
\(j_p=(\alpha_p,\gamma_p)\in J_2\), set
\[
B_{\rm raw}((p,z),Z)=R_{\rm temp}(z,j_p),
\]
\[
B_{\rm raw}((i,u),Z)=
 \begin{cases}
 R(u,\xi_{\pi_\alpha(i)}),
   &\alpha_1=\cdots=\alpha_s=\alpha,\ v_i(\alpha)=1,\\
 R(u,\xi_{s+1}),&\text{otherwise}.
 \end{cases}
\]
Thus a source row uses its active position when all selectors agree, and
the common inactive position otherwise.

\smallskip\noindent\textit{Two column families.}
Ordinary columns \(\Omega=Y^{s+1}\times J_2^s\) have their product law.
Diagonal columns are
\[
\mathcal D = \{(\operatorname{tail}(\gamma),j,\ldots,j):j=(\alpha,\gamma)\in
J_2\}, \qquad \operatorname{tail}(\gamma)=(\gamma_{q_1+1},\ldots,\gamma_{2s-1}),
\]
with the pushforward of the \(J_2\) law. Give these families the masses
\[
\begin{array}{ccl}
\text{Family}&\text{Mass}&\text{Role}\\
\Omega&\lambda&\text{Reference inputs and the initial selector prefix}\\
\mathcal D&\mu&\text{Collision witnesses and protected inputs}.
\end{array}
\]
In particular, the ordinary-column mass equals the template-row mass
\(\lambda\). Keep the two family tags distinct even when their raw columns
coincide. At every bin \(p\), the label \(j_p\) has the \(J_2\) law in
either family.

\smallskip\noindent\textit{The two interlaces.}
Transpose the raw matrix and name its reference restriction:
\[
N=B_{\rm raw}^{\mathsf T},\qquad
N_{\rm ref}=N[\Omega,\mathcal K],
\]
where the \(\Omega\) law is normalised. Finally set
\[
U=\Int{w}{N},\qquad W=U^{\mathsf T},\qquad
R^+=(\Int{Q}{W})^{\mathsf T}.
\]
\end{definition}

The ordinary inputs force the protocol to assign items to bins; the diagonal
inputs then expose a collision through distinct positions of one base word.
We call the diagonal inputs \emph{protected inputs}. The internal interlace
tracks their destinations under arbitrary speaking orders. On a NO instance,
the argument gives its transpose \(W\) the required full cost, with a
restriction profile still one bit below that cost. The final interlace
restores the profile at the new cost, so the next round can begin.

The following \(s=4\) example compares two bin assignments with the same
activity. The endpoints \(u,v\) and auxiliary item \(a\) are active; item
\(h\), from another group, is inactive. Take \(u,v\) to have active ranks
\(1,2\). In the diagonal family they request \(\gamma_3,\gamma_4\), while
\(h\) requests \(\xi_{s+1}=\gamma_7\).

\begin{figure}[H]
\centering
\begin{tikzpicture}[>=Stealth,font=\small,
 bin/.style={draw,rounded corners=2pt,minimum height=13mm},
 active/.style={circle,draw=answergreen,fill=answergreen!15,
   minimum size=5.5mm,inner sep=0pt},
 inactive/.style={rectangle,draw=boborange,dashed,minimum size=5.5mm,
   inner sep=0pt},
 target/.style={draw,rounded corners=2pt,fill=answergreen!4,
   minimum width=9mm,minimum height=7mm,inner sep=2pt}]
\node[font=\small\bfseries] at (2.2,1.55) {Assignment A};
\node[font=\small\bfseries] at (9.6,1.55) {Assignment B};
\foreach \x/\b/\wide in {0/1/22,1.85/2/12,3.3/3/12,4.75/4/12} {
  \node[bin,minimum width=\wide mm] at (\x,0) {};
  \node[font=\footnotesize] at (\x,.95) {bin \b};
  \node[bin,minimum width=\wide mm] at (\x+7.4,0) {};
  \node[font=\footnotesize] at (\x+7.4,.95) {bin \b};
}
\node[active] at (0,0) {$u$};
\node[active] at (1.85,0) {$v$};
\node[inactive] at (3.3,0) {$h$};
\node[active] at (4.75,0) {$a$};
\node[active] (u) at (6.8,0) {$u$};
\node[active] (v) at (7.4,0) {$v$};
\node[inactive] (h) at (8,0) {$h$};
\node[active] at (12.15,0) {$a$};
\node[target] (gu) at (6.5,-1.8) {$\gamma_3$};
\node[target] (gv) at (7.4,-1.8) {$\gamma_4$};
\node[target] (gh) at (8.3,-1.8) {$\gamma_7$};
\draw[->,answergreen] (u.south) -- (gu.north);
\draw[->,answergreen] (v.south) -- (gv.north);
\draw[->,boborange,dashed] (h.south) -- (gh.north);
\node[align=center,anchor=north] at (7.4,-2.35)
 {One bin, three distinct\\lookup positions};
\node[align=left,text width=57mm,anchor=north west,font=\footnotesize]
 at (-.95,-1.05)
 {Filled circles: active $u,v,a$.\\
  Dashed square: inactive $h$.\\[3pt]
  Activity stays fixed.};
\end{tikzpicture}
\caption{The same active and inactive items can have different bin assignments.
In assignment B, one bin contains the active endpoints \(u,v\) and the inactive
item \(h\). Their diagonal inputs request three distinct positions.
These are two illustrative assignments, not steps of a protocol.}
\label{fig:source-collision}
\end{figure}

The collision in \cref{fig:source-collision} is the local witness after a
protocol has fixed an assignment of items to bins. The lemma below makes that
witness uniform over the allowed deletions, which is the graph-specific input
needed by the amplification proof.

\begin{lemma}
\label{src:item-witness}
A YES graph admits an assignment of \(\mathcal I\) to \([s]\) with at most one
active item in each bin at every coordinate. For a NO graph, delete fewer than
\((s/2)\delta_v n\) vertex items and fewer than \(s/2\) auxiliary items, where
\[
\delta_v=\frac{\gamma|E|}{2\Delta n},\qquad \zeta=\frac{\gamma}{24}.
\]
Every assignment of the remaining items to the \(s\) bins then has at least
\(\zeta q_2\) coordinates at which one bin contains two active items and an
inactive item.
\end{lemma}

The same bin assignment gives a short protocol on YES instances.

\begin{proposition}
\label{src:upper}
If \(\D(R)\le y\), then \(\D(N)\le y+2\log s+\log q_2+1\). For a YES graph, the
final \(+1\) can be omitted. Both bounds remain valid after restricting either
input set.
\end{proposition}

\begin{proof}
For a YES graph, fix a public bin assignment $\chi$ from
\cref{src:item-witness}. Alice holds a template row $(p,z)$ or a source row
$(i,u)$; Bob holds $Z=(\xi,j_1,\ldots,j_s)$ with
$j_p=(\alpha_p,\gamma_p)$. They exchange:
\[
\begin{array}{lll}
\text{Alice} & \text{bin }p\text{ (use }p=\chi(i)\text{ for a source row)}
  & \log s\text{ bits},\\
\text{Bob} & \alpha_p & \log q_2\text{ bits},\\
\text{Alice} & \text{one of }q_1\text{ template codes, active, inactive}
  & \log s\text{ bits}.
\end{array}
\]
The third message has $q_1+2=s$ possibilities. For a template row, write
$z_{\alpha_p}=(r,u)$ and send code $r$; the answer is
$R(u,(\gamma_p)_r)$. For a source row send active precisely when
$v_i(\alpha_p)=1$. If all $\alpha_j$ agree and the code is active, Bob
identifies the unique active item in bin $p$ and uses
$\xi_{\pi_{\alpha_p}(i)}$. Otherwise he uses $\xi_{s+1}$.
Bob can test agreement from his own input. The entry rule therefore reduces
both row types to one execution of a cost-$y$ protocol for $R$.

For an arbitrary graph, use any public assignment and replace the active
code by the $s$ possible values of $\pi_{\alpha_p}(i)$. There are now
$q_1+s+1=2s-1$ codes, so the third message costs one additional bit.
The same rules work for every column tag. Exchanging the players gives the
protocol for $N=B_{\rm raw}^{\mathsf T}$, and restriction preserves its cost.
\end{proof}

This protocol sets the YES baseline. We next prove that collisions force the
extra bit and that the two interlaces preserve it.

\subsection{A collision survives the conversation}

The graph promise must apply to one assignment fixed by the protocol before an
activity coordinate is selected. We use the items and matrices of
\cref{def:source-matrix}, with the widths fixed by \cref{lem:full-schedule}. The
combinatorial count in \cref{src:item-witness} applies to every assignment of
the retained items.

\begin{proof}[Proof of \cref{src:item-witness}]
For a YES graph, let \(c:V\to[3]\) be a proper three-colouring supplied by
\cref{src:gap-colouring}. Assign each vertex \(v\) to bin \(c(v)\), and assign
the \(s-3\) auxiliary labels to the remaining bins. Every real coordinate
activates an edge's two endpoints and those auxiliaries in one group, so all
its active items occupy distinct bins. Padded coordinates activate no items.

Now suppose that \(G\) is NO. By the deletion hypotheses in
\cref{src:item-witness}, fewer than \(s/2\) groups lose an auxiliary item and
fewer than \(s/2\) groups lose more than \(\delta_v n\) vertices. Hence more
than \(3s\) groups retain every auxiliary item and all but at most \(\delta_v
n\) vertices.

Fix such a group. If two auxiliaries share a bin, they collide at every
edge coordinate. Otherwise the auxiliaries occupy \(s-3\) distinct bins,
leaving three free bins. Colour the retained vertices in those free bins by
their bins, and extend the colouring arbitrarily to all other vertices.
At least \(\gamma|E|\) edge occurrences are monochromatic by the NO promise,
and at most \(\Delta\delta_v n=\gamma|E|/2\) meet a deleted vertex.
Every remaining monochromatic edge gives an active collision: an endpoint
in an auxiliary bin collides with its auxiliary; if neither endpoint is in
an auxiliary bin, the two endpoints share a free bin. Thus this group has
at least \(\gamma|E|/2\) edge coordinates with a collision.

For each counted collision, choose one bin containing the two active items.
Discard a group from the count if one of its chosen bins contains retained
items from that group alone. Every discarded group has an exclusive chosen bin, and two groups
cannot have the same exclusive bin; hence at most \(s\) groups are discarded.
In each of the more than \(2s\) remaining groups, every chosen collision
bin also contains a retained item from another group. That item is inactive
at the chosen coordinate. We have therefore found more than
\(\gamma s|E|\) witnesses with two active items and one inactive item.
Since \(q_0<8s|E|\), their fraction after padding exceeds
\(\gamma/8\), and hence the stated \(\zeta=\gamma/24\).
Uniform repetition to width \(q_2\) preserves that fraction.
\end{proof}

Retained coordinate blocks give lower bounds on the conversation that remains
after a prefix. The first estimate below handles a constant retained column
mass; the second permits exponentially smaller mass when more blocks remain.
Their proofs, and the verification of the source profiles and forced prefix, are
in \cref{sec:interlacing_details}.

\begin{lemma}\label{lem:product-blocks}
Suppose $\RP{M}{x}{H}$ holds for a rationally weighted matrix $M$, where
$H\ge2$. Let $u\ge2$ be dyadic with $ux\le1$. Every full-product interlace
restriction retaining at least $u/2+1$ blocks, each of base row mass at least
$ux$, and a column set of mass at least $9/25$, has cost at least $H+\log u$. In
particular, two blocks of base row mass at least $2x$ and column mass at least
$9/25$ suffice for cost $H+1$.
\end{lemma}

\begin{lemma}\label{lem:many-product-blocks}
Suppose $\RP{M}{x}{H}$ holds for a rationally weighted matrix $M$, where
$H\ge3$. Let $u\ge16$ be dyadic with $ux\le1$. Every full-product interlace
restriction retaining at least $9u/16$ blocks, each of base row mass at least
$ux$, and a column set of mass at least $2^{-u/1024}$, has cost at least
$H+\log u$.
\end{lemma}

These estimates give the template profiles and the extra bit forced by three
distinct tail positions.

\begin{lemma}\label{src:profiles}
Use the full-product schedule of \cref{lem:full-schedule}, and suppose
$\D(R)=d\ge3$ and $\RP{R}{\rho}{d}$, where $\rho=2^{-L}$. Write $x=2^{-10L}$ for
the scheduled block threshold, and set
\[
H_1:=d+\log s,\qquad H_2:=H_1+\log q_2,\qquad C:=H_2+\log s+1,\qquad x_{\rm
cap}:=2s\rho.
\]
Then
\[
\D(A)=H_1,\quad \RP{A}{x}{H_1},\qquad \D(R_{\rm temp})=H_2,\quad \RP{R_{\rm
temp}}{x}{H_2},\qquad \RP{N_{\rm ref}}{x}{C-1}.
\]
In the interlace of $R$ on $\Gamma$, the first $q_1$ positions together with any
three distinct tail positions require at least $H_1+1$ bits whenever every
position retains base row mass at least $x_{\rm cap}$ and the word set has mass
greater than $1/2$.
\end{lemma}

The next lemma follows collisions until Alice must speak. Its proof uses
one conserved count: an intact collision counts twice, and after it splits
each of the two resulting parts counts once. Some continuation must then
carry more requests than its remaining communication budget permits.

\begin{lemma}[Following collision requests]\label{lem:collision-tracking}
Let $q\ge2$ be dyadic, let $0\le b<q$ be an integer, and let $x>0$
satisfy $2qx\le1/8$.
Suppose a protocol with at most
$H+\log q$ bits left has $q-b$ usable coordinate fibres, each containing
word mass at least $3/4$. Here a fibre consists of the word inputs having one
fixed coordinate label. Mark $w_0>b$ fibres as collision witnesses.
Assume the following lower bounds, with all masses measured in the original
word and selector laws.
\begin{enumerate}
\item[(i)] Until Alice speaks, a fibre of word mass at least $x$, against the
unchanged selector set, costs at least $H>0$.
\item[(ii)] At the first Alice query, one child retains a selector set on which,
for every dyadic $u\ge2$ with $ux\le1$, any $u/2+1$ distinct fibres
of word mass at least $ux$ cost at least $H+\log u$.
\item[(iii)] A marked fibre retaining word mass greater than $1/2$ costs at least
$H+1$ before Alice speaks.
\end{enumerate}
Then no such protocol exists.
\end{lemma}
\begin{proof}
Follow the fibres through Bob queries and stop at the first Alice query.
At depth $r$, maintain word mass at least $2qx2^{-r}$ in every tracked
fibre, called its support. An unmarked support follows a child retaining
at least half its mass. A marked support splits into two unmarked supports
if both children meet the next threshold; otherwise it follows the qualifying
child and keeps its mark. Thus every mark splits at most once. Before a split,
the discarded word mass is less than
\[
\sum_{r\ge0}2qx2^{-r-1}=2qx\le1/8,
\]
so an unsplit mark always retains word mass greater than $1/2$.

By (i), a support cannot end at a leaf, nor reach Bob depth $\log q+1$:
it would still have mass at least $x$ but only $H-1$ bits left.
The stopping nodes are prefix-free. A node at depth $r_v$ has
\emph{prefix capacity} $\kappa_v=q2^{-r_v}\ge1$, the remaining width budget;
Kraft's inequality gives $\sum_v\kappa_v\le q$.
Write $h_v$ for its support count and $z_v\le h_v$ for its unsplit marks.
Counting an ordinary support once and an unsplit mark twice is conserved:
a split replaces two by one plus one. Consequently
\[
\sum_v(h_v+z_v)=q-b+w_0>q\ge\sum_v\kappa_v.
\]
Some node therefore has $h+z>\kappa$. Its supports have distinct coordinates,
because descendants of a split cannot rejoin, and each has mass at least
$2\kappa x$.

If $\kappa\ge2$, the inequalities $h+z\le2h$ and $h+z>\kappa$ give
$h\ge\kappa/2+1$, since $\kappa$ is even. Thus (ii), with $u=\kappa$, requires
$H+\log\kappa$ bits after Alice's query, but only
$H+\log\kappa-1$ remain. For $\kappa=1$ and $h\ge2$, use (ii) with
$u=2$ and the available masses $2x$. In the remaining case $h=1$,
the inequalities force $z=1$. By (iii), the unsplit collision needs
$H+1$ bits at this stopping node, where only $H$ remain.
\end{proof}

\begin{theorem}[Source hardness after deletion]\label{src:near-full}
Use either the full-product source of \cref{def:source-matrix} or the
indexed source of \cref{det:construction}, with a NO graph and
$\D(R)=d\ge3$. Set $H_1=d+\log s$, $H_2=H_1+\log q_2$,
$C=H_2+\log s+1$, and $x_{\rm cap}=2s\rho$.
Suppose the three profiles and the three-tail bound stated in
\cref{src:profiles} hold with these parameters. Assume the scalar margins
\eqref{eq:width-margins}, that the selector law on $S$ is exactly eight-wise
independent with marginal $J$, and that the ordinary family's bin tuple has
an exactly eight-wise independent law with marginal $J_2$.
Then $\D(N)=C$, and every restriction retaining at least $1-\eta$ of each
input law has cost at least $C$. In the full-product case these hypotheses
are supplied by \cref{src:profiles,lem:full-schedule}.
\end{theorem}
\begin{proof}
We first recover one assignment from the template prefix. Its graph collisions
then satisfy the three lower bounds in \cref{lem:collision-tracking},
which rules out the proposed communication budget.

Suppose a retained restriction has a protocol of depth at most $C-1$,
and transpose it to the raw orientation. Each template bin loses at most
$\eta s/\lambda<1/1000$ of its $S$ law. The ordinary family $\Omega$
retains normalised mass at least $1-\eta/\lambda>99/100$.
Project its retained inputs onto the whole bin tuple (onto its whole index
label for an indexed source), forgetting the tail. This image has mass
at least $99/100$. Choose one retained completion per image label,
independently of the template row. Since template entries ignore the tail,
these representatives realise an interlace of $R_{\rm temp}$.

Apply \cref{cl:separation} to the inherited, unpruned protocol with
$u=s$ and $\xi=999/1000$. Its hypotheses follow from the profile of
$R_{\rm temp}$ and the scalar margins
$999/1000>2sx$ and $999/1000-2(s-1)x>998/1000$.
Every binary history of length below $\log s$ is reached by template
inputs, and every such node is Alice's. They are actual nodes of the
original protocol, so the same complete Alice prefix governs all its
retained inputs, including item rows and diagonal columns. Its final
parts have dominant template sets $A_1,\ldots,A_s\subseteq S$, each of
mass greater than $998/1000$.
The representatives were used only to establish this prefix; henceforth
we use all originally retained columns. The remaining budget is $H_2$,
since $C=H_2+\log s+1$.

We next fix the item assignment before choosing an activity coordinate. Call an
item deficient if less than half of its base row mass remains in the retained
restriction. By the row masses in \cref{def:source-matrix}, each item type has
total mass \(1/512\); hence fewer than \(4096\eta sn\) vertex items and fewer
than \(4096\eta s(s-3)\) auxiliary items are deficient. These bounds are within
the deletion allowances of \cref{src:item-witness} by \eqref{eq:width-margins}. Every nondeficient item can be assigned to
one prefix part receiving at least \(1/(2s)\) of its base mass, by averaging
over these \(s\) prefix parts. The schedule gives
\(4s^2\rho\le1\), so
\[
\frac1{2s}\ge2s\rho=x_{\rm cap}.
\]
Thus the assigned part retains at least \(x_{\rm cap}\) of the item's base row
law. This assignment is fixed before the activity coordinate is selected. By
\cref{src:item-witness}, at least \(\zeta q_2\) coordinates contain two active
items and an inactive item in one bin. Assign each such coordinate to one prefix
part witnessing its collision, and write \(U_j\subseteq[q_2]\) for the resulting
disjoint sets.

We remove the coordinates on which the assigned rows or diagonal words are too
depleted. For a coordinate \(\alpha\), call it a template-bad coordinate for
\(A_j\) if some one of the \(q_1\) blocks of its \(J\)-projection retains less
than \(x_{\rm cap}\) of the base law. Since \(x_{\rm cap}\le1/8\) by
\eqref{eq:width-margins}, such a block misses more than \(7/8\) of its \(1/q_1\)
mass, hence the \(J\)-projection misses at least \(7/(8q_1)\).
If there were $8q_1$ such coordinates, choose one missed event at each
and let $Z$ count these events on a random selector. Pairwise independence
of the coordinates gives $\mathbb EZ\ge7$ and
$\operatorname{Var}Z\le\mathbb EZ$. Every selector in $A_j$ has $Z=0$, so
\[
\Pr(A_j)\le\Pr(Z=0)\le
\frac{\operatorname{Var}Z}{(\mathbb EZ)^2}\le\frac17,
\]
contrary to $\Pr(A_j)>998/1000$. Only pairwise independence was used.
Thus \(A_j\) has
fewer than \(8q_1\) template-bad coordinates. Likewise, fewer than \(4\eta
q_2/\mu\) coordinates retain less than \(3/4\) of their diagonal word law.
The diagonal activity label is uniform, so each such fibre deletes more than
\(\mu/(4q_2)\) of the total column law, whose total deletion is at most
\(\eta\). Let \(B_j\) be the union of these two
bad-coordinate sets and put \(b_j:=|B_j|\). The schedule in
\eqref{eq:width-margins} gives
\[
\sum_{j=1}^s b_j<8sq_1+4s\eta q_2/\mu<\zeta q_2/4.
\]
Since \(\sum_j|U_j|\ge\zeta q_2\), some \(j\) satisfies \(|U_j|>4b_j\). Fix this
\(j\), write \(b:=b_j\), and define \(w_0:=|U_j\setminus B_j|\). Then
\(w_0\ge|U_j|-b>3b\).

Fix a usable collision coordinate \(\alpha\in U_j\setminus B_j\). Let
\(i_1,i_2\) be its two active items and let the inactive item be the third item
in their bin, as supplied by \cref{src:item-witness}. Choose retained template
rows for the first \(q_1\) projected blocks and retained rows of \(i_1,i_2\) and
the inactive item. By the entry rule in \cref{def:source-matrix}, against a
diagonal word \((\alpha,\gamma)\) these rows read
\[
R(u,\gamma_\ell)\quad(\ell\in[q_1]),\qquad
R(u,\gamma_{q_1+\pi_\alpha(i_1)}),\qquad
R(u,\gamma_{q_1+\pi_\alpha(i_2)}),\qquad R(u,\gamma_{q_1+s+1}).
\]
The three tail positions are distinct by the collision definition and the tail
convention in \cref{def:source-matrix}. Each selected row group has base mass at
least \(x_{\rm cap}\), because \(\alpha\notin B_j\) and the item assignment
above was fixed with that threshold. The representatives depend on the row data
and not on \(\gamma\). The assumed three-tail bound therefore requires \(H_1+1\) bits
whenever the retained diagonal-word mass is greater than \(1/2\).

Apply \cref{lem:collision-tracking} with $q=q_2$ and $H=H_1$.
Its usable fibres and marked fibres are exactly the $q_2-b$ coordinates and
$w_0>3b$ witnesses just obtained; the schedule gives $2q_2x\le1/8$.
We check its three lower bounds explicitly.

Until Alice speaks, the selector set $A_j$ is unchanged. Its projection onto
each $J$-coordinate has mass at least $\Pr(A_j)>998/1000>3/5$.
Choosing a surviving selector for every projected label therefore gives
an actual restriction of $A$, so its profile supplies (i).
At the first Alice query, a denser child retains selector mass greater than
$499/1000>9/25$. After transposition, its entries on supported coordinates
are $A(\gamma,z_\alpha)$: they form the interlace of $A$ with those word
sets as row blocks. The exactly eight-wise selector law and \cref{det:exact} supply (ii).
All representatives are chosen from the retained selector set, independently
of the word input. The three-tail bound proved above supplies (iii), before
that Alice query. Thus the tracking lemma excludes depth $C-1$.
The upper protocol in \cref{src:upper} has depth $C$, proving equality.
\end{proof}

\subsection{What remains after restricting tuple inputs}
\label{sec:retained-tuples}

A protocol can create dependencies by restricting its tuple inputs. For
example, after Bob announces that two independent coordinates agree, either
coordinate determines the other. We therefore use independence only in the
original tuple law. The next lemma says that conditioning on an event of
specified mass changes the small marginals substantially at only a bounded
number of coordinates.

For an event-mass lower bound $0<v\le1$ and an error cap $0<\tau\le1$, set
\[
\begin{aligned}
\varepsilon(v,\tau)&:=\min\{\tau,2^{-30}v^{1/256}\},\\
e(v,\tau)&:=8\ceil{2^{20}\varepsilon(v,\tau)^{-2}v^{-1/512}},
\qquad e(v):=e(v,1).
\end{aligned}
\]
The constant independence order $8192=8\cdot1024$ lets us test $1024$
disjoint events, each involving at most eight coordinates.

\begin{lemma}[Conditioning outside exceptional coordinates]\label{det:conditioning}
Let $\nu$ be the law of an indexed exactly $8192$-wise independent tuple list
on $Z^q$ with marginal $\mu$, and let $0<v,\tau\le1$. For every label event
$Y$ with $\nu(Y)\ge v$, there is a set $E\subseteq[q]$, with
$|E|\le e(v,\tau)$, such that
\[
d_{\rm TV}\bigl((\nu\mid Y)_J,\mu^J\bigr)\le\varepsilon(v,\tau)
\qquad
(J\subseteq[q]\setminus E,\ |J|\le8).
\]
Here total variation is the maximum absolute difference between the two
laws' probabilities of an event. Coordinates outside this $E$ are called
\emph{good for $Y$}. The event and exceptional set are proof objects;
neither must be computed by the reduction.
\end{lemma}

Write $\pi_JK$ for the decoded projection of a label set $K$ onto coordinates
$J$. If $K\subseteq Y$ and $J$ consists of at most eight good coordinates,
then the inverse image of $\pi_JK$ contains $K$, so
\[
\mu^J(\pi_JK)
\ge(\nu\mid Y)_J(\pi_JK)-\varepsilon(v,\tau)
\ge(\nu\mid Y)(K)-\varepsilon(v,\tau).
\]
Choose one actual retained label for each projected tuple. On the row blocks
indexed by $J$, these representatives realise the projected matrix. Its
column mass is measured in the original product law $\mu^J$; no product
assumption on $\nu\mid Y$ is needed. Before conditioning, exact marginals
give the same inequality without the error term.

The next two results turn surviving blocks into communication lower bounds.
They apply to full products as well as to the explicit lists constructed
later. Their proofs, including the small-coordinate cases and the moment
estimates behind \cref{det:conditioning}, are in
\cref{sec:conditioning-details}.

\begin{lemma}[Surviving blocks in an exact tuple list]\label{det:exact}
Suppose $d\ge2$ and $\RP{M}{x}{d}$. Let the tuple source for an interlace be
exactly eight-wise independent with the base-column marginal. If an integer
$a\ge1$ satisfies $2^ax\le1$, and a restriction retains at least
$2^{a-1}+1$ blocks, each of base row mass at least $2^ax$, together with a
tuple-label set of source mass at least $9/25$, then its complexity is at
least $d+a$.
\end{lemma}

\begin{theorem}[Surviving blocks after conditioning]\label{det:rareboost}
Suppose $d\ge2$, $\RP{M}{x}{d}$, and $\nu$ is an exactly $8192$-wise
independent tuple source with the base-column law. Let $\nu(Y)\ge v>0$,
choose the exceptional set from \cref{det:conditioning} with any cap
$0<\tau\le1$, and let an integer $a\ge1$ satisfy $2^ax\le1$. If a restriction
retains at least $2^{a-1}+1$ good blocks, each of base row mass at least
$2^ax$, and a column set $K\subseteq Y$ with $(\nu\mid Y)(K)\ge1/2$, then
its complexity is at least $d+a$.
\end{theorem}


\subsection{Where displaced inputs can go}\label{sec:internal-interlace}

We now prove the internal-interlace bound once for both constructions.
Protected inputs are the diagonal inputs in $\mathcal D$. In
$U=\Int{w}{N}$, call a copy an input's \emph{origin}. Measure transferred
mass in units of one complete copy of $N$, without division by $w$;
reference losses within $\Omega$ are normalised. Recall that $\Omega$ has
mass $\lambda=255/256$ and $\mathcal D$ has mass $\mu=1/256$.

Here are the inputs to the common argument. The source profiles are
$\RP{A}{x}{H_1}$, $\RP{R_{\rm temp}}{x}{H_2}$ and
$\RP{N_{\rm ref}}{x}{C-1}$, with $H_2=H_1+\log q_2$,
$C=H_2+\log s+1$, and $H_1,H_2\ge2$.
The source has exact cost $C$, including on every restriction retaining
$1-\eta$ of both laws. The selector law $S$ is exactly $8192$-wise
independent on $J$ coordinates; the ordinary bin tuple has that same
independence order and marginal $J_2$. The internal tuple law $\nu$ has
the raw $\mathcal K$ marginal and the following representation by three
independent complete streams: metadata selects a bin or an item, base
labels have the $X$ law, and selectors have the $S$ law. Each stream is
exactly $8192$-wise independent across origins. The decoder uses
$(p,z_i)$ at a template position and $(a,u_i)$ at an item position.
These hypotheses imply the exact $8192$-wise $\mathcal K$ marginals.
They also imply exact $8192$-wise independence of the flattened selector
coordinates $(i,\alpha)$: group at most $8192$ requests by origin, apply
outer independence, then the independence within each $S$ selector.
Fixing the other complete streams leaves this ambient selector law unchanged.

Both schedules satisfy \eqref{eq:width-margins} and the following bounds,
where $C_{128}=(2^{14}/(\lambda-1/2)^2)^{128}$ is the flag-moment
constant used in the proof of \cref{res:incoming}:
\begin{equation}\label{eq:internal-hypotheses}
\begin{gathered}
\delta_U=\eta/(16w),\qquad
 e_0\ge e(\theta,\eta/2),\qquad E_{\rm exc}=2e_0,\qquad
 K_*=512\mu(\log w+2)s,\\
 E_{\rm exc}\le w/4,\qquad K_*E_{\rm exc}\le\eta w/16,\\
 e(\theta)\le w/128,\qquad e(\theta/4)\le sq_2,\qquad
 C_{128}s^{-128}\le\theta/4.
\end{gathered}
\end{equation}
The numerical consequences of the scalar margins include
$2wx+\eta/(16\lambda)<1$, $wx<1/256$ and
$1-\eta s/(2\lambda)>2sx$; all row thresholds below are at most one.

For full products the stream representation is an auxiliary probability
space, not a change to the column alphabet $\mathcal K^w$. Independently
sample complete metadata, base-label and selector streams, even the values
that the decoder will ignore. Metadata has exactly the type/bin/item law
of $\mathcal K$. The decoder therefore pushes this probability space
forward to the original product law on $\mathcal K^w$. Pulling back a
retained event preserves its mass, and pulling back a protocol preserves
its entries and depth. There is no enumeration or size requirement on
this auxiliary space. In the explicit construction these complete streams
are the actual indexed input labels.

This distinction matters only in the receipt proof: first choose the
actual retained tuple event (or its denser Bob child), pull it back, and
then fix metadata and base streams by averaging. Only afterwards condition
on the surviving selector section. The selector law before that last step
is unchanged; no independence after conditioning is asserted. Selected
positions have template metadata, so their projected selector values
have surviving actual-column representatives. Distinct selected values
decode to distinct actual columns. Unused stream values need not be
represented in the original full-product matrix.

The full-product construction meets these inputs by
\cref{src:profiles,src:near-full,lem:full-schedule} and product independence.
The explicit construction will verify the same inputs in
\cref{sec:deterministic}.

\begin{lemma}\label{res:destinations}
Under the common inputs above, suppose a restriction of $U$ retains row
mass at least $1-\delta_U$ and a tuple event $Y'$ of ambient mass at least
$\theta$, with a protocol of depth at most $C+\log w-1$.
Apply \cref{det:conditioning} to this one event with cap $\eta/2$, removing
at most $e_0$ origins. At an Alice node of depth $r$ and capacity
$\kappa=w2^{-r}$, track each remaining reference support of normalised
mass at least $2\kappa x$ into every child retaining mass at least
$\kappa x$. An unsplit support has weight one; each descendant after its
first split has weight one-half.

Stop at the first Bob query or at depth $\log w$, whichever comes first;
nodes at that depth are called final, regardless of the next speaker.
Complete paths without supports by dummy Alice queries. A first-Bob node
of capacity $\kappa\ge2$ has weight at most $\kappa/2$, and a final node
has at most one support. At least $w-E_{\rm exc}$ final nodes carry
unsplit supports of distinct origins. Call these nodes their origins'
\emph{owners}. All other destinations together have capacity at most
$E_{\rm exc}$.
\end{lemma}
\begin{proof}
Initial deletion removes at most $w\delta_U=\eta/16$ complete copies.
Every origin starts with normalised reference mass at least
$1-\eta/(16\lambda)>2wx$. A support always has a child retaining at
least half its mass. Its first split preserves total weight, and further
splits can only increase it. Descendants of an origin never rejoin.
Thus the total tracked weight is at least $w-e_0$.

Until Bob speaks, the tuple set is the fixed $Y'$. For every good origin,
its projection has ambient $\mathcal K$ mass at least $1-\eta/2>3/5$:
the conditional law gives this image mass one, and the marginal guarantee
loses at most $\eta/2$. One surviving tuple per projected value realizes
a restriction of $N_{\rm ref}$. A support therefore cannot reach a leaf,
by its profile and row mass at least $2x$. Padding support-free paths
gives a complete frontier whose capacities sum to $w$.

At a first-Bob node of capacity $\kappa\ge2$, suppose there were at least
$\kappa/2+1$ supports. Their origins are distinct and good for the same
event $Y'$. The denser Bob child has conditional mass at least one-half.
Apply \cref{det:rareboost} to $N_{\rm ref}$ with $a=\log\kappa$:
each support has mass $2\kappa x\ge\kappa x$. The child requires
$C-1+\log\kappa$ bits, but only $C+\log\kappa-2$ remain.
Consequently there are at most $\kappa/2$ supports and at most that weight.
This also covers $\kappa=2$.

At a final node, two supports of mass $2x$ would, by the same theorem
with $a=1$ and all of $Y'$, require $C$ bits against a budget of $C-1$.
Thus a final node has at most one support. It is an owner if that support
is unsplit; otherwise its weight is at most one-half. If $c$ is the
number of owners, every other node has at most half its capacity in
weight, whence
\[
 w-e_0\le c+\tfrac12(w-c),\qquad w-c\le2e_0=E_{\rm exc}.
\]
Each owner has capacity one. An unsplit support follows only one path,
so owners have distinct origins and all other destinations have the
claimed total capacity.
\end{proof}

Two receipt bounds complete the argument. Their proofs in
\cref{sec:mass_bounds} use the common stream representation above.

\begin{lemma}\label{res:incoming}
In the protocol and frontier of \cref{res:destinations}, the protected
mass from all origins reaching a first-Bob node of capacity $\kappa$ is
less than $K_*\kappa$. At a final node it is less than $K_*$.
Mass is measured in units of one complete copy of $N$.
\end{lemma}

\begin{lemma}\label{res:foreign}
In the same protocol, let $\varepsilon_i$ be the normalised reference
mass missing at owner $i$. Every other origin that has an owner sends at
most $\mu\delta_i$ protected mass to owner $i$, where
\[
 \delta_i=2q_2x+4\varepsilon_i.
\]
\end{lemma}

\begin{lemma}[Conservation of protected mass]\label{lem:protected-accounting}
Measure mass in units of one complete source copy. Suppose $w\ge2$
origins send surviving rows to a complete Alice-only frontier, and at
least $3w/4$ origins have distinct owners. Each copy has reference mass
$\lambda=255/256$ and diagonal mass $\mu=1/256$.
The following properties cannot all hold:
\begin{enumerate}
\item[(i)] Each owner misses more than $\eta$ of its own copy.
\item[(ii)] The missing normalised reference masses satisfy
$\sum_i\varepsilon_i\le\eta/(8\lambda)$, and initial deletion removes
at most $\eta/16$ protected mass from these origins in total.
\item[(iii)] Destinations other than owners receive at most $\eta w/16$
protected mass from these origins.
\item[(iv)] Every other owner-origin sends at most
$\mu(2qx+4\varepsilon_i)$ protected mass to owner $i$, where
$4\mu wqx\le\eta/32$.
\end{enumerate}
\end{lemma}
\begin{proof}
Let $P_i$ be the surviving protected mass from origin $i$ sent outside
its owner, and $a_i$ its initially deleted protected mass. Missing
reference rows, initially deleted protected rows, and protected rows
elsewhere partition the missing part of its copy. Thus
$P_i+a_i+\lambda\varepsilon_i>\eta$, and
\[
 \sum_iP_i>\frac{3\eta w}{4}-\frac{\eta}{16}-\frac{\eta}{8}
 >\frac{\eta w}{2}.
\]
But other owners receive at most
\[
 \mu w\sum_i(2qx+4\varepsilon_i)
 \le w\left(2\mu wqx+\frac{\mu\eta}{2\lambda}\right)
 \le\eta w\left(\frac1{64}+\frac1{510}\right)<\frac{\eta w}{32}.
\]
Together with (iii), this is less than $3\eta w/32$. Every surviving
row has one destination because only Alice has spoken, contradicting
the outgoing mass.
\end{proof}

\begin{figure}[H]
\centering
\begin{tikzpicture}[
  >=Stealth,font=\footnotesize,
  box/.style={draw,rounded corners=2pt,align=center,inner sep=5pt,minimum height=12mm},
  src/.style={box,text width=28mm},
  dst/.style={box,text width=38mm},
  flow/.style={->,semithick}]
\node[src] (origin) at (0,0) {Protected inputs\\from copies with owners};
\node[src,draw=black!55] (deleted) at (0,-1.8)
  {Initially deleted\\$\sum_i a_i\le\eta/16$};
\node[dst,draw=black!55] (own) at (4.7,1.6)
  {Stay at their own owner\\not counted in $P_i$};
\node[dst,fill=aliceblue!6,draw=aliceblue] (outgoing) at (4.7,-.4)
  {Survive outside their own owner\\$\displaystyle\sum_iP_i>\frac{\eta w}{2}$};
\node[dst] (other) at (10.1,.6)
  {Go to other owners\\total $\le\eta w/32$};
\node[dst] (exceptional) at (10.1,-1.4)
  {Go to nonowners\\total $\le\eta w/16$};
\coordinate (firstsplit) at (2.1,0);
\draw[semithick] (origin.east) -- (firstsplit);
\draw[flow,black!55] (origin.south) -- (deleted.north);
\draw[flow,black!55] (firstsplit) |- (own.west);
\draw[flow] (firstsplit) |- (outgoing.west);
\coordinate (secondsplit) at (7.4,-.4);
\draw[semithick] (outgoing.east) -- (secondsplit);
\draw[flow] (secondsplit) |- (other.west);
\draw[flow] (secondsplit) |- (exceptional.west);
\node[box,fill=black!4,text width=127mm] at (4.7,-3.3)
  {\textbf{There is nowhere for all the outgoing mass to go.}\\[3pt]
   Required outgoing mass $>\eta w/2
   \;>\;3\eta w/32=\eta w/32+\eta w/16$
   (total receiving bound).};
\end{tikzpicture}
\caption{The conservation contradiction. Every surviving protected input
outside its own owner must reach one of the two destinations on the right,
whose combined receipt bound is too small. Nonowners are first-Bob or final
nonowner nodes. Deletion and retention at the own owner are accounted for
separately. All masses are in complete-copy units; arrow widths carry no scale.}
\label{fig:reference-destinations}
\end{figure}

\begin{theorem}[Lower bound for the internal interlace]\label{res:alignment}
Under the common inputs of \cref{sec:internal-interlace},
\[
 \D(W)=C+\log w,\qquad \RP{W}{\theta}{C+\log w-1}.
\]
Every restriction of $W$ retaining row mass at least $\theta$ and column
mass at least $1-\delta_U$ has cost at least $C+\log w$.
\end{theorem}
\begin{proof}
Suppose a restriction violates the near-full bound, and transpose to $U$.
Its retained tuple event $Y'$ has mass at least $\theta$.
By \cref{res:destinations}, at least $w-E_{\rm exc}\ge3w/4$ origins
have owners, each with only $C-1$ bits left. Conditioning on $Y'$ with
cap $\eta/2$ gives every owner-origin a column projection of ambient
mass at least $1-\eta/2$. Surviving representatives realize a restriction
of $N$. Since its near-full cost is $C$, the owner must miss more than
$\eta$ of its own row law.

An unsplit reference support discards less than
$\sum_{r=0}^{\log w-1}wx2^{-r}=2(w-1)x$ on its path.
Initial deletion costs at most $\eta/16$ complete copies, so
\[
 \sum_i\varepsilon_i\le\frac{\eta}{16\lambda}+2w(w-1)x
 \le\frac{\eta}{8\lambda}.
\]
This also bounds initially deleted protected mass by $\eta/16$.
Exceptional destinations have total capacity at most $E_{\rm exc}$;
\cref{res:incoming} bounds their receipt by
$K_*E_{\rm exc}\le\eta w/16$.
Finally \cref{res:foreign} and $4\mu wq_2x\le\eta/32$ supply the
last hypothesis of \cref{lem:protected-accounting}, with $q=q_2$.
Its contradiction proves the near-full bound.

For the profile, fix $j\in\{0,1\}$ and transpose a restriction of $W$
with row mass at least $\theta$ and column mass at least
$\beta_j=(3/5)2^{-j}$. In $U$ its tuple event $Y$ has mass at least
$\theta$. Discard the diagonal rows. The average normalised reference
retention is at least $\beta'_j=(\beta_j-(1-\lambda))/\lambda$.
Put $\kappa=w2^{-j}$ and apply \cref{det:conditioning} to $Y$ with
cap one. If $m$ good origins retain reference mass at least $\kappa x$,
block averaging gives $w\beta'_j\le m+e(\theta)+w\kappa x$ and hence
\[
 \frac m\kappa\ge2^j\beta'_j-wx-\frac{4e(\theta)}w
 \ge\frac{(3/5)-4(1-\lambda)}\lambda-\frac1{256}-\frac1{32}
 =\frac{44}{75}-\frac1{256}-\frac1{32}>\frac12.
\]
Thus at least $\kappa/2+1$ good blocks remain. Apply
\cref{det:rareboost} to $N_{\rm ref}$ with $a=\log\kappa$, using all
of $Y$, of conditional mass one. The schedules give $\kappa\ge2$ and
$\kappa x\le1$. The lower bound is
$C-1+\log\kappa=C+\log w-1-j$, proving the profile after transposition.
The near-full bound applies to the whole matrix. Conversely the selector
protocol sends the copy index and then runs the cost-$C$ protocol for
$N$, giving equality $\D(W)=C+\log w$.
\end{proof}

\subsection{Restoring the bound for the next round}

Repeated amplification needs a restriction profile at the exact cost, whereas
\cref{res:alignment} leaves that profile one bit lower. Its stronger bound when
almost all columns remain is the additional hypothesis used by \cref{cl:bridge}
below. The starting matrix must meet the full profile as well; the lookup seed
and the restoration theorem are proved in \cref{sec:interlacing_details}.

\begin{theorem}[Restoration of the restriction profile]\label{cl:bridge}
Let $W$ have $m\ge1$ rows and $n\ge1$ columns, and suppose its exact cost
satisfies $\D(W)\ge3$. Let $\theta>0$ and $0<\delta\le1/64$. Assume that every
restriction retaining row mass at least $\theta$ and column mass at least
$1-\delta$ has cost at least $\D(W)$, and that $\RP{W}{\theta}{\D(W)-1}$ holds.
For an integer $b\ge1$, choose any power of two $Q=2^r$ satisfying
$Q\delta\ge80b$, and assume $Q\theta<1/20$. Then $R^+=(\Int{Q}{W})^{\mathsf T}$
satisfies
\[
\D(R^+)=\D(W)+r,\qquad \RP{R^+}{2^{-b}}{\D(W)+r}.
\]
Independently of the lower-bound hypotheses, any upper bound on $\D(W)$ gives an
upper bound larger by $r$ on $\D(R^+)$. The matrix $R^+$ has exactly $n^Q$ rows
and $Qm$ columns.
\end{theorem}

\begin{lemma}\label{lem:amplifier-rank}
In the source construction, let
$h=2\log s+\log q_2+\log w+\log Q$ and suppose
$|\mathcal I|\le2sq_2$. Then $\rank(R^+)\le2^h\rank(R)$.
The same bound holds if any word alphabet is replaced by an indexed list
of legal words, including products of such lists.
\end{lemma}
\begin{proof}
A width-$t$ interlace stacks $t$ blocks obtained by selecting or repeating
columns of its base matrix, and hence multiplies rank by at most $t$.
The template part therefore has rank at most $sq_1q_2\rank(R)$.
For each item $i$, its block has entries $R(u,\psi_i(Z))$, where the selected
column $\psi_i(Z)$ is independent of $u$; its rank is at most $\rank(R)$.
After the outer interlaces, rank subadditivity gives
\[
\rank(R^+)\le wQ(sq_1q_2+|\mathcal I|)\rank(R)
\le wQsq_2(q_1+2)\rank(R)=2^h\rank(R),
\]
using $q_1+2=s$. Restricting word lists, duplicating indexed labels, and
transposing preserve this argument.
\end{proof}

\begin{proof}[Proof of \cref{thm:amplifier}]
We first apply the restoration theorem to the NO case, then derive the YES upper
bound and the rank bound, which does not use the graph promise.

For a NO graph, \cref{src:near-full} gives \(\D(N)=C\), and \cref{res:alignment}
gives
\[
\D(W)=C+\log w,\qquad \RP{W}{\theta}{C+\log w-1},
\]
together with the near-full restriction bound required by \cref{cl:bridge}. By
\cref{lem:full-schedule}, \(0<\delta_U\le1/64\), \(Q=w^2\) satisfies \(Q\delta_U\ge80L\), and
\(Q\theta=1/512<1/20\). Applying \cref{cl:bridge}
with \(b=L\), \(\delta=\delta_U\), and \(W\) as above yields
\[
\D(R^+)=C+\log w+\log Q,\qquad \RP{R^+}{2^{-L}}{C+\log w+\log Q}.
\]
The definitions in \cref{src:profiles} give \(C=d+2\log s+\log q_2+1\). Thus,
with
\[
h:=2\log s+\log q_2+\log w+\log Q,
\]
these bounds become
\[
\D(R^+)=d+h+1,\qquad \RP{R^+}{2^{-L}}{d+h+1}.
\]
On either source promise, \cref{src:upper} and the two outer selector protocols
from \cref{def:source-matrix} give the upper bound \(\D(R^+)\le\D(R)+h+1\); on a
YES source the source protocol omits the final bit by \cref{src:upper}, giving
\(\D(R^+)\le\D(R)+h\). The overhead estimate \(h=L+O_{\Delta,\gamma}(\log L)\)
and the effective admissibility cutoff are supplied by \cref{lem:full-schedule}.
Recursive input generation and entry evaluation are supplied by
\cref{lem:sample-access}.

The source-coordinate bound $q_2\ge4s|E|$ and $|E|\ge n/2$ imply
$|\mathcal I|=4s(n+s-3)\le2sq_2$.
Thus \cref{lem:amplifier-rank} supplies the rank bound, independently of the
graph promise.
\end{proof}

%% ---- section: sampling ----

\section{Interlacing on small input lists}
\label{sec:small-lists}
\subsection{Constructing exact marginals}

The deterministic reduction replaces each full product by an explicit tuple
list, as defined in \cref{def:exact-tuples}. The list agrees with the product
on every small set of coordinates; \cref{sec:retained-tuples} explains why
this remains useful after a protocol restricts its inputs.

The construction appends one coordinate at a time and removes words through
linear dependences that preserve every small marginal. This is the sequential
construction idea of \citet{koller1994sample}; its implementation and bit bound
are given in \cref{sec:support-algorithm}. We store the rational weights
directly and measure restrictions in these laws.

\begin{lemma}[Exact rational supports]\label{det:code}
Fix an integer $\mathsf D\ge1$. Given a rational law $\mu$ on an alphabet of
size $m$ and an integer $q\ge1$, a deterministic algorithm constructs an exactly
$\mathsf D$-wise independent source $\mathsf{Code}_{\mathsf D}(q,\mu)$ with
support size at most $(1+qm)^{\mathsf D}$.
For fixed $\mathsf D$, the construction time, total output length and
rational-weight bit lengths are polynomial in $q$, $m$ and the binary encoding
length of $\mu$.
\end{lemma}

Every construction call below uses $\mathsf{Code}_{8192}$.

\subsection{Restoring the full restriction profile}

Restoration uses the base matrix's full cost on near-full column sets when
Bob has not yet spoken. Once he has spoken, the conditioned block bound
supplies the extra bit from the weaker profile.

\begin{theorem}\label{det:restore}
Let $0<\phi\le1$, $0<\delta\le1/64$ and $\theta>0$. Suppose $W$, with input laws
$\mu_X,\mu_Y$, satisfies $\D(W)=C\ge3$, $\RP{W}{\theta}{C-1}$, and
\[
\D(W[A,B])\ge C \quad\text{whenever }\mu_X(A)\ge\theta
\text{ and }\mu_Y(B)\ge1-\delta.
\]
Choose a power of two $Q$ satisfying
\[
Q\ge128\bigl(e(\phi,\delta/2)+1\bigr), \qquad Q\theta\le1/20.
\]
Interlace $W$ $Q$ times on any exactly $8192$-wise independent tuple source with
the base-column law, and transpose the result to obtain $R^+$. Then
\[
\D(R^+)=C+\log Q, \qquad \RP{R^+}{\phi}{C+\log Q}.
\]
Independently of the lower-bound hypotheses, a protocol of depth $U$ for $W$
yields one of depth $U+\log Q$ for $R^+$.
\end{theorem}
\begin{proof}
Fix $j\in\{0,1\}$ and a restriction of $R^+$ retaining row mass at least
$\phi$ and column mass at least $(3/5)2^{-j}$. Undo the transpose. The
columns now form an event $Y$ of source mass at least $\phi$, and the rows
have average block mass at least $(3/5)2^{-j}$. Apply
\cref{det:conditioning} to $Y$ with cap $\delta/2$, and put
\[
e=e(\phi,\delta/2),\qquad k=Q2^{-j},\qquad a=\log k.
\]
The width assumptions make $a$ a positive integer. Call a block qualifying
when its coordinate is good and its retained base-row mass is at least
$k\theta$. There are more than $k/2$ qualifying blocks: otherwise the
block-averaging observation gives
$(3/5)Q2^{-j}\le k/2+e+Qk\theta$, and hence
\[
\frac35\le\frac12+\frac{2^je}{Q}+Q\theta
\le\frac12+\frac1{64}+\frac1{20}<\frac35.
\]
Retain the qualifying blocks and all of $Y$.

Suppose this restriction has a protocol of depth at most $C+a-1$.
Follow the mass-halving choices from \cref{det:rareboost} for $a$ messages,
starting with conditional column mass one. Pad an early leaf with dummy
Alice messages if necessary. If $\alpha$ messages are Alice's and $\beta$
are Bob's, then $\alpha+\beta=a$ and the surviving restriction has
\[
\begin{gathered}
\text{more than }2^{\beta-1}\text{ good blocks},\qquad
\text{row mass per block at least }2^\beta\theta,\\
(\nu\mid Y)(K')\ge2^{-\beta},\qquad
\text{at most }C-1\text{ bits remaining}.
\end{gathered}
\]
If $\beta=0$, the columns are still all of $Y$. At least one good block
remains, with row mass at least $\theta$. Its projected column set has base
mass at least $1-\varepsilon(\phi,\delta/2)\ge1-\delta$, by
\cref{det:conditioning}. Choosing actual retained representatives gives a
restriction of $W$ that requires $C$ bits by the near-full hypothesis.

If $\beta\ge1$, apply the short-continuation bound \cref{det:suffix} with
$d=C-1$, $b=\beta-1$, and $x=\theta$. The integer block count is at least
$2^b+1$, and the row and column masses meet its thresholds. Also $d\ge2$ and
$2^{b+1}x=2^\beta\theta\le k\theta\le1$. Again $C$ bits are required.
Both cases contradict the remaining budget.

Thus each restriction requires at least $C+a=C+\log Q-j$ bits. Transposing
back proves the profile and the lower bound on the whole matrix. For the
upper bound, send the block index in $\log Q$ bits, decode that coordinate,
and run the depth-$U$ protocol for $W$. Exchanging the players gives the
same cost for $R^+$. Taking $U=C$ proves equality.
\end{proof}

\section{Deterministic gap tables}\label{sec:deterministic}
\label{sec:deterministic_amplification}

For every requested integer \(k\ge1\), the explicit-list construction below
preserves a deterministic gap of \(k\) bits after \(k\) rounds, with the YES and
NO costs and the complete encoding bound given in \cref{det:reduction}.

We use the subscript $\rm g$ for the explicit-list construction. For the
source graph $G$, choose an admissible dyadic $\mathsf N\ge |V(G)|$ as in
\cref{lem:generated-schedule}, and put $L_{\rm g}=\log\mathsf N$ and
$\rho_{\rm g}=2^{-L_{\rm g}}$. That lemma supplies all subscripted widths and
thresholds below. They depend on $\mathsf N$ and the fixed graph promise,
and are kept unchanged as the base matrix $R$ passes through successive rounds.

The one-round statement records what the construction must preserve. Its
proof below verifies the source profiles and the common internal-interlace
hypotheses, then restores the profile for another round.

\begin{theorem}\label{det:amplifier}
Under the admissible schedule of \cref{lem:generated-schedule}, the
construction of \cref{det:construction} is
deterministic and polynomial in \(\mathsf N\) and the explicit rational encoding
length of \(R\). Set
\[
h_{\rm g}=2\log s_{\rm g}+\log q_{2,\rm g}+\log w_{\rm g}+\log Q_{\rm g}
=12L_{\rm g}+2\ceil{L_{\rm g}/4}.
\]
Then every upper bound \(\D(R)\le y\) gives
\[
\D(R_{\rm g}^{+})\le y+h_{\rm g}+1,
\]
and, on a YES graph, \(\D(R_{\rm g}^{+})\le y+h_{\rm g}\). If the graph is NO,
\(\D(R)=d\ge3\), and \(\RP{R}{\rho_{\rm g}}d\), then
\[
\D(R_{\rm g}^{+})=d+h_{\rm g}+1,\qquad \RP{R_{\rm g}^{+}}{\rho_{\rm g}}{d+h_{\rm
g}+1}.
\]
On either promise, \(\rank(R_{\rm g}^{+})\le2^{h_{\rm g}}\rank(R)\). Every
generated input decodes to an input of the corresponding full-product
construction with the same entries.
\end{theorem}

\subsection{Replacing products by explicit lists}

In \cref{det:construction}, we replace each product list by an explicit rational
support while preserving the selector laws needed for another round. We retain
labels even when decoded inputs repeat, and keep the lists for different roles
independent before exposing any restriction.

\begin{lemma}\label{det:flat}
Let an indexed list on \(Z^{q_{\rm in}}\) be exactly \(8192\)-wise independent
with marginal \(\mu_Z\). Choose its indices at \(q_{\rm out}\) positions using
an exactly \(8192\)-wise independent list whose marginal is the inner index law.
Then the decoded coordinates, indexed by outer--inner pairs, are exactly
\(8192\)-wise independent with marginal \(\mu_Z\). If this complete outer list
is independent of finitely many other complete lists, fixing any
positive-probability joint value of those lists leaves its decoded ambient law
unchanged.
\end{lemma}

\begin{proof}
Write $\nu$ for the inner index law and $z_r(\ell)$ for the inner decoder.
For at most $8192$ prescribed outer--inner pairs, group the requests by outer
position: let $P$ be the positions used and $T_p$ the inner coordinates requested
at position $p$. If $I_p$ is the chosen inner index, then
\[
\Pr[z_r(I_p)=a_{p,r}\text{ for all requested pairs}]
=\prod_{p\in P}\Pr_{\ell\sim\nu}[z_r(\ell)=a_{p,r}\ (r\in T_p)]
=\prod_{p\in P}\prod_{r\in T_p}\mu_Z(a_{p,r}).
\]
The first equality uses outer independence and the second inner independence;
each involves at most $8192$ coordinates. Independence of complete lists also
means that conditioning on values of other lists does not change this law.
\end{proof}

The construction uses independent complete lists of metadata, base labels and
selector indices. In the receipt argument, fix the metadata and base-label
streams first; the ambient selector law is unchanged by \cref{det:flat}.
Then condition on the surviving selector section, using the
exceptional-coordinate bounds.

\begin{definition}\label{det:construction}
Let \(R:X\times Y\to\{0,1\}\) be a rationally weighted base matrix and use the
admissible schedule of \cref{lem:generated-schedule} and the item-activity
construction of \cref{def:source-matrix}.
Replace the word alphabets by the indexed lists
\[
\Gamma_{\rm g}=\mathsf{Code}_{8192}(2s_{\rm g}-1,\mu_Y),\qquad J_{\rm
g}=[q_{1,\rm g}]\times X,\qquad S_{\rm g}=\mathsf{Code}_{8192}(q_{2,\rm
g},\mu_{J_{\rm g}}),
\]
\[
J_{2,\rm g}=[q_{2,\rm g}]\times\Gamma_{\rm g},\qquad \mathcal J_{\rm
g}=\mathsf{Code}_{8192}(s_{\rm g},\mu_{J_{2,\rm g}}).
\]
Integer factors are uniform, and all other factors carry their base or indexed
laws. On decoded labels use
\[
A_{\rm g}(\gamma,(r,u))=R(u,\gamma_r),\qquad R_{\rm temp,g}(z,(\alpha,\gamma))=A_{\rm g}(\gamma,z_\alpha).
\]
Keep the ordinary and diagonal column families separately tagged.
The ordinary family $\Omega_{\rm g}$ has mass $\lambda=255/256$ and pairs
an independent $\mathcal J_{\rm g}$ index with the tail pushforward of the
$\Gamma_{\rm g}$ law. The diagonal family $\mathcal D_{\rm g}$ has mass
$\mu=1/256$ and is the pushforward of $J_{2,\rm g}$ under
\[
(\alpha,\gamma)\longmapsto
(\operatorname{tail}(\gamma),(\alpha,\gamma),\ldots,(\alpha,\gamma)).
\]
The raw-row alphabet is
\[
\mathcal K_{\rm g}=([s_{\rm g}]\times S_{\rm g})\sqcup(\mathcal I_{\rm g}\times
X),
\]
with type masses \(\lambda=255/256\), \(1/512\), and \(1/512\), and the same
conditional item laws as in \cref{def:source-matrix}. These decoded rules define
\(B_{\rm raw,g}\), \(N_{\rm g}=B_{\rm raw,g}^{\mathsf T}\), and the normalised
restriction to the \(\Omega_{\rm g}\) inputs, denoted
\(N_{\rm ref,g}\). The internal interlace uses three independent complete
streams of length \(w_{\rm g}\), each exactly \(8192\)-wise independent:
metadata has the type/bin/item pushforward of the raw-row law, base labels have
law \(\mu_X\), and selector indices have the \(S_{\rm g}\) law. At each
position, the metadata selects either its bin and selector or its item and base
label. Their coordinatewise combination gives an exactly \(8192\)-wise
independent \(\mathcal K_{\rm g}\)-tuple list with the raw-row marginal. This
list defines \(U_{\rm g}\) from \(N_{\rm g}\); set \(W_{\rm g}=U_{\rm
g}^{\mathsf T}\). Interlace \(W_{\rm g}\) at width \(Q_{\rm g}\) using
\(\mathsf{Code}_{8192}\) with its column law, then transpose to obtain \(R_{\rm
g}^{+}\). Retain the complete-list products, rational weights, and index labels
explicitly, including repeated decoded values.
\end{definition}

Fixing the metadata and base-label streams leaves the ambient selector stream
unchanged by \cref{det:flat}. Selecting a surviving section is a further
conditioning event, governed by \cref{det:conditioning}. For projected
restrictions, surviving representatives certify the entries; their masses
are measured in the original projected law.

\subsection{Verifying the source profiles}

We verify the source profiles separately because their retained events can
be much rarer than those in the common internal proof. Once these profiles
are available, the collision and destination arguments already proved apply
to the indexed construction without change.

\begin{lemma}\label{det:profiles}
Suppose \(\D(R)=d\ge3\) and \(\RP{R}{\rho_{\rm g}}d\). Set
\[
H_{1,\rm g}=d+\log s_{\rm g},\qquad H_{2,\rm g}=H_{1,\rm g}+\log q_{2,\rm
g},\qquad C_{\rm g}=H_{2,\rm g}+\log s_{\rm g}+1,\qquad x_{\rm cap,g}=2s_{\rm
g}\rho_{\rm g}.
\]
Then
\[
\D(A_{\rm g})=H_{1,\rm g},\qquad \RP{A_{\rm g}}{x_{\rm g}}{H_{1,\rm g}},
\]
\[
\D(R_{\rm temp,g})=H_{2,\rm g},\qquad \RP{R_{\rm temp,g}}{x_{\rm g}}{H_{2,\rm
g}},\qquad \RP{N_{\rm ref,g}}{x_{\rm g}}{C_{\rm g}-1}.
\]
In the interlace of \(R\) decoded by \(\Gamma_{\rm g}\), retaining the first
\(q_{1,\rm g}\) positions and any three distinct tail positions, each with base
row mass at least \(x_{\rm cap,g}\) against word mass greater than \(1/2\),
requires \(H_{1,\rm g}+1\) bits.
\end{lemma}

\begin{proof}
Suppress the subscript $\rm g$ and put $\lambda=255/256$.
The block-averaging observation proves all three profiles.
For $j\in\{0,1\}$ put $\beta_j=(3/5)2^{-j}$.
Transpose each profile restriction to view it as an interlace. The following
table lists its base, the target width $u$, and a lower bound for the number
$m$ of blocks whose retained base-row mass is at least $u$ times the base
profile threshold.
\[
\begin{array}{c|c|c|c}
\text{profile to prove}&\text{base and threshold}&u&m/u\text{ is at least}\\ \hline
A&R,\rho&s2^{-j}&\frac35(1-2/s)-q_1\rho\\
R_{\rm temp}&A,x&q_2 2^{-j}&\frac35-q_2x\\
N_{\rm ref}&R_{\rm temp},x&s2^{-j}&
\frac{2^j(\beta_j-(1-\lambda))}{\lambda}-sx
\end{array}
\]
For the first row there are $q_1=s-2$ actual blocks. For the last row we
first discard the non-template rows, of total mass $1-\lambda$.
Each displayed bound follows from $q\beta\le m+qt$, at threshold
$t=u\rho$ or $ux$.
The schedule makes all three bounds greater than $9/16$.

In each row, the retained tuple-label event has mass at least $x$.
For $N_{\rm ref}$, take its image on the $\mathcal J$ label, forgetting the
independent tail; this image has mass at least $x$. Entries of template rows
ignore the tail, so one surviving reference label per image label realises
precisely this projected restriction. By \cref{det:conditioning}, at most
\[
e(x)\le s/128\le u/32
\]
coordinates are exceptional. More than $17u/32>u/2$ qualifying blocks remain,
and their integer count is at least $u/2+1$. Apply \cref{det:rareboost} to all
of the conditioned event, whose conditional mass is one. Successively using
the preceding row's profile gives lower bounds
$H_1-j$, $H_2-j$, and $C-1-j$. This proves all three profiles; the schedule
ensures $u\ge s/4\ge16$ and the required row thresholds are at most one.

The $j=0$ lower bounds, together with the protocols sending the $q_1$ or
$q_2$ selector and then solving the base problem, give
$\D(A)=d+\ceil{\log q_1}=H_1$ and $\D(R_{\rm temp})=H_2$.
Finally, a collision uses $q_1+3=s+1=(2s)/2+1$ distinct positions of
$\Gamma$. Its ambient law is at least eight-wise independent, each retained
row mass is $2s\rho$, and its word mass exceeds $1/2>9/25$.
Applying \cref{det:exact} with $a=\log(2s)$ gives $d+\log(2s)=H_1+1$ bits.
\end{proof}

\subsection{Iteration, encoding size and deterministic ETH}

Each round must restore the profile as well as enlarge the gap. The one-round
guarantee \cref{det:amplifier} does this; to iterate it as a uniform reduction,
we also bound the complete rational encoding of every intermediate round.

\begin{proof}[Proof of \cref{det:amplifier}]
We bound the construction cost, apply restoration to the NO lower bound, and
then verify the upper protocols, rank, and decoding. By
\cref{lem:generated-schedule}, every width in \cref{det:construction} is
polynomial in \(\mathsf N\). By \cref{det:code}, each exact-support call of
fixed order \(8192\) takes polynomial time in its width and input-law encoding
length and produces a polynomial-length rational encoding. The construction in
\cref{det:construction} uses a fixed number of such nested calls and products of
complete lists. Composing their polynomial bounds therefore bounds the entire
round, including its indexed labels and rational weights, by a fixed polynomial
in \(\mathsf N\) and the rational encoding length of \(R\).

For a NO graph, \cref{det:profiles} supplies all source profiles and
the three-tail bound required by \cref{src:near-full}. The laws of $S_{\rm g}$
and $\mathcal J_{\rm g}$ are exactly $8192$-wise independent, and the
scalar margins hold by \cref{lem:generated-schedule}. Thus the common source
proof gives $\D(N_{\rm g})=C_{\rm g}$, even after deleting at most
$\eta_{\rm g}$ of each law.

We verify the remaining inputs of \cref{sec:internal-interlace}.
The internal tuple list in \cref{det:construction} is formed from independent
complete metadata, base-label and selector streams with the correct raw-row
marginal. The flattened selector coordinates have exact $8192$-wise
independence by \cref{det:flat}; fixing the other streams leaves their
ambient law unchanged. The ordinary bin law is $\mathcal J_{\rm g}$, and
the diagonal image in every bin is exactly $J_{2,\rm g}$. Finally
\cref{lem:generated-schedule} supplies \eqref{eq:internal-hypotheses} with
$e_{0,\rm g}\ge e(\theta_{\rm g},\eta_{\rm g}/2)$ and
$E_{\rm exc,g}=2e_{0,\rm g}$.
Consequently \cref{res:alignment} gives, for
$D_0=C_{\rm g}+\log w_{\rm g}$,
\[
 \D(W_{\rm g})=D_0,\qquad \RP{W_{\rm g}}{\theta_{\rm g}}{D_0-1},
\]
and cost at least $D_0$ when row mass is at least $\theta_{\rm g}$ and
column mass at least $1-\delta_{U,\rm g}$. Since $D_0\ge3$, these are
precisely the matrix hypotheses of \cref{det:restore}.
Apply that theorem with \(\phi=\rho_{\rm g}\), \(\theta=\theta_{\rm g}\),
\(\delta=\delta_{U,\rm g}\), and width \(Q_{\rm g}\). The numerical hypotheses
follow from \cref{lem:generated-schedule}:
\[
Q_{\rm g}\ge128\bigl(e(\rho_{\rm g},\delta_{U,\rm g}/2)+1\bigr),\qquad Q_{\rm
g}\theta_{\rm g}=1/512<1/20,\qquad 0<\delta_{U,\rm g}\le1/64.
\]
The outer source has the required exact marginals by
\cref{det:construction,det:code}. Thus \cref{det:restore} gives exact cost and
profile value
\[
D_0+\log Q_{\rm g} =d+2\log s_{\rm g}+\log q_{2,\rm g}+\log w_{\rm g}+\log
Q_{\rm g}+1 =d+h_{\rm g}+1,
\]
with profile threshold \(\rho_{\rm g}\), as required.

Every indexed input decodes locally to a legal input of the corresponding
full-product construction, with the same entry. Thus the source protocol of
\cref{src:upper}, followed by the two outer selector protocols, gives cost
$y+h_{\rm g}+1$, or $y+h_{\rm g}$ on YES graphs. The source protocol has
$s_{\rm g}$ local selector choices in the YES case and $2s_{\rm g}-1$
in general; no feature of the ambient product law is used.

The source-coordinate bound implies $|\mathcal I_{\rm g}|\le2s_{\rm g}q_{2,\rm g}$, so \cref{lem:amplifier-rank} also applies to these indexed lists.
It gives $\rank(R_{\rm g}^+)\le2^{h_{\rm g}}\rank(R)$.
If decoded labels repeat, retaining them separately or identifying them
preserves communication cost by \cref{obs:local-decoding}; weights of
identified labels are added. The construction retains the labels and their
rational weights explicitly for use in the next round.
\end{proof}

Iteration uses one public schedule at every round because restoration returns
the same profile hypothesis. The following theorem makes the round loop and its
complete encoding bound explicit.

\begin{theorem}\label{det:reduction}
There are effective absolute constants \(\mathsf A\ge2\) and \(C_0>0\) and a
deterministic algorithm with the following guarantee. Given a source graph \(G\)
and an integer \(k\ge1\), it produces an explicit total Boolean matrix
\(F_{G,k}\) and a public integer \(B_k^{\rm g}\) such that
\[
\text{YES: }\D(F_{G,k})\le B_k^{\rm g},\qquad \text{NO: }\D(F_{G,k})=B_k^{\rm
g}+k.
\]
For every admissible dyadic \(\mathsf N\ge|V(G)|\), its running time and full
binary output length are at most \(\mathsf N^{C_0\mathsf A^k}\). The constants
depend only on the fixed source reduction and construction. The output may be
square with power-of-two side length at least \(\mathsf N\). Its threshold is
\[
B_k^{\rm g}=d_{0,\rm g}+kh_{\rm g},
\]
where
\[
d_{0,\rm g}=\log\operatorname{pow}_2(64(L_{\rm g}+1))+1,\qquad L_{\rm
g}=\log\mathsf N.
\]
\end{theorem}

\begin{proof}
Choose the schedule of \cref{lem:generated-schedule} and the lookup seed of
\cref{cl:seed} with $L_{\rm g}=\log\mathsf N$. Apply
\cref{det:construction} $k$ times, storing each indexed table and its rational
input laws. The same schedule is used throughout. Finally, forget the weights
and duplicate inputs to a square dyadic side at least $\mathsf N$.

The seed has exact cost $d_{0,\rm g}\ge3$ and profile threshold
$\rho_{\rm g}=2^{-L_{\rm g}}$. At each round, \cref{det:amplifier} increases
the YES upper bound by $h_{\rm g}$ and the NO exact cost by $h_{\rm g}+1$,
restoring the same profile threshold. Induction therefore gives
\[
\D(R_i)\le d_{0,\rm g}+ih_{\rm g}\quad\text{on YES},\qquad
\D(R_i)=d_{0,\rm g}+i(h_{\rm g}+1)\quad\text{on NO}.
\]
Forgetting weights and the local duplication of \cref{cor:dyadic-padding}
change neither cost. This proves the gap with
$B_k^{\rm g}=d_{0,\rm g}+kh_{\rm g}$.

For the uniform time bound, let $S_i$ include all indexed labels, table entries,
and rational laws after round $i$. The fixed polynomial construction in
\cref{det:amplifier} gives constants $a,c$ independent of $i$ and $k$ such that
\[
S_{i+1},\ \text{time for round }i+1\le c(\mathsf N+S_i+2)^a.
\]
Choose a fixed power of two $\mathsf A\ge2$ large enough to absorb the
coefficient and the additional term $\mathsf N+2$. With
$Z_i=\mathsf N+S_i+2$, we then have $Z_{i+1}\le Z_i^{\mathsf A}$.
The polynomial seed encoding gives $Z_0\le\mathsf N^{C_1}$ for a fixed
constant $C_1$, and hence
\[
Z_i\le\mathsf N^{C_1\mathsf A^i},\qquad
\text{total round time}\le k\mathsf N^{C_1\mathsf A^k}
\le\mathsf N^{(C_1+1)\mathsf A^k}.
\]
The final square side is less than $2\max\{\mathsf N,m_k,n_k\}\le2Z_k$,
where $m_k,n_k$ are the table dimensions. Printing its entries costs a fixed
polynomial in $Z_k$. Enlarging a single effective constant $C_0$ therefore
covers the schedule, seed, all rounds, threshold computation, and final output
within $\mathsf N^{C_0\mathsf A^k}$. This bound uses one fixed polynomial
per round, not a polynomial whose degree can depend on $k$.
\end{proof}

\begin{proof}[Proof of \cref{intro:deterministic}]
For fixed $k$, apply the graph reduction of \cref{src:gap-colouring} and then
\cref{det:reduction}. The first admissible dyadic size obeys
$\mathsf N\le2\max\{|V(G)|,\mathsf N_0\}$ for an effective constant
$\mathsf N_0$, by \cref{lem:generated-schedule}. Since $k$ is fixed,
$\mathsf N^{C_0\mathsf A^k}$ is polynomial in the formula length.
The resulting public threshold separates costs at most $B_k^{\rm g}$ from
cost exactly $B_k^{\rm g}+k$, on a square table with power-of-two side.
\end{proof}

\begin{proof}[Proof of \cref{det:eth}]
We choose enough rounds to separate the estimator's error intervals while
keeping both construction and estimation subexponential in the sparse SAT size.
Suppose the deterministic estimator excluded by \cref{det:eth} exists for the
coefficient \(c_*\) chosen below. Given a sparse formula on \(N\) variables,
apply \cref{src:gap-colouring} and write \(n=|V(G)|\). Its size bound and
\cref{lem:generated-schedule} give an admissible dyadic \(\mathsf
N\ge\max\{N,n\}\) with \(\mathsf N=N^{O(1)}\): start the schedule search at the
least dyadic size at least \(\max\{N,n\}\), and use its effective cutoff.
Enlarge the effective constant \(\mathsf A\) of \cref{det:reduction} to a power
of two if necessary. For sufficiently large \(N\), choose
\[
k=\floor{\frac{\floor{\log N}}{2\log\mathsf A}},
\]
Then \(k\ge1\), \(k=\Theta(\log N)\), and \(k\log\mathsf A\le(\log N)/2\), so
\(\mathsf A^k\le\sqrt N\). Apply \cref{det:reduction} and use its square output,
with \(T\) entries and threshold \(B_k^{\rm g}\). Its encoding bound and minimum
side give
\[
N^2\le\mathsf N^2\le T\le\mathsf N^{C_0\mathsf A^k},\qquad \log T\le C_0\mathsf
A^k\log\mathsf N=O(\mathsf A^k\log N).
\]
Taking another logarithm yields
\[
\log\log T\le k\log\mathsf A+O(\log\log N).
\]

Choose an effective positive \(c_*\) with \(2c_*\log\mathsf A<1\). The preceding
bound and \(k=\Theta(\log N)\) imply
\[
\frac{2c_*\log\log T}{k} \le2c_*\log\mathsf A+O\!\left(\frac{\log\log N}{\log
N}\right)<1
\]
for all sufficiently large \(N\). Write \(a_T=c_*\log\log T\) and let \(\widehat
D\) be the estimator's output. The promises in \cref{det:reduction} and the
estimator guarantee then give
\[
\begin{array}{ll}
\text{YES:}&\widehat D\le B_k^{\rm g}+a_T<B_k^{\rm g}+k/2,\\
\text{NO:}&\widehat D\ge B_k^{\rm g}+k-a_T>B_k^{\rm g}+k/2.
\end{array}
\]
Thus comparison with the public binary rational \(B_k^{\rm g}+k/2\), permitted
by \cref{def:table-estimation}, decides satisfiability using
\cref{src:gap-colouring}.

By \cref{det:reduction}, construction takes at most
\[
\mathsf N^{C_0\mathsf A^k} =2^{O(\sqrt N\log N)}=2^{o(N)}
\]
time. Any fixed polynomial in \(T\), including the estimator's running time and
the comparison, has the same subexponential bound by the displayed upper bound
on \(T\). The source reduction is polynomial-time by \cref{src:gap-colouring}.
Finally, \(T\ge N^2\) ensures that the estimator's guarantee on every
sufficiently large table applies for all sufficiently large \(N\). Exhaustive
search handles the fixed finite remainder, including the values for which
\(k<1\). We have obtained a deterministic \(2^{o(N)}\)-time algorithm for sparse
SAT, contradicting \cref{def:reth}.
\end{proof}

%% ---- section: alternation ----

\section{Strictly alternating protocols}\label{sec:alternation}
\label{sec:alternating}

We now apply the sparsification principle to a lookup construction designed
for strict alternation. The unrestricted construction's short protocols use
multi-bit messages. Simulating them by inserting the missing turns can double
the cost and consume the additive gap, so we need a protocol and lower bound
adapted to the alternating model.

Start conceptually with a table in which Bob may hold any Boolean array on
the lookup positions. A balanced hypergraph colouring gives a short alternating
protocol. In a NO instance, the absence of a suitable colouring forces many
lookup requests to share Alice's early response strategy, leaving too few
remaining response patterns. Uniform arrays give the pattern-probability
bound needed for this contradiction, but give an exponential column
family. Replacing that family by a small explicit list preserves the pattern
bound used in the proof. This yields a polynomial-size table for every fixed
gap, using the fixed turn order in place of the repeated interlacing amplifier.
The polynomial's degree may grow with the requested gap.

\subsection{Starter, output and lookup inputs}

We first compare the choices of starter and output; these constant losses are
used when transferring the gap to other conventions.

\begin{lemma}\label{yao:comparison}
On nonempty input sets, each cost
$C\in\{D_A^{\rm alt},D_B^{\rm alt},D_{\rm free}^{\rm alt},Y_{\rm free}\}$
satisfies
\[
|C(M)-D_A^{\rm alt}(M)|\le2.
\]
\end{lemma}
\begin{proof}
Prepending one fixed bit changes the starter, so
$|D_A^{\rm alt}-D_B^{\rm alt}|\le1$ and
\[
D_A^{\mathrm{alt}}-1\le D_{\mathrm{free}}^{\mathrm{alt}}\le D_A^{\mathrm{alt}}.
\]
Public output also gives receiver output. Conversely, at a nonzero-depth
terminal transcript the receiver is the next scheduled speaker and can announce
the answer in one bit. Public stopping tells both players when to make this
announcement. At depth zero, choose the designated output player as the starter
and let that player announce the answer. Thus
\[
Y_{\mathrm{free}}\le D_{\mathrm{free}}^{\mathrm{alt}}\le Y_{\mathrm{free}}+1,
\]
and the claimed two-bit bound follows. These transformations preserve public
stopping and use an input-independent starter.
\end{proof}

Next we identify the source promise used by the lookup construction. It supplies
a small colouring for the intended protocol and forbids the larger colouring
that a short competing protocol would induce. We use the balanced hypergraph
conversion of Guruswami and Lee~\cite{guruswami2018balanced}, including distinct
vertices within each edge, recalled in \cref{sec:sources}.

\begin{proposition}[Balanced hypergraph source]\label{yao:source}
Fix an integer $t\ge4$, and set
\[
s_{\rm a}=2^t,\qquad K_{\rm a}=s_{\rm a}^2,\qquad S_{\rm a}=2^{s_{\rm a}}.
\]
There is a deterministic polynomial-time reduction from $3$-SAT to
$K_{\rm a}$-uniform hypergraphs with nonempty vertex sets and distinct vertices
in every edge such that:
\begin{enumerate}
\item a satisfiable formula has an $s_{\rm a}$-colouring with at least $s_{\rm a}-1$ vertices of every colour in every edge;
\item an unsatisfiable formula has no proper $S_{\rm a}$-colouring.
\end{enumerate}
The polynomial may depend on $t$.
\end{proposition}

The table turns this source promise into a lookup problem. Alice's vertex and
one stored address select a bit in Bob's array. Bob specifies which address to
use and an ordered edge. The address supplies the first array coordinate, and
the vertex's position in the edge supplies the second. If the vertex is outside
the selected edge, or the edge index is unused, the answer is zero.

\begin{definition}\label{def:alternating-table}
Let $\mathcal H=(V,E)$ be a source hypergraph with an ordering of the vertices
in each edge. Set
\[
q_{\rm a}=2^{t-1}=s_{\rm a}/2, \qquad H_{\rm a}=2^b,
\]
where $H_{\rm a}$ is the least power of two at least $\max\{1,|E|\}$. Index the
edges by distinct $\alpha\in[H_{\rm a}]$, and let $\iota_\alpha(v)$ be the
position of $v$ in that edge, or zero if $v$ is absent or the index is unused.
Use every address tuple once in
\[
\mathcal C=[H_{\rm a}]^{q_{\rm a}}.
\]
For any labelled Boolean array list $\mathcal W=(w^\ell)$ on
\[
\Omega_{\rm a}=[H_{\rm a}]\times[K_{\rm a}],
\]
define
\[
M_{\rm a}((v,c),(g,\alpha,\ell)):=
\begin{cases}
w^\ell(c_g,\iota_\alpha(v)),&\iota_\alpha(v)>0,\\
0,&\iota_\alpha(v)=0.
\end{cases}
\]
Rows are indexed by a vertex and an address tuple. Column labels lie in
$[q_{\rm a}]\times[H_{\rm a}]\times[|\mathcal W|]$. Equal arrays retain their
separate labels.
\end{definition}

The two selections in \cref{fig:yao-entry} determine the row and column of
Bob's array separately: \(g\) selects an address from Alice's tuple, and her
vertex determines its position in Bob's edge. Their intersection is the
requested bit. A vertex outside the edge needs no array lookup.

\begin{figure}[H]
\centering
\begin{tikzpicture}[>=Stealth,every node/.style={font=\small},
 data/.style={draw,rounded corners=2pt,align=left,inner sep=5pt}]
\matrix (bits) at (8,0) [matrix of math nodes,
 nodes={minimum width=8mm,minimum height=6mm,inner sep=1pt},
 column sep=0pt,row sep=0pt,
 row 1/.style={nodes={font=\footnotesize}},
 column 1/.style={nodes={font=\footnotesize}}] {
  &1&2&3&4\\
 1&1&0&1&0\\
 2&0&1&0&0\\
 3&0&1&1&1\\
 4&1&0&0&1\\
};
\begin{scope}[on background layer]
  \node[fill=aliceblue!8,fit=(bits-2-2)(bits-2-5),inner sep=0pt] {};
  \node[fill=boborange!8,fit=(bits-2-4)(bits-5-4),inner sep=0pt] {};
\end{scope}
\node[draw=aliceblue,fit=(bits-2-1),inner sep=0pt] {};
\node[draw=boborange,fit=(bits-1-4),inner sep=0pt] {};
\node[draw=answergreen,very thick,fit=(bits-2-4),inner sep=0pt] {};
\node[data,draw=aliceblue,fill=aliceblue!5,text width=48mm,anchor=east]
 (address) at ([xshift=-17mm]bits-2-1.west)
 {Alice's tuple: $c=(3,\boxed{1})$\\
  Bob selects $g=2$\\[2pt]
  Address: $c_g=1$};
\node[data,draw=boborange,fill=boborange!5,text width=54mm,anchor=south]
 (position) at ([yshift=15mm]bits-1-4.north)
 {Bob's edge: $e_\alpha=(u,v,\boxed{z},r)$\\
  Alice's vertex: $z$\\[2pt]
  Position: $\iota_\alpha(z)=3$};
\draw[->,aliceblue,semithick] (address.east) -- (bits-2-1.west);
\draw[->,boborange,semithick] (position.south) -- (bits-1-4.north);
\node[below=5mm of bits,align=center]
 {Bob's array $w^\ell$\\[2pt]
  \textcolor{answergreen}{Requested bit: $w^\ell(1,3)=1$}};
\node[data,text width=29mm,anchor=north west]
 (outside) at ([yshift=-12mm]address.south west)
 {Vertex outside\\the selected edge};
\node[data,anchor=west] (zero) at ([xshift=9mm]outside.east) {output $0$};
\draw[->,semithick] (outside.east) -- (zero.west);
\end{tikzpicture}
\caption{The address selects a row and the vertex's edge position selects a
column. Their intersection gives \(w^\ell(1,3)=1\); a vertex outside the edge
bypasses the array and gives zero. The small indices illustrate the entry rule
independently of the source parameters; the arrows describe the lookup, not
a protocol's messages.}
\label{fig:yao-entry}
\end{figure}

\subsection{The protocol supplied by a colouring}

We next use a balanced colouring to build the short protocol. It shortens the
position message after Alice has learnt the edge. The intended schedule
exchanges colour and coordinate, then edge and address, then the position within
the colour, and its public cost is
\[
B_{\rm a}:=(2t-1)+2b+2(t+1)+1=2b+4t+2.
\]

\begin{proposition}\label{prop:alternating-upper}
For the table of \cref{def:alternating-table}, an $s_{\rm a}$-colouring with at
least $s_{\rm a}-1$ vertices of every colour in every edge gives
\[
D_A^{\rm alt}(M_{\rm a})\le B_{\rm a}.
\]
\end{proposition}
\begin{proof}
Fix a row $(v,c)$ and a column $(g,\alpha,\ell)$. With a balanced colouring
$\chi$, an edge contains at least
$s_{\rm a}-1$ vertices of each of the other $s_{\rm a}-1$ colours. Consequently
the vertices whose colour is $\chi(v)$ occupy at most
\[
K_{\rm a}-(s_{\rm a}-1)^2=2s_{\rm a}-1
\]
positions.
Add one message code for ``outside the edge.'' There are then at most
$2s_{\rm a}=2^{t+1}$ slots, encoded in $t+1$ bits. Fix their ordering publicly
for every colour and edge index; an unused edge index has only the outside
slot. Alice determines her slot from $v$, $\alpha$ and $\chi(v)$.

First Alice sends the $t$ bits of $\chi(v)$, while Bob sends
the $t-1$ bits of $g$ between consecutive Alice bits. This uses $2t-1$ bits and
leaves Alice knowing $g$. Next Bob sends the $b$ bits of $\alpha$, while Alice
sends the $b$ bits of $c_g$; this uses $2b$ bits, and Alice knows $\alpha$
before she sends the slot. Finally Bob sends $t+1$ dummy bits and Alice sends
the $t+1$ slot bits, after which Bob sends one answer bit. He answers zero on
the outside slot and otherwise reads the specified array entry. The cost is
\[
(2t-1)+2b+2(t+1)+1=B_{\rm a}.
\]
\end{proof}

\subsection{Few response patterns cannot cover many requests}

Bob's early messages may depend on the whole array, so one realised transcript
does not describe Alice's behaviour across columns. To analyse a competing
protocol, we therefore track Alice's responses on hypothetical Bob histories.
For a fixed row, a complete early response strategy specifies Alice's $r$th bit
after every possible history of $r-1$ Bob bits, including hypothetical
histories; her own preceding bits are then determined. Shorter branches are
completed with fixed alternating bits. For example, her first
two turns require three response bits: the initial bit and one response to each
possible first Bob bit.
In general, her first $t$ turns require
$\sum_{r=1}^t2^{r-1}=2^t-1$ response bits, giving at most
$2^{2^t-1}$ complete early strategies.

A uniformly random Boolean array matches a pattern on $k$ distinct positions
with probability $2^{-k}$. We need only the weaker bound
$2^{-k}+\varepsilon$, which permits a small explicit list. More precisely,
for a labelled list of Boolean arrays
$\mathcal W=(w^\ell)$ on a finite position set $\Omega$, we ask that its
probability law $\nu$ satisfy
\begin{equation}\label{eq:alternating-pattern}
\Pr_{\ell\sim\nu}[w^\ell|_I=z]\le2^{-|I|}+\varepsilon
\qquad
(I\subseteq\Omega,\ z\in\{0,1\}^I).
\end{equation}
A union bound controls any family of patterns. All probabilities below use
this original law, including after we restrict the columns.

The lookup proof keeps a surplus of requests through a prefix of the protocol.
The remaining response tables then describe too few patterns to cover all
retained arrays.

\begin{observation}\label{ex:pattern-mass}
Two remaining Bob--Alice pairs permit at most eight output patterns. On eight
distinct array positions, these patterns have original-law mass at most
$8(2^{-8}+\varepsilon)\le1/16$ when $\varepsilon\le2^{-8}$.
\end{observation}

The next lemma uses $m$ for the total number of Alice turns, and $t$ for the
number whose complete response strategies agree across the rows.

\begin{lemma}[Alternating lookup]\label{yao:small-bias-escape}
Let $m>t\ge1$ be integers with $m-t\ge\ceil{\log m}$. Suppose the array law
$\nu$ satisfies \cref{eq:alternating-pattern} with
$\varepsilon\le2^{-4m}$. If $J\subseteq\Omega$ has at least
$2^{m-t+1}$ positions, then the lookup table
\[
L_J(j,\ell):=w^\ell(j)
\]
has no Alice-first alternating protocol of depth at most $2m-1$ whose
rows share one complete response strategy for their first $t$ turns.
\end{lemma}
\begin{proof}
Suppose such a protocol exists. Pad every shorter branch to depth $2m-1$ with
fixed alternating bits. At the final Alice node, her row and the transcript
determine the old output leaf, so replace her last bit by that output. This
uses no extra bit and, since $m>t$, preserves her first $t$ responses.

Because Alice's first $t$ responses
are common to all rows after every Bob history, the first $2t-1$ bits partition
only the columns, into at most $2^{t-1}$ classes. Retain a class of mass at
least $2^{-(t-1)}$, keeping every row.

Put
\[
h=\ceil{\log m},\qquad Q=2^h,\qquad i=m-t-h\ge0.
\]
For the next $i$ Bob--Alice pairs,
retain a Bob child with at least half the current column mass and an Alice
child with at least half the current rows. The resulting rectangle $X\times Y$
satisfies
\[
|X|\ge2^{m-t+1-i}=2Q,\qquad
\nu(Y)\ge2^{-(m-h-1)}.
\]
Exactly $h$ Bob--Alice pairs remain. Bob's complete suffix response table has
one entry for every hypothetical history of earlier Alice bits, hence
$1+2+\cdots+2^{h-1}=Q-1$ entries. Fixing that table and Alice's maps determines
her final output on every retained row. On any $2Q$ of those rows, there are
therefore at most $2^{Q-1}$ output patterns.

Since $m\le Q<2m$, we have $\varepsilon\le2^{-2Q}$. Correctness,
\cref{eq:alternating-pattern}, and a union bound now give
\[
\nu(Y)\le2^{Q-1}(2^{-2Q}+\varepsilon)\le2^{-Q}.
\]
This contradicts the retained-mass bound because
$Q\ge m>m-h-1$.
\end{proof}

The pattern bound holds for every set of positions, so the same array list
works for every proposed protocol. The law conditioned on a transcript need
not satisfy that bound. The finite-field construction in
\cref{sec:alternating_implementation} supplies the list we need.

\begin{corollary}\label{yao:array-list}
For every fixed $t\ge4$, a nonempty labelled array list on $\Omega_{\rm a}$
whose uniform law satisfies \cref{eq:alternating-pattern} with
\[
\epsilon_{\rm a}:=2^{-4(b+3t)}
\]
can be generated in deterministic time polynomial in $H_{\rm a}$.
\end{corollary}

\subsection{From a hypergraph to an explicit gap}

We colour each vertex by a common early strategy used by many of its address
tuples. A monochromatic source edge then supplies enough distinct requests to
apply the lookup lower bound.

\begin{proposition}\label{yao:lift}
Use the complete address list and the explicit array list of
\cref{yao:array-list} in $M_{\rm a}$. A balanced source colouring gives
\[
D_A^{\rm alt}(M_{\rm a})\le B_{\rm a},
\]
whereas the absence of a proper $S_{\rm a}$-colouring gives
\[
D_A^{\rm alt}(M_{\rm a})\ge B_{\rm a}+2t-8.
\]
For each fixed $t$, the full table can be printed in deterministic time
polynomial in $|V|+H_{\rm a}$.
\end{proposition}
\begin{proof}
The upper bound is \cref{prop:alternating-upper}. For the lower bound, put
$m=b+3t-3$, so $2m-1=B_{\rm a}+2t-9$, and suppose that an Alice-first protocol
has depth at most $2m-1$. Pad it to that depth. For each vertex $v$, partition
its address tuples by Alice's complete response strategy on her first $t$ turns.
There are at most $2^{2^t-1}\le S_{\rm a}$ strategies by the complete response-table
count above. A largest class $C_v$ therefore has
\[
|C_v|\ge H_{\rm a}^{q_{\rm a}}/S_{\rm a}.
\]
Colour $v$ by this class's strategy. The source promise gives an edge
$e_\alpha$ whose vertices all receive the same strategy.

For $v\in e_\alpha$ and $g\in[q_{\rm a}]$, let
$J_{v,g}=\{c_g:c\in C_v\}$ be the set of addresses occurring in coordinate
$g$. Since $C_v$ is contained in the product of these projections,
\[
\prod_{g=1}^{q_{\rm a}}|J_{v,g}|\ge |C_v|
\ge H_{\rm a}^{q_{\rm a}}/S_{\rm a}.
\]
AM--GM, summed over the vertices of $e_\alpha$, gives a coordinate $g$ with
\[
\begin{aligned}
N:=\sum_{v\in e_\alpha}|J_{v,g}|
&\ge K_{\rm a}H_{\rm a}S_{\rm a}^{-1/q_{\rm a}}\\
&=\frac{K_{\rm a}H_{\rm a}}4=2^{m-t+1},
\end{aligned}
\]
where $S_{\rm a}=4^{q_{\rm a}}$.
For each $v\in e_\alpha$ and $u\in J_{v,g}$, choose a tuple $c(v,u)\in C_v$
with $c(v,u)_g=u$. Retain the corresponding rows and all columns
$(g,\alpha,\ell)$. Their entries are
\[
M_{\rm a}((v,c(v,u)),(g,\alpha,\ell))
=w^\ell(u,\iota_\alpha(v)).
\]
These are $N$ distinct lookup positions: different vertices occupy different
positions in the edge, and different addresses give different first coordinates.
All retained rows share the chosen early strategy. To apply
\cref{yao:small-bias-escape}, note that $m-t\ge t+1\ge5$ and
$m\le2(m-t)-1\le2^{m-t}$, so $m-t\ge\ceil{\log m}$; also
$\epsilon_{\rm a}\le2^{-4m}$. The resulting lookup restriction contradicts
that lemma. Integer depths give
$D_A^{\rm alt}(M_{\rm a})\ge2m=B_{\rm a}+2t-8$.

There are $|V|H_{\rm a}^{q_{\rm a}}$ rows and
$q_{\rm a}H_{\rm a}|\mathcal W|$ columns. For fixed $t$, the number
$q_{\rm a}$ is constant and \cref{yao:array-list} constructs the arrays in
polynomial time. Enumerating all tuples and evaluating the displayed entry rule
therefore prints the entire table in polynomial time.
\end{proof}

\begin{proof}[Proof of \cref{yao:main,yao:bpp}]
Given a fixed requested gap $k\ge1$, choose $t=\ceil{k/2}+4$. Then
$2t-8=2\ceil{k/2}\ge k$.
Compose the source reduction with \cref{yao:lift}. It gives, in deterministic
polynomial time, a table with YES cost at most $B_{\rm a}$ and NO cost at least
$B_{\rm a}+k$, proving \cref{yao:main}.

An additive-$u$ estimate of any cost in \cref{yao:comparison} is, by the
triangle inequality, an additive-$(u+2)$ estimate of $D_A^{\rm alt}$ with the
same success probability. Choose any fixed integer $k>2(u+2)$ and apply
\cref{obs:gap-estimation} to the reduction just proved. This gives
$\mathsf P=\NP$ for a deterministic estimator and $\NP\subseteq\BPP$ for a
bounded-error randomised estimator, proving \cref{yao:bpp}.
\end{proof}

The address length $q_{\rm a}=2^{t-1}$ and the source reduction's degree may
both grow with $t$. Thus the construction is polynomial for each fixed gap;
this argument supplies no growing-error theorem for strict alternation.

\section{Remaining approximation questions}
\label{sec:conclusion}

The program has moved from controlling communication in interlaced
problems to encoding exact hardness, amplifying a gap, and preserving it
in small explicit tables. This yields fixed-error NP-hardness in both
models and quantitative growing-error hardness for unrestricted protocols.
The remaining limitation is how the explicit table size grows with the
requested gap. Improving that dependence would give stronger approximation lower bounds.

Our deterministic construction has the upper bound $N^{2^{O(k)}}$ on its
size. Can smaller lists preserve the restriction bounds needed for the
next round, or can another construction improve the $\log\log T$ error scale?

For strict alternation, the construction is polynomial for each fixed
gap, but its polynomial may depend on that gap. A construction with
better control of this dependence could give growing-error hardness
in Yao's model as well.

\appendix
\section{Interlacing estimates and the lookup seed}\label{sec:interlacing-details}
\label{sec:interlacing_details}

We prove the restriction profiles and forced Alice prefix needed in
\cref{src:profiles,src:near-full} using the estimates below.

\subsection{Projection and probability mass}

Fix a Boolean matrix $M:X\times Y\to\{0,1\}$ with positive input laws $\mu$ and
$\nu$. Every mass below is measured under these original laws or their products.
For a nonempty $K\subseteq Y^p$, define its logarithmic mass loss by
$\Delta_p(K)=\log(1/\nu^p(K))$. The empty-coordinate projection has loss zero.
For $p\ge1$ and $0<x,z\le1$, let $\mathcal B(M;p,x,z)$ contain the restrictions
of $\Int{p}{M}$ retaining base row mass at least $x$ in each block and product
column mass at least $z$, and put
\[
\mathcal D_M(p;x,z):=\min_{G\in\mathcal B(M;p,x,z)}\D(G).
\]

\begin{lemma}\label{cl:projection}
Let $K\subseteq Y^p$ be nonempty, and write $\pi_JK$ for its projection onto the
coordinates in $J\subseteq[p]$. Then
\begin{align*}
 \Delta_{\card J}(\pi_JK)&\le\Delta_p(K),\\
 \Delta_{\card{J_0}}(\pi_{J_0}K)+\Delta_{\card{J_1}}(\pi_{J_1}K)
 &\le\Delta_p(K) &&(J_0\sqcup J_1=[p]).
\end{align*}
For every integer $1\le l\le p$, some $l$-element set $J$ satisfies
\[
\Delta_l(\pi_JK)\le\frac lp\Delta_p(K).
\]
Each projected matrix, with arbitrary retained rows in its selected blocks, is
isomorphic as a Boolean matrix to a restriction of the original matrix.
Consequently, for admissible
$p'\le p$, $x'\le x$, and $z'\le z$,
\[
\mathcal D_M(p';x',z')\le\mathcal D_M(p;x,z).
\]
\end{lemma}

\begin{proof}
The first two inequalities follow from
$K\subseteq\pi_J^{-1}(\pi_JK)$ and
$K\subseteq\pi_{J_0}K\times\pi_{J_1}K$, using the product law.

For the averaged bound, let $Z$ have law $Q=\nu^p(\cdot\mid K)$.
For probability laws $P,V$ on a finite set, their relative entropy is
$D_{\rm KL}(P\|V)=\sum_zP(z)\log(P(z)/V(z))$, and entropy is
$\mathsf H(P)=-\sum_zP(z)\log P(z)$, with zero terms omitted.
Conditioning on $K$ gives
\[
D_{\rm KL}(Q\|\nu^p)=\Delta_p(K).
\]
Since $Q_J$ is supported on $\pi_JK$, nonnegativity of relative entropy gives
\[
\Delta_{|J|}(\pi_JK)\le D_{\rm KL}(Q_J\|\nu^J).
\]
For a uniformly chosen $l$-subset $J$, the entropy chain rule gives
\[
\mathbb E_J\mathsf H(Z_J)
\ge\mathbb E_J\sum_{i\in J}\mathsf H(Z_i\mid Z_1,\ldots,Z_{i-1})
=\frac lp\mathsf H(Z).
\]
Also each coordinate contributes to
$-\mathbb E\log\nu^J(Z_J)$ with probability $l/p$. Subtracting the entropy
inequality therefore yields
\[
\mathbb E_J\Delta_l(\pi_JK)
\le\mathbb E_JD_{\rm KL}(Q_J\|\nu^J)
\le\frac lpD_{\rm KL}(Q\|\nu^p)=\frac lp\Delta_p(K).
\]
Some $J$ satisfies the desired bound. The conditioned law $Q$ need not be a
product; all projection masses above use the original product law.

For each tuple in $\pi_JK$, choose one tuple in $K$ above it. The entries on the
selected row blocks depend only on the selected coordinates, so these
representatives realise the projected Boolean matrix
(\cref{obs:local-decoding}). Their original weights need not be preserved:
the projection estimate uses $\nu^J$, while representatives certify the
communication lower bound. Selecting $p'\le p$ blocks leaves projected mass
at least $z$ and base row mass at least $x$, proving the final monotonicity.
\end{proof}

\subsection{Two or three surviving blocks}

The first-message arguments needed below use only the two thresholds in the
restriction profile. We state their three small cases together.

\begin{lemma}[Small block bounds]\label{lem:small-blocks}
Suppose $\RP{M}{x}{H}$ with $H\ge2$. The following restrictions of a
full-product interlace have the indicated lower bounds, whenever the displayed
row threshold is at most one. Row mass is measured separately in each block.
\[
\begin{array}{c|c|c|c}
\text{blocks}&\text{row mass per block}&\text{product column mass}&\text{cost}\\
\hline
2&2x&9/50&H\\
2&2x&9/25&H+1\\
3&4x&9/50&H+1
\end{array}
\]
\end{lemma}
\begin{proof}
For the first row, \cref{cl:projection} gives a singleton column projection
of mass at least $\sqrt{9/50}>3/10$. Its profile cost is at least $H-1\ge1$,
so the restriction is nonconstant. If Alice speaks first, follow this block
into a child retaining base row mass at least $x$; the suffix costs at least
$H-1$. If Bob speaks first, his denser child has column mass at least $9/100$,
and hence a singleton projection of mass at least $3/10$. Its suffix again
costs at least $H-1$. Including the first bit gives $H$.

For the second row, the first bound already gives nonconstancy. If Alice
speaks, a singleton projection has column mass at least $3/5$; a child
retaining half this block's rows costs at least $H$. If Bob speaks, his denser
child has column mass at least $9/50$, so the first bound gives the same suffix
cost. Including the first bit gives $H+1$.

For the third row, projecting onto any two blocks and applying the first bound
excludes a constant restriction. If Alice speaks, assign each block to a child
retaining half its rows. One child receives two blocks of mass at least $2x$,
and the first bound gives suffix cost $H$. If Bob speaks, his denser child
retains column mass at least $9/100$. The averaged projection inequality gives
two of its three coordinates with product mass at least
\[
(9/100)^{2/3}>9/50.
\]
Their row masses are still at least $4x\ge2x$, so the first bound again gives
suffix cost $H$. Including the first bit proves the claim. All projected
restrictions are realised by surviving representatives as in
\cref{cl:projection}; their masses use the original product laws.
\end{proof}

\begin{proof}[Proof of \cref{lem:product-blocks}]
A full product is exactly eight-wise independent. Apply \cref{det:exact} with
$a=\log u$: the hypotheses are precisely $u/2+1$ retained blocks, base row mass
at least $ux$, and tuple mass at least $9/25$. This includes $u=2$ and gives
$H+\log u$ bits.
\end{proof}

At the exact base-cost budget, an early Bob query leaves too many surviving
copies for the remaining depth. The next theorem formalises this forced prefix
and identifies the dominant block in each resulting rectangle. It applies both
to full products and to exact tuple lists. Its only additional input is the
exact block bound \cref{det:exact}, whose proof uses \cref{lem:small-blocks}
and the pair-projection estimate, independently of this theorem.

\begin{theorem}\label{cl:separation}
Suppose $\RP{M}{\rho}{d}$ with $d\ge2$. Let $u\ge2$ be dyadic and
$1\ge\xi>2u\rho$. Interlace $M$ at width $u$ using an exactly eight-wise independent
tuple list with the base-column marginal, as in \cref{def:exact-tuples}.
Every restriction retaining base row mass at least $\xi$ in every block
and tuple-label mass at least $18/25$, with a protocol of depth at most
$d+\log u$, has a complete prefix of $\log u$ Alice queries.
Its $u$ final parts are in bijection with the blocks: each part retains more
than $\xi-2(u-1)\rho$ of its assigned block and less than $2\rho$ of every
other block. Full products are a special case, and fixed queries count as
queries.
\end{theorem}

\begin{proof}
At depth $r<\log u$, maintain $\kappa=u2^{-r}$ assigned blocks at each node,
each retaining mass at least $\xi2^{-r}>2\kappa\rho$, with the original
retained column set unchanged. The root has this property. A singleton
projection has base-column mass at least $18/25>3/5$, by the exact marginal.
One retained label per projected value realises that base restriction, so
the profile excludes a leaf. The inequality $\xi>2u\rho$ also ensures
that every row threshold used below is at most one.

If Bob speaks, one child retains column mass at least $9/25$ and all
$\kappa$ assigned blocks. By \cref{det:exact} with $a=\log\kappa$,
it needs $d+\log\kappa$ bits, but has only
$d+\log u-r-1=d+\log\kappa-1$ left. Thus Alice speaks. Assign each block
to a child retaining at least half its mass. If a child received more than
$\kappa/2$ blocks, it would contain at least $\kappa/2+1$ blocks of mass
$\xi2^{-r-1}>\kappa\rho$, and the same lemma would again exceed its budget.
Both children therefore receive exactly $\kappa/2$ blocks and remain live,
proving the invariant and the complete prefix. In particular, a fixed Alice
query is impossible, since it would assign every block to one child.

At depth $\log u$, each part has one assigned block of mass greater than
$2\rho$. Two blocks of mass at least $2\rho$ would require $d+1$ bits by
\cref{det:exact} with $a=1$, whereas only $d$ remain. Every other block thus has
mass less than $2\rho$. Each original block is assigned to exactly one part;
its mass elsewhere totals less than $2(u-1)\rho$, proving the final bound.
All column masses in this argument are measured in the original tuple law;
no independence after restriction is asserted.
\end{proof}

\subsection{Absorbing column loss}

For a restriction $G$ of $\Int{p}{M}$ with surviving columns $K$,
$\Delta_p(K)=\log(1/\nu^p(K))$ records the loss created by several Bob
messages. The following induction charges this loss to surviving copies and uses
only lower bounds on the base matrix; exact costs enter later through a matching
selector protocol.

\begin{lemma}\label{cl:defect-induction}
Let $C\ge3$, $0<\theta\le1/4$, and $0<\delta\le1/64$.
Suppose
\[
\mathcal D_M(1;\theta,1-\delta)\ge C,\qquad\RP{M}{\theta}{C-1}.
\]
Let $e\in\{0,1\}$, let $a\ge2e$ be an integer with $2^a\theta\le1$, and let
$p>2^{a-1}$ be an integer. If a restriction $G$ of $\Int{p}{M}$ retains base
row mass at least $2^a\theta$ in every block and its nonempty column set $K$ satisfies
\[
\Delta_p(K)\le(p-2^{a-1})\delta+\frac{4e}{3},
\]
then $\D(G)\ge C+a-e$. Here $C$ is only a lower bound; equality $\D(M)=C$ is not
assumed.
\end{lemma}

\begin{proof}
We prove the lower bound by simultaneous induction on
\[
\mathfrak m(a,e):=2a+1-e.
\]
The induction tracks the number of surviving blocks and the logarithmic
column-density loss that they can absorb. Alice's bit halves row thresholds;
Bob's bit increases the defect by at most one, after which \cref{cl:projection}
supplies a suitable projection.

The profile $\RP{M}{\theta}{C-1}$ and $C-1\ge2$ allow
\cref{det:exact} on full products. Its cases $a=1,2$ give
\[
\mathcal D_M(2;2\theta,9/25)\ge C,\qquad
\mathcal D_M(3;4\theta,9/25)\ge C+1.
\]
The numerical margin used to invoke these bounds is
\[
3\delta+\frac43\le\frac{265}{192}<\frac75<-\log(9/25).
\]
The last inequality follows from \(2^7\cdot9^5<25^5\).

The initial induction states are as follows. For \((a,e)=(0,0)\), the averaged
projection bound in \cref{cl:projection} gives a singleton \(J\subseteq[p]\)
with
\[
\Delta_1(\pi_JK)\le\frac{p-1/2}{p}\delta<\delta.
\]
Since \(2^{-v}\ge1-v\) for \(0\le v\le1\), the projected column density is
greater than \(1-\delta\), and the near-full hypothesis gives cost at least
\(C\).

For \((a,e)=(2,1)\), the defect allowance and \cref{cl:projection} give a
three-coordinate projection with defect at most
\[
\frac3p\left((p-2)\delta+\frac43\right)
\le3\delta+\frac43<-\log(9/25).
\]
Its column mass exceeds \(9/25\), and each block retains mass at least
\(4\theta\). The three-copy bound above gives cost at least
\(C+1=C+2-1\).
There is one further base case, \((a,e,p)=(3,1,5)\). Its defect is at most
\(\delta+4/3<-\log(9/25)\), and its rows have mass at least \(8\theta\).
Apply \cref{det:exact} with $a=3$ to these five blocks, giving $C+2$ bits.

For \((a,e)=(1,0)\), a singleton projection has defect at most
\((p-1)\delta/p<\delta\), so its column density is greater than \(1-\delta\);
the near-full hypothesis makes this restriction nonconstant. After a Bob bit,
the denser child has a two-coordinate projection of defect at most
\[
\frac2p\left((p-1)\delta+1\right) \le2\delta+1<3\delta+\frac43,
\]
so the two-copy bound gives cost at least \(C\) after that bit. Every other
valid state has \(a\ge2\), \(p\ge3\), and base row mass at least \(4\theta\); its
three-coordinate projection has defect at most \(3\delta+4/3\), so the
three-copy bound makes the restriction nonconstant before the first-bit
analysis.

Consider such a non-base state and a protocol node for the restriction. If Alice
speaks, assign every block to a child retaining at least half its row mass. For
child \(i\in\{0,1\}\), let \(p_i\) be the number of assigned blocks and let
\(L_i\) be the defect after projecting the retained column set onto those
blocks. Then \(p_0+p_1=p\), and the partition inequality in \cref{cl:projection}
gives \(L_0+L_1\le\Delta_p(K)\). Put
\[
s_i=\max\{p_i-2^{a-2},0\}.
\]
Since \(s_0+s_1\ge p-2^{a-1}>0\), at least one \(s_i\) is positive. The
allowances of the children with \(s_i>0\) satisfy
\[
 \begin{aligned}
\sum_{i:s_i>0}\left(s_i\delta+\frac{4e}{3}\right)
 &\ge (p-2^{a-1})\delta+\frac{4e}{3}\\
&\ge\sum_{i:s_i>0}L_i.
 \end{aligned}
\]
At least one eligible child contributes the constant \(4e/3\); the
final inequality uses the assumed defect bound on \(K\) and
\(L_0+L_1\le\Delta_p(K)\). Consequently some child with \(s_i>0\) meets its
allowance. Each assigned block retains mass at least \(2^{a-1}\theta\), and
there are \(p_i>2^{a-2}\) such blocks. This is state \((a-1,e)\); the exceptional
state \((a,e)=(2,1)\) was already proved. The induction hypothesis gives cost at least
\(C+a-1-e\) in that child, and Alice's bit gives \(C+a-e\). The same argument
includes the transition \((1,0)\to(0,0)\).

If Bob speaks, one child has at least half the current column mass, so its defect
is at most \(\Delta_p(K)+1\). When \(e=0\) and \(a\ge2\), this child satisfies
state \((a,1)\), because
\[
\Delta_p(K)+1\le(p-2^{a-1})\delta+\frac43.
\]
Its induction cost is at least \(C+a-1\), and adding Bob's bit proves the claim.

It remains to handle \(e=1\) and \(a\ge3\). The five-block base case already
handles \(p=5\), so \(p\ge6\). Put \(p'=\ceil{p/2}\). Then
\[
\frac12\le\frac{p'}p\le\frac47,\qquad
\frac{p'}p(p-2^{a-1})\le p'-2^{a-2},\qquad
\frac47\left(\frac43+1\right)=\frac43.
\]
By the averaged projection bound in \cref{cl:projection}, some \(p'\)
coordinates of the denser child therefore have defect at most
\[
\frac{p'}p\left((p-2^{a-1})\delta+\frac43+1\right)
\le(p'-2^{a-2})\delta+\frac43.
\]
The projected child retains base row mass at least \(2^{a-1}\theta\) in each
of its \(p'>2^{a-2}\) blocks, so it is state \((a-1,1)\). Its induction cost is at
least \(C+a-2\); adding Bob's bit gives \(C+a-1\), as required.
\end{proof}

\begin{proof}[Proof of \cref{lem:many-product-blocks}]
Let \(G\) be the given restriction. Project to \(p=9u/16\) retained blocks and
apply \cref{cl:defect-induction} with
\[
p=\frac{9u}{16},\qquad a=\log u,\qquad \delta=\frac1{64},\qquad \theta=x,\qquad
C=H,\qquad e=0.
\]
Since \(u\ge16\) and \(ux\le1\), we have \(x\le1/16\le1/4\),
\(2^a\theta=ux\le1\), and \(p=9u/16>u/2=2^{a-1}\). The profile \(\RP{M}{x}{H}\)
implies \(\RP{M}{x}{H-1}\), and its \(j=0\) case gives \(\mathcal
D_M(1;x,1-\delta)\ge H\) because \(1-\delta=63/64\ge3/5\). Thus all hypotheses
of \cref{cl:defect-induction} are available.

The retained column density \(2^{-u/1024}\) gives defect at most \(u/1024\). For
the chosen parameters,
\[
(p-2^{a-1})\delta =\left(\frac{9u}{16}-\frac u2\right)\frac1{64}
=\frac{u}{1024}.
\]
If more than \(9u/16\) blocks are retained, select exactly \(p\) of them. By
\cref{cl:projection}, projecting the retained column set to these blocks does
not increase its defect, and surviving representatives realise an actual
restriction. Applying \cref{cl:defect-induction} to this projected restriction
gives
\[
\D(G)\ge C+a=H+\log u.
\]
The original restriction contains this projected restriction, so its cost is at
least the same by the monotonicity conclusion of \cref{cl:projection}.
\end{proof}

\begin{proof}[Proof of \cref{cl:bridge}]
Put $C=\D(W)$ and $f=\ceil{b/\delta}$. The assumptions imply
\[
f\delta\ge b,\qquad (f-1)/Q<1/80,\qquad f<Q/40,\qquad Q\ge5120.
\]
Indeed, $b/\delta\ge64$, $Q\delta\ge80b$, and
$f\le b/\delta+1<2b/\delta$. Fix $j\in\{0,1\}$ and put
\[
a_j=r-j,\qquad p_j=2^{a_j-1}+f,\qquad t_j=2^{a_j}\theta.
\]
These choices give $a_j\ge0$, $2^{a_j-1}<p_j<Q$, and $t_j<1$.

Transpose a profile restriction of $R^+$ to $\Int{Q}{W}$. It retains tuple
mass at least $2^{-b}$ and total block-row mass at least $(3/5)2^{-j}$.
If fewer than $p_j$ blocks had base row mass at least $t_j$, total row mass
would be less than
\[
\frac{p_j-1}{Q}+t_j
<2^{-j-1}+\frac1{80}+\frac1{20\cdot2^j}
\le\frac35\,2^{-j},
\]
a contradiction; the last inequality uses $j\le1$. Select $p_j$ qualifying
blocks $J$. For the retained tuple set $K$, \cref{cl:projection} gives
\[
\Delta_{p_j}(\pi_JK)\le\Delta_Q(K)\le b
\le f\delta=(p_j-2^{a_j-1})\delta.
\]
The projected restriction therefore meets \cref{cl:defect-induction} with
$e=0$: the near-full and weaker profile hypotheses are those for $W$,
$\theta<1/4$ follows from $Q\theta<1/20$, and the row threshold is $t_j$.
It requires $C+a_j=C+r-j$ bits. Representatives realise it as a Boolean
restriction of the original matrix, so transposing back proves
$\RP{R^+}{2^{-b}}{C+r}$.

The selector protocol costs $r+\D(W)$ bits, giving equality on the whole
matrix and the claimed transfer of any upper bound. By definition the
transpose has $n^Q$ rows and $Qm$ columns.
\end{proof}

\subsection{The initial lookup matrix and template layers}

The lookup seed starts from $R_0(z,j)=z_j$ on $\{0,1\}^{p_0}\times[p_0]$, and
rank after deleting query coordinates supplies the restriction profile needed
for amplification. Its proof forces leaves of each output value and then uses
the elementary query-and-answer protocol for equality.

\begin{proof}[Proof of \cref{cl:seed}]
We first lower-bound the rank of every retained lookup submatrix, then count
leaves for the two output values, and finally record the exact size and upper
bounds of the seed.

Write \(a=\log p_0\). Fix \(j\in\{0,1\}\), a retained row set
\(V\subseteq\{0,1\}^{p_0}\) with \(\card V\ge2^{p_0-L}\), and a retained query
set \(J\subseteq[p_0]\) with \(\card J\ge\ceil{(3/5)2^{-j}p_0}\). Let \(v\) be
the dimension of the real span of \(V\). A \(v\)-dimensional subspace has \(v\)
coordinate functionals whose restrictions form a basis, so projection to those
\(v\) coordinates is injective on the subspace. Since the Boolean cube has at
most \(2^v\) points in those coordinates, \(V\) has at most \(2^v\) points. The
lower bound on \(\card V\) therefore implies \(v\ge p_0-L\).

Deleting the \(p_0-\card J\) coordinates outside \(J\) can lower the row rank by
at most \(p_0-\card J\). Thus the restricted lookup matrix has rank \(r\)
satisfying
\[
r\ge \card J-L \ge \frac35\,2^{-j}p_0-L >2^{a-j-1}.
\]
The strict inequality follows from
\[
\frac35\,2^{-j}p_0-2^{a-j-1} =\frac{p_0}{10\cdot2^j}\ge\frac{p_0}{40}>L,
\]
where \(p_0/40-L>0\) is part of the choice in the lemma statement.

Let \(F\) denote the restricted lookup matrix. Consider a protocol for \(F\).
The indicator of each 1-leaf rectangle has rank at most one, and the sum of
these indicators is the restricted matrix, so there are at least \(r\)
one-leaves. Let \(\mathbf1\) be the all-ones matrix; it has rank one. The
complement of the restricted matrix therefore has rank at least \(r-1\), by
\(\rank(\mathbf 1-F)\ge\rank(F)-\rank(\mathbf1)\), so there are at least \(r-1\)
zero-leaves. Hence
\[
\#\text{leaves}\ge r+(r-1)=2r-1>2^{a-j}.
\]
A depth-\(t\) binary protocol has at most \(2^t\) leaves, so the restriction has
depth at least \(a+1-j\). At \(j=0\), this gives \(\D(R_0)\ge a+1\). Conversely,
Bob sends the query coordinate in \(a\) bits and Alice sends the queried bit, so
\(\D(R_0)\le a+1\); hence \(\D(R_0)=a+1=d_0\). The unit-vector rows form the
identity submatrix, so \(\rank(R_0)=p_0=2^{d_0-1}\).

By minimality of the dyadic choice, \(p_0<128(L+1)\). If \(L=\ceil{\log n}\) and
\(n\ge4\), then \(L+1\le\log n+2\le2\log n\), so \(p_0<256\log n\) and
\(2^{p_0}<n^{256}\). Enumerating the \(2^{p_0}\) words and the \(p_0\) query
coordinates therefore prints the table deterministically in polynomial time.
\end{proof}

\begin{proof}[Proof of \cref{src:profiles}]
Apply the block-averaging observation to the three interlaces in order.
Fix $j\in\{0,1\}$,
put $\beta_j=(3/5)2^{-j}$, and write $\lambda=255/256$ for the template mass.
After transposition, the relevant data are
\[
\begin{array}{c|c|c|c|c}
\text{target}&\text{base}&q&u&\text{average block retention}\\ \hline
A&R&q_1=s-2&s2^{-j}&\beta_j\\
R_{\rm temp}&A&q_2&q_2 2^{-j}&\beta_j\\
N_{\rm ref}&R_{\rm temp}&s&s2^{-j}&(\beta_j-(1-\lambda))/\lambda
\end{array}
\]
The retained tuple mass is at least $x$ in each case. For $N_{\rm ref}$,
discard the non-template rows after transposition and project each reference
label onto its $J_2^s$ selector tuple, forgetting its tail. The displayed
average retention follows by subtracting non-template mass; projection keeps
mass at least $x$ and representatives preserve the template entries.

Count blocks at thresholds $u\rho$, $ux$, and $ux$, respectively.
The bound $q\beta\le h+qt$ gives the following lower bounds for
their counts divided by $u$:
\[
\frac35(1-2/s)-q_1\rho,\qquad
\frac35-q_2x,\qquad
\frac{2^j(\beta_j-(1-\lambda))}{\lambda}-sx.
\]
All three exceed $9/16$ by \cref{lem:full-schedule}. The same schedule gives
$u\ge16$, row thresholds at most one, and
$x=2^{-10L}\ge2^{-u/1024}$. Applying \cref{lem:many-product-blocks}
successively therefore gives costs at least
\[
d+\log s-j=H_1-j,\qquad
H_1+\log q_2-j=H_2-j,\qquad
H_2+\log s-j=C-1-j.
\]
These are the three asserted profiles. The $j=0$ cases, together with the
selector protocols of costs $d+\log s$ and $H_1+\log q_2$, give
$\D(A)=H_1$ and $\D(R_{\rm temp})=H_2$.

Finally, the first $q_1$ positions and three distinct tails give
$q_1+3=s+1=(2s)/2+1$ blocks. Each retains base row mass at least $2s\rho$,
and the retained word mass exceeds $1/2>9/25$. Apply
\cref{lem:product-blocks} with parameter $2s$ and the profile of $R$ to get
$d+\log(2s)=H_1+1$ bits. Projected representatives realise the required
Boolean restriction by \cref{cl:projection}.
\end{proof}

\section{Bounds on transferred inputs}\label{sec:mass-details}
\label{sec:mass_bounds}

We prove the two receipt bounds in \cref{sec:internal-interlace} under its common
hypotheses. All parameters are written without a construction subscript.
Protected inputs belong to the diagonal family, of mass $\mu=1/256$ per source
copy. The template and ordinary families each have mass $\lambda=255/256$.

\subsection{Two selectors at one destination}

Two independently chosen selectors force an extra bit even when their inner
indices coincide. Only eight-wise independence in their original laws is needed.

\begin{lemma}[Two tagged selectors]\label{res:paired}
Let $A:\Gamma\times J\to\{0,1\}$ satisfy $\RP{A}{x}{H_1}$ with $H_1\ge2$,
and let $q_2$ be dyadic with $2q_2x\le1/8$. Take two independent labelled
sources $S_o,S_f$ of $J^{q_2}$, each exactly eight-wise independent with the
base-column marginal. Put $H_2=H_1+\log q_2$ and give
$J_2=[q_2]\times\Gamma$ its uniform-index times base-row law.
Rows are selector pairs $(z_o,z_f)\in S_o\times S_f$. Ordinary and
other-copy columns are separately tagged copies of $J_2$, with entries
\[
A(\gamma,(z_o)_\alpha)\quad\text{and}\quad
A(\gamma,(z_f)_\alpha),
\]
respectively. The tags remain distinct even when the inner indices coincide.
Retain a pair event of original product mass at least $49/64$, ordinary columns
of normalised mass at least $1-\varepsilon$, and other-copy columns of
normalised mass at least $\delta$, where
\[
\varepsilon\ge0,\qquad 0<\delta\le1,\qquad
\delta>2q_2x+\varepsilon.
\]
Then the retained matrix requires at least $H_2+1$ bits.
\end{lemma}

\begin{proof}
Set $\lambda_0=2q_2x$. Let $b$ count ordinary $\Gamma$-fibres of retained
mass below $\lambda_0$, and let $f$ count other-copy fibres of mass at least
$\lambda_0$. Averaging their masses gives
\[
1-\varepsilon\le1-\frac b{q_2}(1-\lambda_0),\qquad
\delta\le\lambda_0+\frac f{q_2}(1-\lambda_0).
\]
Consequently
\[
b\le\frac{\varepsilon q_2}{1-\lambda_0},\qquad
f\ge\frac{q_2(\delta-\lambda_0)}{1-\lambda_0}>b.
\]
Thus more than $q_2$ tagged fibres have mass at least $\lambda_0$.

Suppose a protocol has depth at most $H_2$. Follow each such fibre through
column-player queries, choosing a child retaining at least half its current
mass. At depth $r$ its mass is at least $\lambda_0 2^{-r}$, while the pair
set of rows is unchanged. Every tagged selector coordinate has projection
mass at least $49/64>3/5$. The profile of $A$ therefore excludes a leaf and
also excludes depth $\log q_2+1$: there the fibre mass is at least $x$, so
$H_1$ bits are needed but only $H_1-1$ remain. Stop each fibre at its first
row-player query.

At a stopping node of depth $r$ put $\kappa=q_2 2^{-r}$. These capacities are
positive integers, and Kraft's inequality gives total capacity at most $q_2$.
Some node therefore holds at least $\kappa+1$ tagged fibres, each of mass at
least $2\kappa x$. Its denser row child retains pair mass at least
$49/128>9/25$.

The two sources together are exactly eight-wise independent on tagged
coordinates: a request for at most eight coordinates factors between the
independent sources, and within either source uses at most eight coordinates.
After transposition, the fibres are distinct blocks of an interlace of $A$.
Keep the surviving pair labels as tuple labels, so their mass is unchanged.
Apply \cref{det:exact} with width $2\kappa$: there are at least
$(2\kappa)/2+1$ blocks of base row mass at least $2\kappa x$, and
$2\kappa x\le2q_2x\le1/8$. The lower bound is
\[
H_1+\log(2\kappa)=H_1+\log\kappa+1,
\]
whereas the child has only $H_2-r-1=H_1+\log\kappa-1$ bits left.
This contradiction includes $\kappa=1$. All masses used here are measured
in the original laws; no independence after conditioning is asserted.
\end{proof}

\begin{claim}\label{cl:two-stream}
Use the common source laws and profiles of \cref{sec:internal-interlace}, and put
$\lambda_0=2q_2x$, $\tau_0=2sx$, $C=H_2+\log s+1$, and
\[
a_{\min}=\frac{99/100+\lambda^2-1}{\lambda^2}.
\]
Retain ordinary columns of normalised $\Omega$-mass at least
$1-\varepsilon$, and separately tagged other-copy columns whose
measure-preserving map to $J_2$ has retained image mass at least $\delta$.
Renumber bins so that the distinguished bin is one. It consists of a set
$P_*\subseteq S^2$ of product mass at least $a_{\min}$; every other bin retains
$S$-mass at least $\tau_0$. Write ordinary coordinates as
$Z=(\xi,j_1,\ldots,j_s)$, with $j_p=(\alpha_p,\gamma_p)$. Require entries
\[
\begin{array}{c|cc}
 & Z\in\Omega&\text{other-copy image }(\alpha,\gamma)\\ \hline
(z_o,z_f)\in P_*&A(\gamma_1,(z_o)_{\alpha_1})&A(\gamma,(z_f)_\alpha)\\
(p,z),\ p\ne1&A(\gamma_p,z_{\alpha_p})&\text{arbitrary}.
\end{array}
\]
If $\varepsilon\ge0$, $0<\delta\le1$, and $\delta>\lambda_0+4\varepsilon$, then this matrix
requires at least $C$ bits.
\end{claim}

\begin{proof}
Suppose the restriction has a protocol of depth at most $C-1$. Track one
selector support per bin until the first column-player query. An ordinary
support follows a row child retaining at least half its current mass. The
distinguished support carries the pair set $P_*$. At depth $r$, split it into
two ordinary supports if both children have first-selector projection mass at
least $\tau_0 2^{-r-1}$; otherwise keep the pair in the qualifying child.
A discarded pair set has mass no larger than its first-selector projection,
so all unsplit steps together discard less than
\[
\sum_{r\ge0}\tau_0 2^{-r-1}=\tau_0.
\]
The common numerical bounds give
$\tau_0\le10^{-3}$ and $a_{\min}-\tau_0>49/64$. An intact pair therefore
retains mass greater than $49/64$, and every support at depth $r$ has selector
mass at least $\tau_0 2^{-r}$.

The hypothesis on $\delta$ implies $\varepsilon<1/4$, so the ordinary columns
retain normalised mass greater than $3/4$. Their projection onto any bin has
at least that mass. The profile $\RP{R_{\rm temp}}{x}{H_2}$ excludes a
support-bearing leaf and depth
$\log s+1$: at that depth a support has mass at least $x$, but only $H_2-1$
bits remain. Hence every support reaches a first column query at depth at
most $\log s$.

Give an ordinary support demand one and an intact pair demand two. Following
one child preserves demand, and splitting preserves it because $2=1+1$.
If a stopping node $v$ has $h_v$ supports, of which $z_v$ are intact pairs,
then
\[
\sum_v(h_v+z_v)=s+1,\qquad
\sum_v\kappa_v\le s,\qquad \kappa_v=s2^{-r_v}.
\]
The second inequality is Kraft's inequality. Some node therefore has
$h+z>\kappa$. Its supports have distinct bin coordinates, since descendants
in different branches never rejoin, and each has selector mass at least
$\tau_0 2^{-r}=2\kappa x$. Since $z\le h$, for dyadic $\kappa\ge2$ this gives
$h\ge\kappa/2+1$; for $\kappa=1$, either $h\ge2$ or $h=z=1$.

Handle $\kappa=1,h=z=1$ before restricting columns or projecting the pair.
Both tagged column families are still unchanged. Project ordinary columns
onto $j_1$ and other-copy columns onto their $J_2$ images, choosing surviving
representatives. These maps preserve the ambient laws, so the projected masses
are at least $1-\varepsilon$ and $\delta$. The intact row pair has mass greater
than $49/64$. Thus \cref{res:paired} requires $H_2+1$ bits, while the stopping
node has only $C-1-r=H_2$ bits left. The original law on the two selector
labels is the product of their exact selector laws, as required by that lemma.

In all other cases choose the denser column child, retaining ordinary mass
greater than $3/8>9/25$. Replace any intact pair by its first-selector
projection, choosing one surviving row for each projected label. The entries
now form an interlace of $R_{\rm temp}$ on the supported bins. Its tuple
source has exactly $8192$-wise independent bin coordinates in either
construction. Forget the unused tail while keeping whole selector-tuple labels
in the product case and whole selector-list indices in the explicit case.
Their image event has at least the retained mass in its original law; one
surviving completion per image label realises its entries. Thus
\cref{det:exact} applies in both constructions. For $\kappa\ge2$,
use width $\kappa$: the $\kappa/2+1$ blocks have row mass at least
$2\kappa x\ge\kappa x$, so the required cost is $H_2+\log\kappa$ against a
child budget of $H_2+\log\kappa-1$. For $\kappa=1$ there are at least two
blocks of mass at least $2x$; width two gives the stronger bound $H_2+1$.
The schedules ensure all these row thresholds are at most one. This completes
the contradiction without using other-copy entries outside the distinguished
bin.
\end{proof}

\subsection{Total receipt at a destination}

\begin{proof}[Proof of \cref{res:incoming}]
Put $\kappa=1$ at a final node. Suppose a destination receives protected mass
at least $K_*\kappa$. If $d_i\in[0,1]$ is the fraction of origin $i$'s diagonal
family reaching it, then $K_*=512\mu L_ws$, with $L_w=\log w+2$, gives
\[
D_0:=\sum_i d_i\ge512L_ws\kappa.
\]
Set $d_i=0$ outside the origins being counted. Discard terms below $D_0/(2w)$;
their sum is less than $D_0/2$. The remaining terms occupy at most $L_w$
dyadic classes. One class $V$, of size $m$ and upper endpoint $\delta\le1$,
therefore satisfies
\[
\delta/2<d_i\le\delta\quad(i\in V),\qquad
m\delta\ge\frac{D_0}{2L_w},\qquad
\delta\ge\frac{D_0}{2w},\qquad m\ge256s\kappa.
\]
All incoming row sets are fixed by the Alice-only transcript. At a first-Bob
node choose a child whose column event $Y^*$ has original mass at least
$\theta/2$; at a final node set $Y^*=Y'$.

For each $i\in V$, let $T_i$ indicate a template position in the metadata
stream. The centred variables $T_i-\lambda$ are $256$-wise independent in
the original law. By \cref{lem:conditioning-moment},
\[
\mathbb E\left|\sum_{i\in V}(T_i-\lambda)\right|^{256}
\le(2^{14}m)^{128}.
\]
If fewer than $m/2$ positions are template positions, the absolute sum is
greater than $(\lambda-1/2)m$. Markov's inequality and the scheduled bound give
\[
\Pr\left\{\sum_{i\in V}T_i<m/2\right\}
\le C_{128}m^{-128}\le\theta/4,
\qquad C_{128}=\left(\frac{2^{14}}{(\lambda-1/2)^2}\right)^{128}.
\]
Thus $Y^*$ intersects the event of at least $m/2$ template positions in mass
at least $\theta/4$.

Average this intersection over the complete metadata and base-label streams.
Some fixed configuration leaves at least $m/2$ template positions in $V$, with
fixed bins $p_i$, and a selector-label section $Z$ of original mass at least
$\theta/4$. Fix only those two other streams. The remaining selector stream,
flattened to coordinates $(i,\alpha)$, still has its original exactly
$8192$-wise independent $J$ marginals. Now apply \cref{det:conditioning} to $Z$
with cap one. It removes at most $e(\theta/4)\le sq_2$ coordinates.

At every fixed template position $i$, the diagonal input $(\alpha,\gamma)$
reads $A(\gamma,z_{i,\alpha})$, regardless of the chosen bin $p_i$. Its retained
$J_2$-mass is $d_i>\delta/2$. More than $\delta q_2/4$ of its $\Gamma$-fibres
have mass at least $\delta/4$: otherwise those fibres contribute at most
$\delta/4$ and the remaining fibres less than $\delta/4$ to the total.
Consequently there are more than $m\delta q_2/8$ such blocks $(i,\alpha)$.
Set $K_0=2s\kappa q_2$. The class bounds and $w^2q_2x\le1$ give
\[
\frac{m\delta q_2}{8}\ge32s\kappa q_2,
\qquad
K_0x\le\frac{2s\kappa}{w^2}\le\frac\delta4.
\]
After removing the at most $sq_2$ exceptional coordinates, more than $K_0/2$
blocks remain. Retain $K_0/2+1$ of them and all selector labels in $Z$.
Choose an actual incoming row for every retained diagonal image. Distinct
$(i,\alpha,\gamma)$ images give distinct rows, and their entries on $Z$ are
exactly $A(\gamma,z_{i,\alpha})$. The row masses use the original $\Gamma$
law; the column event has mass one in the fixed conditional law on $Z$.
The chosen completions certify these entries, not an identity of their old
weights with the projected weights.

Apply \cref{det:rareboost} with base profile $\RP{A}{x}{H_1}$ and parameter
$a=\log K_0$. Its row threshold is at most $\delta/4\le1$, and it requires
\[
H_1+\log K_0=C+\log\kappa
\]
bits. The first-Bob child has only $C+\log\kappa-2$ bits left, and a final
node only $C-1$. This contradiction proves the receipt bound for every
collection of origins, including those excluded from the owner argument.
\end{proof}

\subsection{Transfer to another origin's owner}

\begin{proof}[Proof of \cref{res:foreign}]
Let $\varepsilon_i$ be the owner's missing normalised ordinary mass and put
$\delta_i=2q_2x+4\varepsilon_i$. If $\delta_i\ge1$, the conclusion is automatic:
an origin has protected mass only $\mu$. Otherwise, a transfer exceeding
$\mu\delta_i$ supplies a diagonal set from origin $j\ne i$ of actual normalised
mass $\delta>\delta_i$. We keep this strict inequality throughout the argument.

Both origins are good coordinates for $Y'$, and only Alice has spoken before
the owner. By \cref{det:conditioning}, the projection of $Y'$ onto these two
raw-row coordinates has original product mass greater than $99/100$.
Restrict to the event that both are templates. Its mass is $\lambda^2$, so the
retained relative mass is at least
\[
a_{\min}=\frac{99/100+\lambda^2-1}{\lambda^2}.
\]
Average over the two bin labels and fix a pair of bins attaining this relative
mass. Their selector section is a set $P_*\subseteq S^2$ of original product
mass at least $a_{\min}$. Distinguish the chosen bin at $i$. The good singleton
projection at $i$ loses at most $\eta/2$ of its raw-row law; hence every other
bin retains selector mass at least
\[
1-\frac{\eta s}{2\lambda}>2sx.
\]
Choose one surviving completion in $Y'$ for every selector pair in $P_*$ and
for every retained selector in the other bins. The bin at $i$ distinguishes
the latter families from the pairs; the projected labels distinguish the
representatives within each family. These choices realise entries of an
actual restriction and do not assert independence of the law conditioned on
$Y'$.

After transposition, the ordinary columns from origin $i$ have normalised mass
$1-\varepsilon_i$. The diagonal columns from $j$ have $J_2$-image mass
$\delta>\delta_i$, since that image map preserves their original law. A pair
$(z_o,z_f)$ reads $A(\gamma_{p_i},(z_o)_{\alpha_{p_i}})$ on ordinary columns and
$A(\gamma,(z_f)_\alpha)$ on the diagonal columns. Other bins have the ordinary
template entries, and their diagonal entries are immaterial. This is the
restriction in \cref{cl:two-stream}, after renaming the distinguished bin.
The claim requires $C$ bits because $\delta>2q_2x+4\varepsilon_i$, whereas the
owner has only $C-1$ bits left. This contradiction proves the non-strict
transfer bound, including the case $\varepsilon_i=0$.
\end{proof}

\section{Conditioning and surviving blocks}
\label{sec:conditioning-details}

The higher moments below use the original law, where independence is exact.
The conditional fourth moment uses only the approximate marginals outside
the exceptional coordinates.

\subsection{A moment count and the exceptional coordinates}

\begin{lemma}\label{lem:conditioning-moment}
Let $r\ge1$ be an integer. If $X_1,\ldots,X_m$ have mean zero,
$|X_i|\le1$, and every subfamily of at most $2r$ variables is independent, then
\[
\mathbb E\left|\sum_{i=1}^m X_i\right|^{2r}\le(r^2m)^r.
\]
In particular, for centred indicators with $1024$-wise independence,
\[
\mathbb E\left|\sum_iX_i\right|^{1024}\le(2^{18}m)^{512},
\qquad
\mathbb E\left|\sum_iX_i\right|^{256}\le(2^{14}m)^{128}.
\]
\end{lemma}
\begin{proof}
The case $m=0$ is immediate. Expand the $2r$th power. A term with an index
occurring exactly once has expectation zero, by centring and independence.
Every remaining term uses at most $r$ distinct indices. There are at most
$m^rr^{2r}$ such terms: encode one by its distinct indices in increasing
order, pad that list to length $r$ by repeating its first index, and record
one of the $r$ positions for each of its $2r$ factors. This encoding is
injective if repeated padding positions are never used. Each expectation
has absolute value at most one. Taking $r=512$ and $r=128$ proves the
displayed special cases.
\end{proof}

\begin{proof}[Proof of \cref{det:conditioning}]
Put $\varepsilon=\varepsilon(v,\tau)$. Choose a maximal disjoint family
$J_1,\ldots,J_m$ of sets of at most eight coordinates whose conditional
marginals differ from their product marginals by more than $\varepsilon$.
If $m=0$, nothing must be deleted. Otherwise, total variation supplies an
indicator $f_i$ on each $J_i$ with
$\mathbb E_{\nu\mid Y}f_i-\mathbb E_\nu f_i>\varepsilon$; take its complement
if necessary. Under the original law $\nu$, any subfamily of at most $1024$
such indicators is independent, because its disjoint sets use at most
$8192$ coordinates.

For $S=\sum_i(f_i-\mathbb E_\nu f_i)$, conditional Jensen's inequality and
\cref{lem:conditioning-moment} give
\[
v(m\varepsilon)^{1024}
<\nu(Y)\mathbb E_{\nu\mid Y}|S|^{1024}
\le\mathbb E_\nu|S|^{1024}
\le(2^{18}m)^{512}.
\]
Thus $m<2^{18}\varepsilon^{-2}v^{-1/512}$.
Delete $E=\bigcup_iJ_i$, so $|E|\le8m\le e(v,\tau)$. Every remaining set of
at most eight coordinates has the desired marginal: otherwise it could be
added to the disjoint family, contradicting maximality.
\end{proof}

\subsection{Some pair has a large projection}

\begin{lemma}\label{lem:pair-projections}
Fix disjoint coordinate pairs $J_1,\ldots,J_m$ in a tuple list with marginal
$\mu$ and law $\nu$. Suppose every projection of a label set $K$ has product
mass $\mu^2(\pi_{J_i}K)<9/25$.
For an exactly eight-wise independent list,
\[
\nu(K)<\frac2{m^2}<\frac9{50m}\qquad(m\ge16).
\]
For an exactly $8192$-wise independent list, let $K\subseteq Y$ with
$\nu(Y)\ge v>0$, and choose the pairs outside the exceptional set from
\cref{det:conditioning}, with any cap $0<\tau\le1$. Then
\[
(\nu\mid Y)(K)<\frac1{4m}\qquad(m\ge16).
\]
Thus a retained set of conditional mass at least $1/(4m)$ forces one of
these pairs to have product projection mass at least $9/25$.
\end{lemma}
\begin{proof}
Write $P_i=\pi_{J_i}K$ and put
\[
p_i=\mu^2(P_i)<9/25,\qquad
X_i=\ind[z_{J_i}\in P_i]-p_i,\qquad S=\sum_{i=1}^mX_i.
\]
On $K$ every indicator equals one, so $S>(16/25)m$. Under the product law,
the centred indicators are independent, with
$\mathbb EX_i^2\le1/4$ and $\mathbb EX_i^4\le1/4$. Only terms with no
singleton index survive the fourth-power expansion. Hence
\[
\mathbb E_{\rm prod}S^4
=\sum_i\mathbb EX_i^4+6\sum_{i<j}\mathbb EX_i^2\mathbb EX_j^2
\le\frac m4+\frac6{16}\binom m2\le\frac{m^2}4.
\]
Exact eight-wise independence gives the same fourth moment under $\nu$,
since each term uses at most four pairs. Markov's inequality gives
$\nu(K)<2/m^2$, because $1/[4(16/25)^4]<2$.
For $m\ge16$, $2/m^2\le1/(8m)<9/(50m)$.

Now suppose the pairs are good for the fixed event $Y$. Every term in $S^4$
lies in $[-1,1]$ and depends on at most eight good coordinates. Total
variation changes its expectation by at most $2\varepsilon$, where
$\varepsilon=\varepsilon(v,\tau)$. Summing the $m^4$ terms and applying
Markov therefore gives
\[
p:=(\nu\mid Y)(K)
\le\frac2{m^2}+\frac{2\varepsilon}{(16/25)^4}
\le\frac2{m^2}+16\varepsilon.
\]
For the other estimate, return to the original law. Any subfamily of at most
$1024$ pair indicators is independent under $\nu$, and $K\subseteq Y$ with
$\nu(Y)\ge v$. Thus \cref{lem:conditioning-moment} gives
\[
p=\frac{\nu(K)}{\nu(Y)}
\le\frac{(2^{18}m)^{512}}{v((16/25)m)^{1024}}
\le\frac{(2^{20}/m)^{512}}v.
\]
Combining the bounds and using
$\min\{u,w\}\le u^{511/512}w^{1/512}$ yields
\[
\begin{aligned}
mp&\le\frac2m+
\min\{16\varepsilon m,2^{20\cdot512}v^{-1}m^{-511}\}
\\
&\le\frac2m+(16\varepsilon)^{511/512}2^{20}v^{-1/512}.
\end{aligned}
\]
The last term is at most
\[
2^{20-26\cdot511/512}v^{255/131072}<\frac1{32},
\]
by $\varepsilon\le2^{-30}v^{1/256}$ and $0<v\le1$. For $m\ge16$, this gives
$mp<1/8+1/32<1/4$, as required.
\end{proof}

\subsection{Small projections and short continuations}

The two bounds from \cref{lem:small-blocks} that we need require $d+1$ bits:
\[
\begin{array}{c|c|c}
\text{number of blocks}&\text{row mass per block}&\text{product column mass}\\
\hline
2&2x&9/25\\
3&4x&9/50.
\end{array}
\]
Here $\RP{M}{x}{d}$, $d\ge2$, and the indicated row threshold is at most
one. For selected coordinates $J$, choose one surviving label for each tuple
in $\pi_JK$. By \cref{obs:local-decoding}, these representatives transfer
the projected full-product lower bound to the labelled matrix. The projected
mass is measured under $\mu^J$, even when the labels have a conditional law.

\begin{proof}[Proof of \cref{det:exact}]
Suppose a protocol has depth at most $d+a-1$. At each Alice message,
assign every current block to a child retaining at least half its row mass,
and keep a child assigned at least half the blocks. At each Bob message,
keep a child retaining at least half the current column mass under the
original law $\nu$. Follow these choices for $a-1$ messages, padding an
early leaf with dummy Alice messages if needed. If $A$ messages are
Alice's and $B$ are Bob's, then $A+B=a-1$ and the remaining restriction has
\[
\begin{gathered}
\text{at least }2^B+1\text{ blocks},\qquad
\text{row mass per block at least }2^{B+1}x,\\
\nu(K)\ge(9/25)2^{-B},\qquad
\text{at most }d\text{ bits remaining}.
\end{gathered}
\]
The block count follows by rounding up
$(2^{a-1}+1)/2^A>2^B$. Every set $J$ of at most eight surviving coordinates
has $\mu^J(\pi_JK)\ge\nu(K)$ by exact marginals.

If $B=0$, two blocks meet the first threshold above; if $B=1$, three meet
the second. If $B=2$, five coordinates have joint product mass at least
$9/100$, so \cref{cl:projection} supplies a pair of mass at least
$(9/100)^{2/5}>9/25$. If $B=3$ or $B=4$, eight coordinates have mass at
least $9/400$, so some pair has mass at least $(9/400)^{1/4}>9/25$.
If $B\ge5$, choose $m=2^{B-1}\ge16$ disjoint pairs from the $2m+1$
surviving blocks. Were every pair projection smaller than $9/25$,
\cref{lem:pair-projections} would give
$\nu(K)<9/(50m)=(9/25)2^{-B}$, a contradiction.
In each case the selected blocks meet one of the two row thresholds above,
which are at most one because $B+1\le a$. They require $d+1$ further bits,
contradicting the available $d$.
\end{proof}

The conditioned version needs a short-continuation estimate for each possible
number of Bob messages.

\begin{lemma}\label{det:suffix}
Suppose $d\ge2$ and $\RP{M}{x}{d}$. Let $\nu$ be an exactly $8192$-wise
independent tuple source with the base-column law, let $\nu(Y)\ge v>0$, and
choose the good coordinates from \cref{det:conditioning} with cap
$0<\tau\le1$. If an integer $b\ge0$ satisfies $2^{b+1}x\le1$, then any
restriction retaining at least $2^b+1$ good blocks, each of base row mass at
least $2^{b+1}x$, and a column set $K\subseteq Y$ with
$(\nu\mid Y)(K)\ge2^{-b-1}$ requires at least $d+1$ bits.
\end{lemma}
\begin{proof}
Keep $Y$ and its good coordinates fixed. Set
$\varepsilon=\varepsilon(v,\tau)\le2^{-30}$. For any set $J$ of at most
eight good coordinates, \cref{det:conditioning} gives
\[
\mu^J(\pi_JK)\ge(\nu\mid Y)(K)-\varepsilon
\ge2^{-b-1}-\varepsilon.
\]
For $b\le4$ this is at least $(9/25)2^{-b}$, since the smallest margin is
$7/800>2^{-30}$. The same small-coordinate checks in the proof of
\cref{det:exact} therefore supply the required pair or triple.
For $b\ge5$, choose $m=2^{b-1}\ge16$ disjoint good pairs. Since
$(\nu\mid Y)(K)\ge1/(4m)$, \cref{lem:pair-projections} supplies a pair of
product projection mass at least $9/25$.

Every selected pair retains row mass at least $2x$; the triple at $b=1$
retains at least $4x$. These thresholds are at most one by
$2^{b+1}x\le1$. Thus \cref{lem:small-blocks} requires $d+1$ bits in every
case. Throughout, retained label masses use the single law $\nu\mid Y$;
no independence after further restriction is asserted.
\end{proof}

\begin{proof}[Proof of \cref{det:rareboost}]
Suppose a protocol has depth at most $d+a-1$. At each Alice message, assign
every current block to a child retaining at least half its row mass, then
keep a child assigned at least half the blocks. At each Bob message, keep
a child with at least half the column mass in the fixed law $\nu\mid Y$.
The good coordinates remain good throughout.

Follow these choices for $a-1$ messages, padding an early leaf with dummy
Alice messages if needed. If $A$ messages are Alice's and $B$ are Bob's,
then $A+B=a-1$ and the surviving restriction has
\[
\begin{aligned}
\text{number of good blocks}
&\ge\ceil{(2^{a-1}+1)/2^A}\ge2^B+1,\\
\text{row mass per block}&\ge2^{a-A}x=2^{B+1}x,\\
(\nu\mid Y)(K')&\ge2^{-B-1}.
\end{aligned}
\]
Apply \cref{det:suffix} with $b=B$. Its row threshold is at most one because
$B+1\le a$. It requires $d+1$ further bits, while only $d$ remain.
\end{proof}

\section{Constructing rational input lists}
\label{sec:support-algorithm}

The construction follows the sequential compression idea of
\citet{koller1994sample}: append an independent coordinate, then replace the
resulting distribution by one with fewer positive weights and the same
small marginals. The compression uses only a linear dependence among
zero--one vectors.

\begin{proof}[Proof of \cref{det:code}]
Fix the independence order $\mathsf D$ and put $P=(1+qm)^{\mathsf D}$. Start
with the empty word of weight one. At each step, append an independent
$\mu$-coordinate to every word currently retained. On the resulting candidate
set $\mathcal U\subseteq Y^i$, the weights satisfy
\[
\sum_{u\in\mathcal U}w_u\ind[u_J=a]
=\prod_{j\in J}\mu(a_j)
\qquad
(J\subseteq[i],\ |J|\le\mathsf D,\ a\in Y^J).
\tag{*}\label{eq:support-cylinders}
\]
The empty set $J$ gives $\sum_uw_u=1$. Appending the coordinate preserves old
marginals, and a marginal involving the new coordinate factors by independence.

There are at most $P$ equalities in \eqref{eq:support-cylinders}, since the
number of assignments on at most $\mathsf D$ coordinates is at most
$\sum_{j=0}^{\mathsf D}(qm)^j\le(1+qm)^{\mathsf D}$. For a candidate word $u$,
let $a(u)$ be its zero--one column in this system. Discard zero weights.
Whenever the remaining columns are dependent, compute a nonzero integer
relation $\sum_uc_ua(u)=0$ and set
\[
t=\min_{c_u>0}\frac{w_u}{c_u},
\qquad
w'_u=w_u-tc_u.
\]
The constant-one row gives $\sum_uc_u=0$, so the relation has both signs.
The update preserves every equality, leaves all weights nonnegative, and
makes at least one positive weight zero. Discard the new zero weights and
repeat until the remaining columns are independent. There are then at most
$P$ words, so the next extension
has at most $mP$ candidates. After $q$ stages, the equalities give the required
$\mathsf D$-wise independence.

To bound the arithmetic, choose each relation using a column basis and one
further column. If the basis has size $r\le P$, select $r$ independent rows.
These span the row space, so a relation on them holds in every row.
The determinant minors of the resulting $r\times(r+1)$ zero--one matrix give
the relation, with integer coefficients of absolute value at most
$\max\{1,r^{r/2}\}$ by Hadamard's bound. Thus each coefficient uses
$O(P\log(P+1))$ bits, and exact linear algebra finds it in polynomial time.
Fixed orders for the columns, rows, and tied pivots make these choices
deterministic.

Let $L_\mu$ be the input law's total binary encoding length and $Q_\mu$ the
product of its denominators. Maintain a common denominator $Q$ for the current
weights. An extension multiplies $Q$ by at most $Q_\mu$. If $w_u=a_u/Q$
and $j$ attains the minimum ratio, one elimination gives
\[
w'_u=\frac{a_uc_j-a_jc_u}{Qc_j}.
\]
It adds at most $O(P\log(P+1))$ denominator bits. There are at most $qmP$
eliminations and $q$ extensions, so every intermediate weight has bit length
\[
O\bigl(qL_\mu+qmP^2\log(P+1)\bigr).
\]
Numerators obey the same bound because the weights lie in $[0,1]$.
All linear algebra, ratio comparisons, and updates therefore have polynomial
cost for fixed $\mathsf D$. Together with the at most $P$ length-$q$ words,
the same bound proves polynomial total output length.
\end{proof}

\section{Printing a sample of the implicit table}\label{sec:sampling-algorithm}
\label{app:implementation}

We prove the access guarantee by counting complete input labels first, then
implementing their elementary draws with fixed blocks of fair bits. Although
one entry visits only one recursive entry at each level, generating its input
labels must also produce all the coordinates that it does not visit.

\begin{proof}[Proof of \cref{lem:sample-access}]
Fix the admissible graph, put $L=\ceil{\log n}$, and use the same schedule at
every level. Write $R:X\times Y\to\{0,1\}$ for the current base matrix.

\emph{Labels and entries.}
The two interlaces $U=\Int{w}{N}$, $W=U^{\mathsf T}$ and
$R^+=(\Int{Q}{W})^{\mathsf T}$ give the successor labels
\[
\mathbf a=((i_t,Z_t))_{t\in[Q]},\qquad
\mathbf b=(j,(k_i)_{i\in[w]}),
\]
where $i_t\in[w]$ and $j\in[Q]$ are uniform, $Z_t$ are tagged raw columns,
and $k_i\in\mathcal K$ are raw rows. All are independent in their prescribed
laws. The entry rule is
\[
R^+(\mathbf a,\mathbf b)=B_{\rm raw}(k_{i_j},Z_j).
\]
For a template raw row $(p,z)$ and raw column
$Z=(\xi,j_1,\ldots,j_s)$, with $j_p=(\alpha_p,\gamma_p)$ and
$z_{\alpha_p}=(r,u)$, this entry is $R(u,(\gamma_p)_r)$.
For an item row $(a,u)$, it is $R(u,\xi_{\pi_\alpha(a)})$ when all activity
indices equal $\alpha$ and $v_a(\alpha)=1$, and $R(u,\xi_{s+1})$ otherwise.
Thus, after reading the stored selectors and activity data, evaluation makes
exactly one base-entry call.

Choose a template row with probability $255/256$; vertex and auxiliary types
each have probability $1/512$. A template row uses a uniform bin and
$q_2$ independent labels from $J=[q_1]\times X$; an item row uses a uniform
item of its type and one base row.

For a raw column, choose $\Omega$ with probability $255/256$ and
$\mathcal D$ otherwise. For $\Omega$,
draw $s+1$ independent base columns for $\xi$ and $s$ independent labels
from $J_2=[q_2]\times Y^{2s-1}$. For $\mathcal D$, draw one
$j=(\alpha,\gamma)\in J_2$, copy it into all $s$ positions, and put
$\xi=\operatorname{tail}(\gamma)$. These are the actual full-product labels;
no extra streams used in the lower-bound proof are generated.

The recursive draw counts are therefore
\[
\begin{array}{c|c|c}
\text{label}&\text{base-row draws}&\text{base-column draws}\\ \hline
\text{raw template row}&q_2&0\\
\text{raw item row}&1&0\\
\Omega\text{ column}&0&(s+1)+s(2s-1)=2s^2+1\\
\mathcal D\text{ column}&0&2s-1\\
\text{successor row}&0&\le Q(2s^2+1)\\
\text{successor column}&\le wq_2&0
\end{array}
\]
Selectors, copying and activity-list scans have polynomial cost in
$n+s+q_2+w+Q+L+1$, with no colouring required. Every width is $2^{O(L)}$ by
\cref{lem:full-schedule}. Hence common computable integer bounds $J_i$ for
label length, generation time, elementary draw count and entry-evaluation
time can be chosen with
\[
J_0=O(L),\qquad J_{i+1}\le2^{c_0L}J_i,\qquad
J_k=2^{O((k+1)L)},
\]
where $c_0$ is fixed and effective. Here an elementary draw is charged one
operation; its bit implementation is accounted for below. At the lookup seed,
a row consists of $p_0=O(L)$ independent bits, a column is a uniform index in
$[p_0]$, and an entry can be read in $O(L)$ time. This proves the initial bound
and ties the resource estimate to the seed iteration.

Since $h=O(L)$ and $d_0=O(L)$, we have $B_k=O((k+1)L)$. The direct sample
size in \cref{thm:compression} then gives $m=2^{O((k+1)L)}$.
Generate and store all $2m$ labels, then evaluate their $m^2$ intersections.
Including the stored tuples and printed entries, this costs
\[
O((2m+m^2)J_k)=2^{O((k+1)L)}=n^{O(k+1)}.
\]
All constants are independent of $k$.

\emph{Fixed-bit implementation.}
Let $N_{\rm draw}\ge1$ be a computable upper bound on the elementary draws
for all $2m$ labels. The preceding bounds give
$N_{\rm draw}=2^{O((k+1)L)}$. A tape-independent bound on every elementary
alphabet size and mixture denominator is
\[
A=\max\{512,4s(n+s),q_2,w,Q,p_0\}=2^{O(L)}.
\]
It covers type draws, item indices, bins, activity and interlace selectors,
and seed draws. All base inputs are generated recursively.
Set $b=\ceil{\log(200A N_{\rm draw})}$ and reserve $N_{\rm draw}$ blocks of
$b$ fair bits. For a draw from $[a]$, read the next block as a uniform integer
$U\in\{0,\ldots,2^b-1\}$ and return
\begin{equation}\label{eq:rounded-draw}
\operatorname{Draw}(a)=1+\floor{aU/2^b}.
\end{equation}
For a rational mixture, draw from its denominator and map consecutive
intervals to types with their prescribed numerator lengths. Ignore unused
blocks.

Every result lies in $[a]$. Each value has either $\floor{2^b/a}$ or
$\ceil{2^b/a}$ preimages, so the draw's total-variation distance from uniform
is at most
\[
a/2^b\le A/2^b\le1/(200N_{\rm draw}).
\]
Mapping to a mixture type cannot increase this distance. Conditional on a
common history of the ideal and implemented generators, couple their next
draw with disagreement probability at most this bound. Fresh blocks permit
this even when the history determines the next alphabet. Before the first
disagreement the generators take the same branches; summing over its at most
$N_{\rm draw}$ possible positions bounds the joint total variation by
$1/200<1/100$.

All elementary choices are legal, so the generation and copying rules assemble
actual inputs on every tape. There are exactly $m$ labels of each kind and a
fixed reserved tape length. Entry evaluation is exact. Finally,
$b=O((k+1)L)$, and each rounded draw uses arithmetic on
$O(b+\log A)$-bit integers. This polynomial overhead preserves the bound
$n^{O(k+1)}$ for time and random bits on every tape, completing the proof.
\end{proof}

\section{Effective amplification widths}
\label{sec:parameters}

The widths must leave room for source coordinates, deleted inputs, and the
final restoration step. We use generous margins. Their exact exponents do
not enter the conclusions, except that the full-product overhead must remain
$L+O(\log L)$.

Write $n=|V|$, $\lambda=255/256$, $\mu=1/256$, and $\zeta=\gamma/24$.
The same value $\lambda$ is the mass of the raw template rows and, separately,
of the ordinary column family. The source promise gives
\[
 n/2\le |E|\le\Delta n/2,\qquad
 \delta_v=\frac{\gamma|E|}{2\Delta n}\ge\frac{\gamma}{4\Delta}.
\]
To make the conditioning budget easy to compute, for an integer $b\ge0$ and rational
$0<\tau\le1$ define
\[
\overline e(2^{-b},\tau)
 :=\ceil{2^{90}\tau^{-2}2^{\ceil{b/64}}},\qquad
\overline e(v):=\overline e(v,1).
\]
The definition in \cref{det:conditioning} gives
\[
e(v,\tau)
\le8+2^{23}\tau^{-2}v^{-1/512}+2^{83}v^{-5/512}
\le2^{84}\tau^{-2}v^{-1/64}.
\]
Thus $e(v,\tau)\le\overline e(v,\tau)$ for dyadic $v$, and
$\overline e(v,\tau)=O(\tau^{-2}v^{-1/64})$. Both schedules use
\begin{equation}\label{eq:common-budgets}
\begin{gathered}
\delta_U=\eta/(16w),\qquad \theta=(512Q)^{-1},\\
e_0=\overline e(\theta,\eta/2),\qquad E_{\rm exc}=2e_0,\qquad
K_*=512\mu(\log w+2)s=2(\log w+2)s,\\
C_{128}=\left(\frac{2^{14}}{(\lambda-1/2)^2}\right)^{128}.
\end{gathered}
\end{equation}
Here $E_{\rm exc}$ bounds exceptional destination capacity and $K_*$
bounds incoming protected mass per unit capacity, as in
\cref{res:destinations,res:incoming}.

In both schedules $s,q_2,w,Q$ are powers of two and $q_1=s-2$.
The following common tests pay for the source restrictions and the shared
internal-interlace proof:
\begin{equation}\label{eq:width-margins}
\begin{gathered}
 \text{Source coordinates and collision survival:}\\
 s\ge256,\qquad
 s\eta\le2^{-16}\min\{1,\delta_v,\mu\zeta\},\\
 q_2\ge\operatorname{pow}_2(4s|E|),\qquad
 4s^2\rho\le1,\qquad 8sq_1\le\zeta q_2/16,\\
 \text{Protected inputs and transfers:}\\
 \max\{q_2x,sx,wx\}\le10^{-4},\qquad
 2w^2x\le\eta/(16\lambda),\\
 w^2q_2x\le1,\qquad 4\mu wq_2x\le\eta/32,\\
 \text{Conditioning, owner loss, and incoming mass:}\\
 \overline e(\theta)\le w/128,\qquad
 \overline e(\theta/4)\le sq_2,\qquad
 C_{128}s^{-128}\le\theta/4,\\
 E_{\rm exc}\le w/4,\qquad K_*E_{\rm exc}\le\eta w/16.
\end{gathered}
\end{equation}
As usual, $\operatorname{pow}_2(a)$ is the least power of two at least $a$.

Several further bounds follow automatically. No isolated vertices gives
$q_2\ge4s|E|\ge2sn\ge2(n+s-3)$, since
$sn-n-s+3=(s-1)(n-1)+2>0$. Also
$s\rho\le1/(4s)\le1/1024<10^{-3}$, and
$2wx+\eta/(16\lambda)<1$ follows from $wx\le10^{-4}$ and $s\eta\le2^{-16}$.
For $j=0,1$ and $\beta_j=(3/5)2^{-j}$, the source population ratios obey
\[
\begin{aligned}
\frac35(1-2/s)-q_1\rho
 &\ge\frac35-\frac6{5\cdot256}-10^{-3}>\frac9{16},\\
\frac35-q_2x&\ge\frac35-10^{-4}>\frac9{16},\\
\frac{2^j(\beta_j-(1-\lambda))}{\lambda}-sx
 &\ge\frac{44}{75}-10^{-4}>\frac9{16}.
\end{aligned}
\]
The same conservative ratio $44/75=((3/5)-4(1-\lambda))/\lambda$
applies to the reference profile. After conditioning, its surviving fraction
is at least
\[
\frac{44}{75}-wx-\frac{4e(\theta)}w
\ge\frac{44}{75}-\frac1{256}-\frac1{32}>\frac12.
\]
For item deletions and discarded coordinates, the common tests also give
\[
\begin{gathered}
4096\eta sn\le\delta_v n/16<(s/2)\delta_v n,\qquad
4096\eta s(s-3)\le(s-3)/16<s/2,\\
8sq_1+4s\eta q_2/\mu
\le(1/16+2^{-14})\zeta q_2<\zeta q_2/4,\qquad
2s\rho\le1/(2s)\le1/8.
\end{gathered}
\]
Template bins retain more than $999/1000$ of their mass because
$\eta s/\lambda<10^{-3}$. Projecting the retained ordinary columns onto
whole selector labels preserves mass at least $1-\eta/\lambda>99/100$.
Consequently the forced-prefix argument has the required margins
$999/1000>2sx$ and $999/1000-2(s-1)x>998/1000$.
The remaining fixed comparisons are
\[
\begin{gathered}
\frac{99}{200}-2(1-\lambda)>\frac{12}{25},\qquad
\frac{99/100+\lambda^2-1}{\lambda^2}-10^{-3}>\frac{49}{64},\\
\frac1{64}+\frac{\mu}{2\lambda}
=\frac1{64}+\frac1{510}<\frac1{32}.
\end{gathered}
\]
The last inequality bounds all foreign receipts by $\eta w/32$:
the transfer threshold is $2q_2x+4\varepsilon_i$, and
$\sum_i\varepsilon_i\le\eta/(8\lambda)$.
Together with the exceptional allowance $\eta w/16$, this is less than
the outgoing mass $\eta w/2$.

Finally, the flag test comes directly from \cref{lem:conditioning-moment}.
For $m\ge s$ template indicators of mean $\lambda$ with $256$-wise independence,
\[
\Pr[\text{fewer than }m/2\text{ template flags}]
\le\frac{(2^{14}m)^{128}}{((\lambda-1/2)m)^{256}}
=C_{128}m^{-128}\le\theta/4.
\]
The cap $\eta/2$ at an owner also leaves each other selector bin with mass
at least $1-\eta s/(2\lambda)>999/1000>2sx$.
These checks apply to both schedules.

\subsection{Full products}

The source-coordinate width contributes the only leading term to the
logarithmic overhead. All other widths are powers of $\log n$.

\begin{lemma}\label{lem:full-schedule}
Put $L=\ceil{\log n}$ and $\ell=\ceil{\log(L+1)}$. Choose
\[
\begin{aligned}
s&=2^{16+\ell},&q_1&=s-2,&q_2&=2^Ls^2,\\
w&=2^{8\ell},&\rho&=2^{-L},&x&=2^{-10L},\\
\eta&=\frac{\min\{1,\delta_v,\mu\zeta\}}{2^{16}s},&
Q&=w^2,
\end{aligned}
\]
and use \eqref{eq:common-budgets}. Call the source admissible when
\eqref{eq:width-margins} holds and
\[
Q\delta_U\ge80L,\qquad L\ge16,\qquad d_0\le L,\qquad h\le7L,
\]
where $d_0$ is the seed cost in \cref{cl:seed} and
\[
h=2\log s+\log q_2+\log w+\log Q=L+64+28\ell.
\]
Every sufficiently large promised source is admissible, with an effective
finite exceptional range depending only on $\Delta,\gamma$. All choices and
tests have polynomial bit complexity, depend on neither the base matrix nor
the round count, and give $h=L+O_{\Delta,\gamma}(\log L)$.
\end{lemma}

\begin{proof}
All implicit constants below are effective and depend only on the source
promise. Since $L+1\le2^\ell<2(L+1)$,
\[
s=\Theta(L),\quad \eta=\Theta(L^{-1}),\quad
q_2=2^L\Theta(L^2),\quad w=\Theta(L^8),\quad
Q=\Theta(L^{16}),\quad \theta=\Theta(L^{-16}).
\]
The conditioning and receipt budgets therefore satisfy
\[
\begin{gathered}
\overline e(\theta),\overline e(\theta/4)=O(L^{1/4}),\qquad
E_{\rm exc}=O(L^{9/4}),\qquad K_*=O(L\log L),\\
\frac{K_*E_{\rm exc}}{\eta w}=O(L^{-15/4}\log L),\qquad
\frac{s^{-128}}\theta=O(L^{-112}),\qquad
\frac{L}{Q\delta_U}=O(L^{-6}).
\end{gathered}
\]
Each budget is eventually below its tested allowance. For the remaining
smallness tests, the losses divided by their variable allowances are
\[
\begin{array}{c|c@{\qquad}c|c}
\text{quantity}&\text{upper bound}&\text{quantity}&\text{upper bound}\\ \hline
s^2\rho&O(L^2 2^{-L})&sq_1/q_2&O(2^{-L})\\
\max\{q_2x,sx,wx\}&O(L^8 2^{-9L})&
w^2q_2x&O(L^{18}2^{-9L})\\
w^2x/\eta&O(L^{17}2^{-10L})&
wq_2x/\eta&O(L^{11}2^{-9L}).
\end{array}
\]
The source coordinates fit because
$4s|E|\le2\Delta s2^L\le q_2$ once $s\ge2\Delta$.
Since $q_2$ is dyadic, this also proves its required dyadic ceiling.
The product-block width tests follow directly from $q_2\ge s$ and
$s/4096\ge16(L+1)>10L$.

For restoration,
\[
0<\delta_U\le2^{-20}/(sw)<1/64,\qquad Q\theta=1/512<1/20.
\]
Together with $Q\delta_U\ge80L$, these are the width conditions in
\cref{cl:bridge}. All logarithms defining $h$ are integers, and substitution
gives its displayed exact value. The seed has $d_0=\ell+7$, so $h\le7L$,
$d_0\le L$, and $L\ge16$ also hold eventually.

These estimates give a uniform effective cutoff. For $\ell\ge2$, the bounds
$2^{\ell-2}\le L<2^\ell$ turn every inverse-power or logarithmic error above
into $C(\ell+1)^r2^{-a\ell}$, and every $CL^d2^{-L}$ into
$C2^{d\ell-2^{\ell-2}}$, with effective rational $C$, fixed $a>0$, and fixed
nonnegative integers $r,d$. Each decreases to zero beyond a computable
point. Searching beyond those points until all tolerances hold gives a
cutoff valid at every larger size. The actual tests use only rational
arithmetic, integer powers and ceilings, with operands of polynomial bit
length. No approximation of a fractional power is required.
\end{proof}

\subsection{Explicit supports}

For explicit supports we can afford widths polynomial in the graph size.
This also leaves room for the separate restoration argument.

\begin{lemma}\label{lem:generated-schedule}
Fix a rational constant
\[
0<c_{\rm del}\le2^{-16}\min\{1,\gamma/(4\Delta),\mu\zeta\}.
\]
For a dyadic $\mathsf N\ge n$, choose
\[
\begin{gathered}
L_{\rm g}=\log\mathsf N,\qquad
s_{\rm g}=2^{\ceil{L_{\rm g}/4}},\qquad q_{1,\rm g}=s_{\rm g}-2,\qquad
q_{2,\rm g}=\mathsf N^4,\\
w_{\rm g}=\mathsf N^2,\qquad Q_{\rm g}=\mathsf N^6,\qquad
\rho_{\rm g}=\mathsf N^{-1},\qquad x_{\rm g}=\mathsf N^{-8},\qquad
\eta_{\rm g}=c_{\rm del}/s_{\rm g}.
\end{gathered}
\]
Use \eqref{eq:common-budgets} with these subscripted quantities. Call
$\mathsf N$ admissible when \eqref{eq:width-margins} holds and
\[
\overline e(x_{\rm g})\le s_{\rm g}/128,\qquad
Q_{\rm g}\ge128\bigl(\overline e(\rho_{\rm g},\delta_{U,\rm g}/2)+1\bigr).
\]
There is an effective dyadic $\mathsf N_0$, depending only on the graph
promise and $c_{\rm del}$, such that every dyadic $\mathsf N\ge\mathsf N_0$
passes uniformly for graphs with $n\le\mathsf N$. The first passing dyadic
$\mathsf N\ge n$ is therefore at most $2\max\{n,\mathsf N_0\}$.
The tests have polynomial bit complexity and ensure every numerical
hypothesis of the explicit-support construction and restoration.
\end{lemma}

\begin{proof}
Omit subscripts in the proof. All implicit constants are effective and depend
only on the fixed graph promise and $c_{\rm del}$. We have
\[
\mathsf N^{1/4}\le s<2\mathsf N^{1/4},\qquad
\eta=\Theta(\mathsf N^{-1/4}),\qquad
\delta_U=\Theta(\mathsf N^{-9/4}),\qquad
\theta=2^{-9}\mathsf N^{-6}.
\]
The conservative budgets fit with room to spare:
\[
\begin{array}{c|c|c}
\text{budget}&\text{upper bound}&\text{available width}\\ \hline
\overline e(x)&O(\mathsf N^{1/8})&s=\Theta(\mathsf N^{1/4})\\
\overline e(\theta)&O(\mathsf N^{1/8})&w=\mathsf N^2\\
\overline e(\theta/4)&O(\mathsf N^{1/8})&sq_2=\Theta(\mathsf N^{17/4})\\
\overline e(\rho,\delta_U/2)+1&O(\mathsf N^5)&Q=\mathsf N^6\\
E_{\rm exc}&O(\mathsf N^{3/4})&w=\mathsf N^2.
\end{array}
\]
For the two nonconstant caps the exponents are
$1/2+6/64<3/4$ and $9/2+1/64<5$. Since
$K_*=O(\mathsf N^{1/4}\log\mathsf N)$,
\[
\frac{K_*E_{\rm exc}}{\eta w}
=O(\mathsf N^{-3/4}\log\mathsf N)\longrightarrow0.
\]
Also $s^{-128}/\theta=O(\mathsf N^{-26})$, proving the flag test.
The remaining vanishing losses are
\[
\begin{array}{c|c@{\qquad}c|c}
\text{quantity}&\text{upper bound}&\text{quantity}&\text{upper bound}\\ \hline
s^2\rho&O(\mathsf N^{-1/2})&
sq_1/q_2&O(\mathsf N^{-7/2})\\
\max\{q_2x,sx,wx\}&O(\mathsf N^{-4})&
w^2x/\eta&O(\mathsf N^{-15/4})\\
wq_2x/\eta&O(\mathsf N^{-7/4})&
\eta&O(\mathsf N^{-1/4}).
\end{array}
\]
The remaining product condition is exact: $w^2q_2x=1$, which is allowed.
The choice of $c_{\rm del}$ supplies the deletion inequality. Source
coordinates fit because $4s|E|\le2\Delta s\mathsf N=O(\mathsf N^{5/4})$,
whereas $q_2=\mathsf N^4$.

To check the source-profile consumers explicitly, their block parameter
is $u=s2^{-j}$ or $q_2 2^{-j}$ for $j=0,1$, so $u\ge s/2\ge128$.
The population fraction exceeds $9/16$, and removing at most
$e(x)\le\overline e(x)\le s/128\le u/32$ coordinates leaves more than
$u/2$ blocks. In the first profile $u\rho\le s\rho\le1$; in the other two,
$ux\le\max\{q_2x,sx\}\le1$. The common checks above supply the outer
profile, incoming-mass, owner, and transfer margins.

For restoration,
\[
0<\delta_U\le2^{-20}/(sw)<1/64,\qquad
w\delta_U=\eta/16,\qquad Q\theta=1/512<1/20.
\]
Together with the tested lower bound on $Q$, these are the hypotheses of
\cref{det:restore} with $\phi=\rho$ and $\delta=\delta_U$.
All widths are polynomial in $\mathsf N$, and the independence order stays
fixed at $8192$. By \cref{det:code}, the fixed number of support constructions
in one round has a fixed polynomial cost in $\mathsf N$ and the complete
base encoding, including its rational weights; its exponent is independent
of the round number. The overhead is
\[
h_{\rm g}=2\log s+\log q_2+\log w+\log Q
=12L_{\rm g}+2\ceil{L_{\rm g}/4}.
\]

Every asymptotic comparison above has a positive power gap. Its error ratio
is bounded by $C(1+L_{\rm g})2^{-aL_{\rm g}}$ for effective rational $C$
and fixed rational $a>0$. These bounds decrease to zero once
$L_{\rm g}+1\ge2/a$; searching beyond this point gives a uniform effective
cutoff. Rational exponents can be removed by taking a fixed integer power.
The actual admissibility tests use only integer exponents, rational
arithmetic and ceilings, on operands with $O(L_{\rm g})$ bits.
Searching successive dyadic sizes from $n$ therefore takes polynomial time
and finds a passing size by $2\max\{n,\mathsf N_0\}$.
None of these choices uses the base matrix, its optimum, or the number of rounds.
\end{proof}

\section{Colouring promises from SAT}
\label{sec:sources}

\subsection{Almost-linear bounded-degree graph colouring}

The graph source must preserve a constant colouring gap and the almost-linear
size of sparse SAT. Both properties survive the standard encoding; we give the
accounting because the size bound determines the final approximation ratio.

\begin{proof}[Proof of \cref{src:gap-colouring}]
Apply Dinur's PCP construction, amplification, and a final bounded-degree
reduction~\cite[Lemma~8.3, Theorem~1.5, and Lemma~1.7]{dinur2007pcp}.
After removing isolated variable blocks, the resulting constraint system has a
fixed alphabet, maximum degree $d$, perfect completeness, and a fixed gap
$\rho>0$: every assignment on a NO input rejects at least $\rho|F|$ constraint
occurrences. Here $F$ counts occurrences, including repetitions, and the number
of variable blocks is at most $2|F|$.

Encode each alphabet value by one Boolean variable. For each block impose
exactly one true value, and for each constraint occurrence forbid its rejected
value pairs. Convert the clauses to width three with private auxiliary
variables. A true clause has a satisfying extension; a false clause leaves at
least one converted clause false. Fixed alphabet and degree give fixed
constants $C,k$ such that the resulting formula has $m\le C|F|$ clauses,
$n\le3m$ used variables, and at most $k$ occurrences of each variable.
In particular, $m\le kn$. Perfect completeness is preserved.

On a NO input, decode every valid block and assign arbitrary values to invalid
blocks. At least $\rho|F|$ constraints are rejected.
If half have valid endpoints, their separate forbidden-pair clauses give at
least $\rho|F|/2$ false clauses. Otherwise at least
$\rho|F|/(2\max\{1,d\})$ distinct blocks are invalid, and each falsifies one of
its own validity clauses. Thus every Boolean assignment violates at least
$\epsilon m$ clauses for some fixed rational $0<\epsilon\le1$.

The bounded-occurrence colouring construction gives a loopless graph $G_0$ of
fixed maximum degree, with at most $C_{\rm g}n$ edges, and preserves
completeness~\cite[Theorem~14 and its proof, pp.~98--99]{bogdanov2002lower}.
We use the proof's global decoding estimate: if $r\le kn/8$ edges are deleted,
a proper colouring of the remaining graph yields an assignment violating at
most $8r$ clauses. For any colouring of a NO graph, delete its $r$
monochromatic edge occurrences. If $r<\epsilon m/8$, then
$r<kn/8$, so decoding contradicts the Boolean gap. Consequently
\[
r\ge\epsilon m/8,\qquad
|E(G_0)|\le C_{\rm g}n\le3C_{\rm g}m.
\]
The monochromatic-edge fraction is therefore at least
$\delta:=\epsilon/(24C_{\rm g})>0$.

Remove isolated vertices and add a disjoint edge, obtaining $H$. On a NO
instance $G_0$ has an edge, so the added edge reduces the gap by at most a
factor two. Choose a fixed rational $0<\gamma\le\delta/2$ and a fixed degree
bound $\Delta$ for $H$. Trivial inputs, including empty formulas, can instead
use fixed promise graphs, decreasing $\gamma$ and increasing $\Delta$ if
necessary. The graph is now nonempty, loopless, and has no isolated vertices.
Keep the construction order on labels and count parallel edge occurrences
separately.

For a sparse formula on $N$ variables, the binary input length is
$L_{\rm enc}=O(N\log(N+2))$. Dinur's starting PCP has size
$L_{\rm enc}(\log L_{\rm enc})^{O(1)}$ and inverse-polylogarithmic gap.
The $O(\log\log L_{\rm enc})$ amplification rounds each multiply size by a
constant; the final degree reduction and both encodings have constant-factor
cost~\cite[Theorem~8.1 and its proof; Lemma~1.7]{dinur2007pcp}. Hence
\[
|V(H)|,|E(H)|\le C' N(\log(N+2))^q
\]
for fixed $C',q$. Taking
$\max\{1,\ceil{N/|E(H)|}\}$ disjoint labelled copies gives a graph $G$ with
\[
N\le |E(G)|\le N+|E(H)|,\qquad
|V(G)|/2\le |E(G)|\le\Delta|V(G)|/2.
\]
Copies preserve completeness, the gap, and the graph conventions. Enlarging a
fixed integer $p$ absorbs the constants and gives
$|V(G)|,|E(G)|\le N(\log(N+2))^p$. All steps are deterministic and
polynomial-time.
\end{proof}

\subsection{Balanced edges with distinct vertices}

The balanced source counts vertex occurrences within an edge. The lookup
construction needs distinct vertices, so we replace repeated occurrences by
distinct labelled clones.

\begin{theorem}[Guruswami--Lee, {\cite[Theorem~1.1]{guruswami2018balanced}}]
\label{yao:gl}
For fixed integers $Q_{\rm GL},k_{\rm GL}\ge2$ and $\epsilon_{\rm GL}>0$, it is
NP-hard to distinguish $Q_{\rm GL}k_{\rm GL}$-uniform hypergraphs admitting a
$k_{\rm GL}$-colouring with at least $Q_{\rm GL}-1$ occurrences of each colour
in every edge from hypergraphs having no independent set of measure at least
$\epsilon_{\rm GL}$.
\end{theorem}

\begin{proof}[Proof of \cref{yao:source}]
Use the parameters
\[
s_{\rm a}=2^t,\qquad K_{\rm a}=s_{\rm a}^2,\qquad
S_{\rm a}=2^{s_{\rm a}}
\]
and apply \cref{yao:gl} with
$Q_{\rm GL}=k_{\rm GL}=s_{\rm a}$ and
$\epsilon_{\rm GL}=1/(2S_{\rm a})$.
The YES case has the required balanced $s_{\rm a}$-colouring. In the NO case
a proper $S_{\rm a}$-colouring would have an independent colour class of
measure at least $1/S_{\rm a}$, contradicting the source promise.

Replace each vertex $u$ by $K_{\rm a}S_{\rm a}$ labelled clones. For each
ordered occurrence edge $(u_1,\ldots,u_{K_{\rm a}})$, include every tuple
\[
(u_1^{(j_1)},\ldots,u_{K_{\rm a}}^{(j_{K_{\rm a}})})
\]
of pairwise distinct clones, in the displayed order. Retain repeated edge
occurrences as separately labelled edges. Every occurrence edge has a lift,
since there are at least $K_{\rm a}$ clones of each vertex. Giving clones
their source vertex's colour preserves the colour multiset of every edge,
and hence the YES promise.

Conversely, suppose the cloned hypergraph has a proper $S_{\rm a}$-colouring.
Give each source vertex a colour appearing on at least $K_{\rm a}$ of its
clones. A monochromatic occurrence edge in this induced colouring would lift
to a monochromatic edge of distinct clones: each source vertex occurs at
most $K_{\rm a}$ times, so choose different clones of that colour for its
positions. This contradicts properness. The induced colouring would therefore
be proper, contradicting the NO promise.

Each vertex produces $K_{\rm a}S_{\rm a}$ clones and each edge produces at
most $(K_{\rm a}S_{\rm a})^{K_{\rm a}}$ lifts. These are constants for fixed
$t$, so enumeration is deterministic and polynomial-time, with degree allowed
to depend on $t$. If the vertex set is empty, add one isolated labelled
vertex. The resulting nonempty vertex set and ordered edges of distinct
vertices give exactly the position map required by
\cref{def:alternating-table}.
\end{proof}

\section{Finite-field arrays}\label{sec:arrays}
\label{sec:alternating_implementation}

We use the polynomial-root construction of Alon, Goldreich,
H{\aa}stad and Peralta~\cite[Construction~3 and Proposition~3]{alon1992smallbias}.
The conversion from parity bounds to pattern probabilities is also recorded
there, with credit to Vazirani~\cite[Lemma~1 and Corollary~1]{alon1992smallbias}.
The following statement gives the exact interface needed by the lookup proof.

\begin{lemma}\label{yao:finite-field-list}
For integers $n,r\ge1$, a deterministic algorithm prints a nonempty labelled
list of $n$-bit arrays whose uniform law satisfies
\[
\Pr[w|_I=z]\le2^{-|I|}+2^{-r}
\qquad(I\subseteq\{0,\ldots,n-1\},\ z\in\{0,1\}^I).
\]
Its running time and complete output length are polynomial in $n2^r$.
\end{lemma}
\begin{proof}
Let $R=2^d$ be the least power of two at least $n2^r$. Then $d\ge1$.
In a fixed binary basis of $\mathbb F_R$, write $\operatorname{vec}(y)$ for
the coordinate vector of $y$. For each seed
$(u,y)\in\mathbb F_2^d\times\mathbb F_R$, define
\[
w^{u,y}(j)=\langle u,\operatorname{vec}(y^j)\rangle_{\mathbb F_2},
\qquad 0\le j<n,
\]
where $y^0=1$, also when $y=0$. Give these $R^2$ seed labels the uniform law;
equal arrays keep separate labels.

For a nonempty position set $U$, the polynomial
\[
P_U(Z)=\sum_{j\in U}Z^j
\]
is nonzero and has degree at most $n-1$. Conditional on $y$, the parity of the
bits at $U$ is $\langle u,\operatorname{vec}(P_U(y))\rangle$. Its sign has
mean zero over uniform $u$ unless $P_U(y)=0$, when its mean is one. The
polynomial root bound therefore gives
\[
\left|\mathbb E_{u,y}(-1)^{\sum_{j\in U}w^{u,y}(j)}\right|
=\Pr_y[P_U(y)=0]\le\frac{n-1}{R}\le2^{-r}.
\]
This includes $n=1$, when the only polynomial is the constant one.
Now expand the pattern indicator:
\[
\ind[w|_I=z]=2^{-|I|}\sum_{U\subseteq I}
(-1)^{\sum_{j\in U}(w(j)+z(j))}.
\]
The empty-set term contributes $2^{-|I|}$, and the remaining terms contribute
at most $(1-2^{-|I|})2^{-r}$, proving the pattern bound, also for empty $I$.

For constructibility, enumerate monic binary polynomials of degree $d$ until
finding an irreducible one. Test each by division by every monic polynomial of
degrees $1$ through $\floor{d/2}$: a reducible polynomial has a factor in that
range, and an irreducible polynomial exists in every positive degree. For
$d=1$ the test range is empty. There are $R$ candidates and fewer than $2R$
divisors, so the search takes polynomial time in $R$. The quotient field then
supports arithmetic in polynomial time in $d$. Enumerate all $R^2$ seeds and
compute successive powers and binary dot products to print each array.
Since $R<2n2^r$, this takes polynomial time and output length in $n2^r$,
including the field representation and seed labels.
\end{proof}

\begin{proof}[Proof of \cref{yao:array-list}]
Number $\Omega_{\rm a}$ by $0,\ldots,n-1$ and apply
\cref{yao:finite-field-list} with
\[
n=K_{\rm a}H_{\rm a},\qquad r=4(b+3t).
\]
Here $2^{-r}=\epsilon_{\rm a}$ and
$n2^r=K_{\rm a}2^{12t}H_{\rm a}^5$, which is polynomial in $H_{\rm a}$
for fixed $t$. The integer parameters have $O_t(b+1)$ bits, and the bound
includes $H_{\rm a}=1$.
\end{proof}

\bibliographystyle{plainnat}
\bibliography{references}
\end{document}
